% Las conductas de riesgo de los jóvenes : identificarlas en el medio escolar y social — Espagnol 1ère A — Cours signé OnBuch+
% Programme officiel MINESEC. Compiler avec : tectonic EA03-conductas-de-riesgo-identificar.tex
\newcommand{\DOCMATIERE}{Espagnol}
\newcommand{\DOCNIVEAU}{1ère A}
\newcommand{\DOCMODULE}{Módulo 1 : Vie familiale et intégration sociale}
\newcommand{\DOCLECON}{EA03}
\newcommand{\DOCTITRE}{Las conductas de riesgo de los jóvenes : identificarlas en el medio escolar y social}
% OnBuch+ — préambule « cours signés » ENRICHI (compilé avec tectonic / XeLaTeX)
% Système visuel « Pop » d'OnBuch+ : fond crème (jamais blanc), bordures encre
% épaisses, ombres « dures » décalées, titres Archivo Black en capitales.
% Version MATHÉMATIQUES : graphes (pgfplots/tikz), boîtes pédagogiques
% (définition, propriété, méthode, attention, le savais-tu, exercice résolu,
% exercice, TP, bilan), page de garde et signature OnBuch+ ancrée en bas.
%
% Macros attendues dans main.tex AVANT \input{preamble} :
%   \DOCMATIERE  \DOCNIVEAU  \DOCMODULE  \DOCLECON (numéro)  \DOCTITRE
\documentclass[11pt]{article}

\usepackage{fontspec}
\usepackage[spanish,french]{babel} % français = langue principale ; l'espagnol via \esp{..} / espagnol
\usepackage[a4paper,margin=2cm,top=3.4cm,bottom=2.8cm,headheight=1.4cm,footskip=1.3cm]{geometry}
\usepackage{amsmath,amssymb,amsfonts}
\usepackage{mathtools} % \xmapsto, \coloneqq... souvent utilisés par les modèles
\usepackage{cancel} % simplifications barrées (bilans de forces)
% \abs{z} (module, valeur absolue) et \norm{u} (norme), utilisés par les modèles
\providecommand{\abs}{}\let\abs\relax\DeclarePairedDelimiter{\abs}{\lvert}{\rvert}
\providecommand{\norm}{}\let\norm\relax\DeclarePairedDelimiter{\norm}{\lVert}{\rVert}
\usepackage[table]{xcolor} % [table] : fournit aussi \rowcolors (alternance de lignes)
\usepackage{pagecolor}
\usepackage{tikz}
\usetikzlibrary{arrows.meta,positioning,calc,shapes.geometric,shapes.misc,shapes.arrows,decorations.pathmorphing,decorations.pathreplacing,decorations.markings,patterns,shadows,mindmap,trees,fit,backgrounds,matrix,intersections,angles,quotes}
\usepackage{pgfplots}
\usepackage[version=4]{mhchem} % équations nucléaires \ce{^{235}_{92}U}
\usepackage[european,siunitx]{circuitikz} % schémas électriques
\pgfplotsset{compat=1.18}
\usepgfplotslibrary{fillbetween,groupplots} % aires entre courbes (calcul intégral)
\usepackage{enumitem}
\usepackage{fancyhdr}
% [most] charge toutes les bibliothèques tcolorbox (listings, minted, raster,
% documentation...) et gonfle inutilement la mémoire TeX sur un document long
% (~300 boîtes) : on ne charge que ce qui est réellement utilisé.
\usepackage[skins,breakable]{tcolorbox}
\usepackage{graphicx}
\usepackage{colortbl}
\usepackage{booktabs}
\usepackage{tabularx}
\usepackage{diagbox} % case d'en-tête diagonale des tableaux à double entrée
% Colonnes extensibles centrée / gauche / droite (C, L, R), souvent utilisées par les modèles
\newcolumntype{C}{>{\centering\arraybackslash}X}
\newcolumntype{L}{>{\raggedright\arraybackslash}X}
\newcolumntype{R}{>{\raggedleft\arraybackslash}X}
\usepackage{array}
\usepackage{multicol}
\usepackage{siunitx}
% parse-numbers=false : accepte aussi \SI{2,5 \times 10^{-3}}{...}
\sisetup{locale=FR,output-decimal-marker={,},group-minimum-digits=4,parse-numbers=false}
\DeclareSIUnit\ohm{\ensuremath{\symup{\Omega}}} % U+2126 absent de la police
\DeclareSIUnit\division{div} % réglages d'oscilloscope (ms/div, V/div)
\DeclareSIUnit\tr{tr} % tours (vitesse de rotation, ex. \SI{1500}{\tr\per\minute})
\DeclareSIUnit\molar{M} % concentration molaire (ex. \SI{3}{\molar}, \SI{1}{\micro\molar})
\DeclareSIUnit\an{an} % année (le modèle confond parfois avec \year, un compteur TeX) — ex. \SI{5}{\kilo\meter\per\an}
\DeclareSIUnit\FCFA{FCFA} % franc CFA (ex. \SI{500}{\FCFA\per\kilo\gram})
\DeclareSIUnit\degreeBrix{\textdegree Brix} % teneur en sucre d'un sirop/jus (réfractométrie)
\DeclareSIUnit\month{mois} % le modèle utilise parfois \month, un compteur TeX, comme unité de durée
\usepackage{qrcode}
% \degree : commande fréquemment utilisée par les modèles pour un angle ou une
% température en dehors de \SI{}{} (non fournie par les paquets chargés).
\providecommand{\degree}{^{\circ}}
% \qed (fin de démonstration), utilisé par les modèles sans amsthm
\providecommand{\qed}{\hfill$\square$}
% \barre : double barre des valeurs interdites dans un tableau de variations (tkz-tab)
\providecommand{\barre}{\ensuremath{\|}}
% environnement proof (amsthm non chargé) utilisé par les modèles
\ifdefined\proof\else\newenvironment{proof}[1][Démonstration]{\par\noindent\textit{#1.}\ }{\qed\par}\fi
\usepackage{lastpage}
\usepackage{hyperref}
\hypersetup{hidelinks}

% ============================================================
% Palette « Pop » OnBuch+ — mêmes hex que la classe P (plus_ui.dart)
% ============================================================
\definecolor{poporange}{HTML}{F59321}
\definecolor{popdark}{HTML}{E07A0C}
\definecolor{popcream}{HTML}{FFF6E8}
\definecolor{popink}{HTML}{1C1714}
\definecolor{poppurple}{HTML}{7A5AE0}
\definecolor{popgreen}{HTML}{2E9E6A}
\definecolor{poppink}{HTML}{E0457B}
\definecolor{popblue}{HTML}{2F7DE1}
\definecolor{popgold}{HTML}{C98A12}
\definecolor{popmuted}{HTML}{8A7F76}
\definecolor{popline}{HTML}{EADFD2}
\definecolor{popgray}{HTML}{8A7F76}
\colorlet{popgrey}{popgray}
% Alias tolérants (le modèle nomme parfois les couleurs différemment)
\colorlet{popwhite}{white}
\colorlet{popblack}{popink}
\colorlet{popred}{poppink}
\colorlet{popyellow}{popgold}
% Teintes claires (fonds de tableaux/encadrés) — le modèle nomme parfois
% les couleurs avec un suffixe L (ex. popblueL) sans les avoir définies.
\colorlet{poporangeL}{poporange!20}
\colorlet{popdarkL}{popdark!20}
\colorlet{poppurpleL}{poppurple!20}
\colorlet{popgreenL}{popgreen!20}
\colorlet{poppinkL}{poppink!20}
\colorlet{popblueL}{popblue!20}
\colorlet{popgoldL}{popgold!20}
\colorlet{popredL}{popred!20}
\colorlet{popgrayL}{popgray!20}
% Teintes claires pour fonds de tableaux / zones
\colorlet{popblueL}{popblue!10!white}
\colorlet{poporangeL}{poporange!14!white}
\colorlet{popgreenL}{popgreen!12!white}
\colorlet{poppurpleL}{poppurple!12!white}
\colorlet{poppinkL}{poppink!12!white}
\colorlet{popgoldL}{popgold!14!white}

\pagecolor{popcream}

% ============================================================
% Typographie
% ============================================================
\defaultfontfeatures{Ligatures=TeX}
\setmainfont{PlusJakartaSans-Medium.ttf}[Path=../../../fonts/,BoldFont=PlusJakartaSans-Bold.ttf]
% Maths en sans-serif (Fira Math) avec chiffres et lettres droites de Plus
% Jakarta, pour que les formules se fondent dans le texte.
\usepackage[math-style=ISO,bold-style=ISO]{unicode-math}
\setmathfont{FiraMath-Regular.otf}[Path=../../../fonts/]
\setmathfont{PlusJakartaSans-Medium.ttf}[Path=../../../fonts/,range={up/num,up/{Latin,latin},\mathup}]
\usepackage{icomma} % virgule décimale sans espace en mode math
% Caractères mathématiques Unicode absents des polices de texte (Plus Jakarta,
% Archivo Black) : rendus par leur équivalent mathématique, en texte comme en
% formule — sinon ils s'affichent en carré vide (ex. « Arithmétique dans ℤ »).
\usepackage{newunicodechar}
\newunicodechar{ℝ}{\ensuremath{\mathbb{R}}}
\newunicodechar{ℤ}{\ensuremath{\mathbb{Z}}}
\newunicodechar{ℂ}{\ensuremath{\mathbb{C}}}
\newunicodechar{ℕ}{\ensuremath{\mathbb{N}}}
\newunicodechar{ℚ}{\ensuremath{\mathbb{Q}}}
\newunicodechar{𝔻}{\ensuremath{\mathbb{D}}}
\newunicodechar{α}{\ensuremath{\alpha}}
\newunicodechar{β}{\ensuremath{\beta}}
\newunicodechar{γ}{\ensuremath{\gamma}}
\newunicodechar{δ}{\ensuremath{\delta}}
\newunicodechar{ε}{\ensuremath{\varepsilon}}
\newunicodechar{θ}{\ensuremath{\theta}}
\newunicodechar{λ}{\ensuremath{\lambda}}
\newunicodechar{μ}{\ensuremath{\mu}}
\newunicodechar{σ}{\ensuremath{\sigma}}
\newunicodechar{τ}{\ensuremath{\tau}}
\newunicodechar{φ}{\ensuremath{\varphi}}
\newunicodechar{ψ}{\ensuremath{\psi}}
\newunicodechar{ω}{\ensuremath{\omega}}
\newunicodechar{Σ}{\ensuremath{\Sigma}}
% majuscules produites par \MakeUppercase dans les titres (α-aminés -> Α)
\newunicodechar{Α}{\ensuremath{\alpha}}
\newunicodechar{Β}{\ensuremath{\beta}}
\newunicodechar{Γ}{\ensuremath{\Gamma}}
\newunicodechar{Δ}{\ensuremath{\Delta}}
\newunicodechar{Ε}{\ensuremath{\varepsilon}}
\newunicodechar{Θ}{\ensuremath{\Theta}}
\newunicodechar{Λ}{\ensuremath{\Lambda}}
\newunicodechar{Μ}{\ensuremath{\mu}}
\newunicodechar{Τ}{\ensuremath{\tau}}
\newunicodechar{Φ}{\ensuremath{\Phi}}
\newunicodechar{Ψ}{\ensuremath{\Psi}}
\newunicodechar{Ω}{\ensuremath{\Omega}}
\newunicodechar{↦}{\ensuremath{\mapsto}}
\newunicodechar{∈}{\ensuremath{\in}}
\newunicodechar{∉}{\ensuremath{\notin}}
\newunicodechar{∧}{\ensuremath{\wedge}}
\newunicodechar{⇔}{\ensuremath{\Leftrightarrow}}
\newunicodechar{⟺}{\ensuremath{\Longleftrightarrow}}
\newunicodechar{⇒}{\ensuremath{\Rightarrow}}
\newunicodechar{⟹}{\ensuremath{\Longrightarrow}}
\newunicodechar{∘}{\ensuremath{\circ}}
\newunicodechar{∪}{\ensuremath{\cup}}
\newunicodechar{∩}{\ensuremath{\cap}}
\newunicodechar{≡}{\ensuremath{\equiv}}
\newunicodechar{⊂}{\ensuremath{\subset}}
\newunicodechar{⊥}{\ensuremath{\perp}}
\newunicodechar{∃}{\ensuremath{\exists}}
\newunicodechar{∀}{\ensuremath{\forall}}
\newunicodechar{∛}{\ensuremath{\sqrt[3]{\,}}}
\newunicodechar{∜}{\ensuremath{\sqrt[4]{\,}}}
\newunicodechar{⃗}{\ensuremath{{}^{\to}}}
\newunicodechar{ⁿ}{\ifmmode^{n}\else\textsuperscript{\ensuremath{n}}\fi}
\newunicodechar{ˣ}{\ifmmode^{x}\else\textsuperscript{\ensuremath{x}}\fi}
\newunicodechar{ᵇ}{\ifmmode^{b}\else\textsuperscript{\ensuremath{b}}\fi}
\newunicodechar{ʳ}{\ifmmode^{r}\else\textsuperscript{\ensuremath{r}}\fi}
\newunicodechar{ᵘ}{\ifmmode^{u}\else\textsuperscript{\ensuremath{u}}\fi}
\newunicodechar{ʸ}{\ifmmode^{y}\else\textsuperscript{\ensuremath{y}}\fi}
\newunicodechar{ᵃ}{\ifmmode^{a}\else\textsuperscript{\ensuremath{a}}\fi}
\newunicodechar{ᵉ}{\ifmmode^{e}\else\textsuperscript{\ensuremath{e}}\fi}
\newunicodechar{ᵏ}{\ifmmode^{k}\else\textsuperscript{\ensuremath{k}}\fi}
\newunicodechar{ᵖ}{\ifmmode^{p}\else\textsuperscript{\ensuremath{p}}\fi}
\newunicodechar{ᴺ}{\ifmmode^{N}\else\textsuperscript{\ensuremath{N}}\fi}
\newunicodechar{ᶜ}{\ifmmode^{c}\else\textsuperscript{\ensuremath{c}}\fi}
\newunicodechar{ʰ}{\ifmmode^{h}\else\textsuperscript{\ensuremath{h}}\fi}
\newunicodechar{ᵠ}{\ifmmode^{q}\else\textsuperscript{\ensuremath{q}}\fi}
\newunicodechar{ₙ}{\ifmmode_{n}\else\textsubscript{\ensuremath{n}}\fi}
\newunicodechar{ₐ}{\ifmmode_{a}\else\textsubscript{\ensuremath{a}}\fi}
\newunicodechar{ᵢ}{\ifmmode_{i}\else\textsubscript{\ensuremath{i}}\fi}
\newunicodechar{ᵣ}{\ifmmode_{r}\else\textsubscript{\ensuremath{r}}\fi}
\newunicodechar{ₖ}{\ifmmode_{k}\else\textsubscript{\ensuremath{k}}\fi}
\newunicodechar{ᵦ}{\ifmmode_{\beta}\else\textsubscript{\ensuremath{\beta}}\fi}
\newunicodechar{ₓ}{\ifmmode_{x}\else\textsubscript{\ensuremath{x}}\fi}
\newfontfamily\popdisplay{ArchivoBlack-Regular.ttf}[Path=../../../fonts/]
\newfontfamily\poplogo{SpaceGrotesk-Bold.ttf}[Path=../../../fonts/]
\newfontfamily\popsemi{PlusJakartaSans-SemiBold.ttf}[Path=../../../fonts/]
\newfontfamily\popxbold{PlusJakartaSans-ExtraBold.ttf}[Path=../../../fonts/]
\newfontfamily\popxboldls{PlusJakartaSans-ExtraBold.ttf}[Path=../../../fonts/,LetterSpace=3.0]

\setlist[enumerate]{leftmargin=1.7em,itemsep=5pt,topsep=4pt}
\setlist[itemize]{leftmargin=1.7em,itemsep=4pt,topsep=4pt,label={\color{poporange}\textbullet}}
\setlength{\parindent}{0pt}
\setlength{\parskip}{5pt}
\renewcommand{\arraystretch}{1.25}

% pgfplots : style Pop par défaut
\pgfplotsset{
  every axis/.append style={axis line style={popink,line width=1pt},tick style={popink},
    grid style={popline,line width=0.5pt},label style={font=\small},tick label style={font=\footnotesize}},
  popcurve/.style={poporange,line width=1.8pt,no markers,smooth},
}
% Même style disponible dans un \draw TikZ simple (hors pgfplots)
\tikzset{popcurve/.style={poporange,line width=1.8pt}}
\tikzset{popgrid/.style={popline,very thin}}
\colorlet{popcurve}{poporange} % le nom du style est parfois utilisé comme couleur

% ============================================================
% Pill / squiggle
% ============================================================
\newcommand{\poppill}[3]{%
  \tikz[baseline=-0.55ex]\node[draw=popink,line width=1.2pt,rounded corners=2.6pt,%
    fill=#2,inner xsep=6.5pt,inner ysep=3pt]{%
    \popxboldls\fontsize{7}{7}\selectfont\color{#3}\MakeUppercase{#1}};%
}
\newcommand{\popsquiggle}[1]{%
  \tikz[baseline=-0.4ex]{
    \draw[#1,line width=2.6pt,line cap=round]
      plot [smooth] coordinates {(0,0.09) (0.34,0.01) (0.68,0.21) (1.02,0.01) (1.36,0.10)};
  }%
}
% Mot-clé surligné (marqueur orange sous le texte)
\newcommand{\cle}[1]{{\popxbold\color{popdark}#1}}

% ============================================================
% En-tête / pied
% ============================================================
\pagestyle{fancy}
\fancyhf{}
\renewcommand{\headrulewidth}{0pt}
\renewcommand{\footrulewidth}{0pt}
\lhead{\raisebox{-0.15ex}{\poplogo\fontsize{13}{13}\selectfont\color{popink}OnBuch\color{poporange}+}}
\rhead{\poppill{\DOCMATIERE\ \textbullet\ \DOCNIVEAU\ \textbullet\ Leçon \DOCLECON}{white}{popink}}
\lfoot{\popxbold\fontsize{7.6}{7.6}\selectfont\color{popmuted}\MakeUppercase{Cours signé OnBuch+}\ \textbullet\ {\popsemi\DOCTITRE}}
\rfoot{\tikz[baseline=-0.6ex]\node[rounded corners=8.5pt,fill=popink,minimum height=17pt,minimum width=17pt,inner xsep=5pt,inner sep=0]{\popxbold\fontsize{8}{8}\selectfont\color{popcream}\thepage\,/\,\pageref*{LastPage}};}

% ============================================================
% Titres
% ============================================================
\newcommand{\chaptitre}[2]{%
  \begin{tcolorbox}[enhanced,colback=poporange,colframe=popink,boxrule=2.6pt,arc=11pt,
      drop shadow=popink,left=15pt,right=15pt,top=12pt,bottom=11pt,before skip=3pt,after skip=18pt]
    \poppill{#2}{popink}{popcream}\par\vspace{9pt}
    {\raggedright\hyphenpenalty=10000\exhyphenpenalty=10000\popdisplay\fontsize{22}{24}\selectfont\color{popcream}\MakeUppercase{#1}\par}
  \end{tcolorbox}
}
% Section numérotée : pastille orange avec le numéro + titre Archivo
\newcounter{popsec}
\newcommand{\coursec}[1]{%
  \stepcounter{popsec}\setcounter{popsub}{0}%
  \par\vspace{16pt}\noindent
  \begin{minipage}{\linewidth}
  \tikz[baseline=-0.7ex]\node[draw=popink,line width=1.6pt,fill=poporange,rounded corners=4pt,minimum size=22pt,inner sep=0]{\popdisplay\fontsize{12}{12}\selectfont\color{popcream}\thepopsec};%
  \hspace{8pt}{\popdisplay\fontsize{15}{17}\selectfont\color{popink}\MakeUppercase{#1}}\par
  \vspace{4pt}\noindent{\color{poporange}\rule{36pt}{3.6pt}}
  \end{minipage}\par\nopagebreak\vspace{8pt}
}
\newcounter{popsub}
\newcommand{\courssub}[1]{%
  \stepcounter{popsub}%
  \par\vspace{9pt}\noindent{\popxbold\fontsize{11.5}{13}\selectfont\color{popink}{\color{poporange}\thepopsec.\thepopsub}\hspace{6pt}#1}\par\nopagebreak\vspace{4pt}
}

% ============================================================
% Boîtes pédagogiques — même recette : fond blanc, bordure épaisse colorée,
% ombre dure encre, badge plein. Titre optionnel [#1].
% ============================================================
\tcbset{popbase/.style={breakable,enhanced,colback=white,boxrule=2.3pt,arc=9pt,
  shadow={3pt}{-3pt}{0pt}{popink},left=14pt,right=14pt,top=11pt,bottom=11pt,before skip=11pt,after skip=15pt,
  fontupper=\popsemi\fontsize{10.6}{14.5}\selectfont\color{popink},
  fontlower=\popsemi\fontsize{10.6}{14.5}\selectfont\color{popink}}}
% Coche dessinée (le glyphe ✓ n'existe pas dans les polices du document)
\newcommand{\popcheck}[1]{\tikz[baseline=-0.5ex]\draw[#1,line width=1.6pt,line cap=round,line join=round](-0.13,0.0)--(-0.04,-0.09)--(0.13,0.1);}
\providecommand{\coche}{\popcheck{popgreen}}
\providecommand{\resultat}[1]{\par\smallskip\noindent\fcolorbox{popgreen}{popgreenL}{\parbox{\dimexpr\linewidth-2\fboxsep-2\fboxrule\relax}{\textbf{Résultat.} #1}}\par\smallskip}
\newcommand{\popboxhead}[3]{\poppill{#1}{#2}{white}\ifx\relax#3\relax\else\hspace{6pt}{\popxbold\fontsize{10.6}{12}\selectfont\color{popink}#3}\fi\par\vspace{7pt}}

\newtcolorbox{aretenir}[1][]{popbase,colframe=popgreen,before upper={\popboxhead{À retenir}{popgreen}{#1}}}
% Texte en espagnol (document de lecture, dialogue, poème, extrait) : cadre
% d'étude. Un titre optionnel (source, auteur) s'affiche dans l'en-tête.
\newtcolorbox{textebox}[1][]{popbase,colframe=popblue,before upper={\popboxhead{Texto}{popblue}{#1}\begin{otherlanguage}{spanish}},after upper={\end{otherlanguage}}}
% Marques ✓ / ✗ (forme correcte / fautive) souvent utilisées par les modèles
\providecommand{\cross}{\textcolor{poppink}{\ensuremath{\times}}}
\providecommand{\xmark}{\cross}\providecommand{\crossmark}{\cross}\providecommand{\cmark}{\ensuremath{\checkmark}}
% Ligne de réponse / justification à compléter (inventée par les modèles)
\providecommand{\justification}[1]{\par\noindent\textit{Justification~:} \dotfill\par}
% \textipa (package tipa non chargé) : la police ne couvre pas l'API -> simple passage
\providecommand{\textipa}[1]{#1}
% Soulignement (ulem non chargé) : \uline, \ul
\providecommand{\uline}[1]{\underline{#1}}\providecommand{\ul}[1]{\underline{#1}}
% Mot ou phrase en espagnol à l'intérieur d'un paragraphe français
\newcommand{\esp}[1]{\ifmmode\text{\textit{#1}}\else\begingroup\selectlanguage{spanish}\itshape #1\endgroup\fi}
\newtcolorbox{exemplebox}[1][]{popbase,colframe=poporange,before upper={\popboxhead{Exemple}{poporange}{#1}}}
\newtcolorbox{definition}[1][]{popbase,colframe=popblue,before upper={\popboxhead{Définition}{popblue}{#1}}}
\newtcolorbox{propriete}[1][]{popbase,colframe=poppurple,before upper={\popboxhead{Propriété}{poppurple}{#1}}}
\newtcolorbox{methode}[1][]{popbase,colframe=popgold,before upper={\popboxhead{Méthode}{popgold}{#1}}}
\newtcolorbox{attention}[1][]{popbase,colframe=poppink,before upper={\popboxhead{Attention}{poppink}{#1}}}
\newtcolorbox{experience}[1][]{popbase,colframe=popblue,colback=popblueL,before upper={\popboxhead{Expérience}{popblue}{#1}}}
% « Le savais-tu ? » : carte violette pleine (contexte, culture, Cameroun)
\newtcolorbox{savaistu}[1][]{popbase,colback=poppurple,colframe=popink,
  fontupper=\popsemi\fontsize{10.4}{14.2}\selectfont\color{white},
  before upper={\poppill{Le savais-tu\,?}{white}{poppurple}\ifx\relax#1\relax\else\hspace{6pt}{\popxbold\color{white}#1}\fi\par\vspace{7pt}}}
% Exercice résolu : énoncé en haut, \tcblower puis la solution sur fond crème
\newcounter{popexo}\newcounter{popexr}
\newtcolorbox{exoresolu}[1][]{popbase,colframe=poporange,colbacklower=popcream,
  segmentation style={popink,line width=1.2pt,dashed},
  before upper={\stepcounter{popexr}\popboxhead{Exercice résolu \thepopexr}{poporange}{#1}},
  before lower={\poppill{Solution}{popink}{popcream}\par\vspace{6pt}}}
% Exercice (non résolu en place) — la correction est dans la partie Corrigés
\newtcolorbox{exercice}[1][]{popbase,colframe=popink,
  before upper={\stepcounter{popexo}\popboxhead{Exercice \thepopexo}{popink}{#1}}}
% Corrigé
\newtcolorbox{corrige}[1][]{popbase,colframe=popgreen,colback=popgreenL,
  before upper={\popboxhead{Corrigé}{popgreen}{#1}}}
% Objectifs / prérequis / bilan
\newtcolorbox{objectifs}{popbase,colframe=popink,colback=poporangeL,
  before upper={\popboxhead{Objectifs de la leçon}{popink}{}}}
\newtcolorbox{prerequis}{popbase,colframe=popink,colback=popblueL,
  before upper={\popboxhead{Prérequis}{popblue}{}}}
\newtcolorbox{bilan}{popbase,colframe=popink,colback=white,boxrule=2.8pt,
  before upper={\popboxhead{Fiche bilan}{poporange}{L'essentiel à connaître pour l'examen}}}
% Figure encadrée avec légende
\newtcolorbox{popfigure}[1][]{enhanced,breakable=false,colback=white,colframe=popline,boxrule=1.4pt,arc=8pt,
  left=10pt,right=10pt,top=10pt,bottom=8pt,before skip=10pt,after skip=12pt,halign=center,
  lower separated=false,halign lower=center,
  fontlower=\popsemi\fontsize{9}{11.5}\selectfont\color{popmuted},
  #1}
\newcommand{\legende}[1]{\par\vspace{6pt}{\popsemi\fontsize{9}{11.5}\selectfont\color{popmuted}#1}}

% ============================================================
% Signature OnBuch+
% ============================================================
\newcommand{\poplogoinline}[1]{{\poplogo\fontsize{#1}{#1}\selectfont\color{popink}OnBuch\color{poporange}+}}
\newcommand{\signatureonbuch}{%
  \begin{tcolorbox}[enhanced,colback=popink,colframe=popink,boxrule=2.2pt,arc=10pt,
      drop shadow=poporange,left=14pt,right=14pt,top=10pt,bottom=10pt,before skip=0pt,after skip=0pt]
    \tikz[baseline=-0.7ex]\node[circle,fill=popgreen,draw=popcream,line width=1.3pt,minimum size=22pt,inner sep=0]{\popcheck{white}};%
    \hspace{9pt}\begin{minipage}[c]{0.62\linewidth}
      {\popxbold\fontsize{12}{13}\selectfont\color{popcream}Signé\ \textbullet\ {\poplogo OnBuch\color{poporange}+}}\par\vspace{2pt}
      {\raggedright\popsemi\fontsize{8}{10.5}\selectfont\color{popline}\MakeUppercase{Équipe pédagogique OnBuch+}\\\MakeUppercase{Conforme au programme officiel MINESEC}\par}
    \end{minipage}\hfill
    {\poplogo\fontsize{22}{22}\selectfont\color{popcream}OnBuch\color{poporange}+}
  \end{tcolorbox}%
}

% ============================================================
% Page de garde
% #1 titre · #2 matière · #3 niveau · #4 module · #5 n° leçon · #6 durée ·
% #7 description / objectifs (texte ou itemize)
% ============================================================
\newcommand{\pagegarde}[7]{%
  \thispagestyle{empty}
  \newgeometry{a4paper,margin=1.7cm,top=1.6cm,bottom=1.4cm}%
  \begin{tikzpicture}[remember picture,overlay]
    \fill[poporange!18] ([xshift=-3.2cm,yshift=-3.6cm]current page.north east) circle (4.6cm);
    \fill[poppurple!14] ([xshift=2.4cm,yshift=7.6cm]current page.south west) circle (3.8cm);
  \end{tikzpicture}%
  {\poplogo\fontsize{30}{30}\selectfont\color{popink}OnBuch\color{poporange}+}\hfill
  \raisebox{0.6ex}{\poppill{Cours signé \textbullet\ Premium}{popink}{popcream}}\par
  \vspace{1pt}\popsquiggle{poppurple}\par
  \vspace{8pt}
  \begin{tcolorbox}[enhanced,colback=poporange,colframe=popink,boxrule=2.8pt,arc=13pt,
      drop shadow=popink,left=18pt,right=18pt,top=12pt,bottom=13pt,after skip=10pt]
    \poppill{#2 \textbullet\ #3}{popink}{popcream}\hspace{5pt}\poppill{#4}{white}{popink}\par\vspace{8pt}
    {\popdisplay\fontsize{14}{14}\selectfont\color{popink}LEÇON #5}\par\vspace{4pt}
    {\raggedright\hyphenpenalty=10000\exhyphenpenalty=10000\popdisplay\fontsize{25}{27}\selectfont\color{popcream}\MakeUppercase{#1}\par}
    \vspace{8pt}
    {\popxbold\fontsize{9.5}{9.5}\selectfont\color{popink}\MakeUppercase{Durée conseillée : #6}}
  \end{tcolorbox}
  \begin{tcolorbox}[enhanced,colback=white,colframe=popink,boxrule=2.2pt,arc=10pt,
      drop shadow=popink,left=16pt,right=16pt,top=10pt,bottom=10pt,after skip=8pt]
    \poppill{Dans cette leçon}{poppurple}{white}\par\vspace{5pt}
    \popsemi\fontsize{9.3}{12.6}\selectfont\color{popink}#7
  \end{tcolorbox}
  \noindent\begin{minipage}[t]{0.487\linewidth}
    \begin{tcolorbox}[enhanced,colback=popgreen,colframe=popink,boxrule=2.2pt,arc=10pt,
        drop shadow=popink,left=11pt,right=11pt,top=10pt,bottom=10pt,height=3.1cm,valign=center]
      \tcbox[colback=white,colframe=popink,boxrule=1.3pt,arc=5pt,boxsep=0pt,nobeforeafter,
        left=3pt,right=3pt,top=3pt,bottom=3pt,valign=center]{\qrcode[height=1.9cm]{https://play.google.com/store/apps/details?id=com.luvvix.onbuch&hl=fr}}%
      \hspace{8pt}\begin{minipage}[c]{0.5\linewidth}
        {\popdisplay\fontsize{11}{12.5}\selectfont\color{white}\MakeUppercase{Télécharge OnBuch}}\par\vspace{4pt}
        {\raggedright\popsemi\fontsize{8.4}{11}\selectfont\color{white}Gratuit \textbullet\ Google Play\\Tuteur IA, annales, concours, résultats\par}
      \end{minipage}
    \end{tcolorbox}
  \end{minipage}\hfill
  \begin{minipage}[t]{0.487\linewidth}
    \begin{tcolorbox}[enhanced,colback=poppurple,colframe=popink,boxrule=2.2pt,arc=10pt,
        drop shadow=popink,left=11pt,right=11pt,top=10pt,bottom=10pt,height=3.1cm,valign=center]
      \tikz[baseline=-0.7ex]\node[circle,fill=white,draw=popink,line width=1.3pt,minimum size=1.6cm,inner sep=0]{\popdisplay\fontsize{24}{24}\selectfont\color{poppurple}+};%
      \hspace{8pt}\begin{minipage}[c]{0.6\linewidth}
        {\popdisplay\fontsize{11}{12.5}\selectfont\color{white}\MakeUppercase{Passe à OnBuch+}}\par\vspace{4pt}
        {\raggedright\popsemi\fontsize{8.4}{11}\selectfont\color{white}Tous les cours signés, Tuteur IA illimité, fiches PDF — onglet OnBuch+ de l'app\par}
      \end{minipage}
    \end{tcolorbox}
  \end{minipage}
  \vspace{8pt}
  \popcontenu
  \vfill
  \signatureonbuch
  \restoregeometry
  \newpage
}

% Rangée « Ce cours contient » (page de garde)
\newcommand{\poptile}[3]{% #1 couleur #2 gros texte #3 légende
  \begin{tcolorbox}[enhanced,colback=#1,colframe=popink,boxrule=2pt,arc=9pt,shadow={2.5pt}{-2.5pt}{0pt}{popink},
      left=6pt,right=6pt,top=9pt,bottom=9pt,halign=center,valign=center,height=1.75cm,width=\linewidth,nobeforeafter]
    {\popdisplay\fontsize{12.5}{13}\selectfont\color{white}\MakeUppercase{#2}}\par\vspace{3pt}
    {\centering\popsemi\fontsize{7.8}{9.5}\selectfont\color{white}#3\par}
  \end{tcolorbox}}
\newcommand{\popcontenu}{%
  \par\noindent{\popxboldls\fontsize{7.5}{7.5}\selectfont\color{popmuted}CE COURS SIGNÉ CONTIENT}\par\vspace{7pt}
  \noindent\begin{minipage}[t]{0.235\linewidth}\poptile{popblue}{Cours}{complet, illustré, pas à pas}\end{minipage}\hfill
  \begin{minipage}[t]{0.235\linewidth}\poptile{poporange}{Méthodes}{et exercices résolus}\end{minipage}\hfill
  \begin{minipage}[t]{0.235\linewidth}\poptile{poppink}{Exercices}{type examen + corrigés}\end{minipage}\hfill
  \begin{minipage}[t]{0.235\linewidth}\poptile{popgreen}{Fiche bilan}{l'essentiel à retenir}\end{minipage}\par}

% Sommaire
\newtcolorbox{sommaire}{enhanced,colback=white,colframe=popink,boxrule=2.2pt,arc=10pt,
  drop shadow=popink,left=18pt,right=18pt,top=13pt,bottom=13pt,before skip=6pt,after skip=20pt,
  before upper={\poppill{Sommaire}{poppurple}{white}\par\vspace{9pt}\popsemi\fontsize{10.6}{14.5}\selectfont\color{popink}}}

% Signature de fin de cours — poussée tout en bas de la dernière page
\newcommand{\findecours}{%
  \par\vfill\nopagebreak
  \begin{center}\popsquiggle{poporange}\end{center}\vspace{4pt}
  \signatureonbuch
}
\AtBeginDocument{\def\llbracket{\mathopen{[\mkern-3.5mu[}}\def\rrbracket{\mathclose{]\mkern-3.5mu]}}} % intervalles d'entiers (glyphe ⟦ absent de la police)
% Glyphes absents de Fira Math : redéfinis à partir de glyphes disponibles
\AtBeginDocument{%
  \def\setminus{\mathbin{\backslash}}\let\smallsetminus\setminus
  \def\vdots{\mathord{\vbox{\baselineskip3.2pt\lineskiplimit0pt\kern4pt\hbox{.}\hbox{.}\hbox{.}}}}%
  \def\ddots{\mathinner{\mkern1mu\raise6.5pt\hbox{.}\mkern2mu\raise3.6pt\hbox{.}\mkern2mu\raise0.7pt\hbox{.}\mkern1mu}}%
  \def\triangle{\mathord{\scalebox{0.9}{$\symup{\Delta}$}}}%
  \def\nabla{\mathord{\raisebox{\depth}{\scalebox{1}[-1]{$\symup{\Delta}$}}}}%
  \def\checkmark{\popcheck{popdark}}%
  \def\star{\mathbin{\ast}}\def\bigstar{\mathord{\ast}}%
  \def\diamond{\mathbin{\scalebox{0.6}{$\Diamond$}}}%
  \def\bigcup{\mathop{\mathchoice{\raisebox{-0.3em}{\scalebox{1.7}{$\cup$}}}{\scalebox{1.25}{$\cup$}}{\cup}{\cup}}\displaylimits}%
  \def\bigcap{\mathop{\mathchoice{\raisebox{-0.3em}{\scalebox{1.7}{$\cap$}}}{\scalebox{1.25}{$\cap$}}{\cap}{\cap}}\displaylimits}%
  \def\lesseqqgtr{\mathrel{\lessgtr}}\def\bigoplus{\mathop{\oplus}\displaylimits}%
}
\providecommand{\ssi}{si et seulement si} % abréviation fréquente
\AtBeginDocument{\providecommand{\wideparen}{\overparen}} % arc de cercle

\begin{document}

\pagegarde{Las conductas de riesgo de los jóvenes : identificarlas en el medio escolar y social}{Espagnol}{1ère A}{Módulo 1 : Vie familiale et intégration sociale}{EA03}{2 h}{Cette leçon de type texte propose la lecture et le commentaire d'un article sur les comportements à risque des jeunes (drogue, alcool, harcèlement, sexting...). Elle permet d'acquérir le lexique thématique, de décrire une situation avec le présent de l'indicatif et la périphrase « estar + gérondif », et de rédiger un court paragraphe descriptif.
\vspace{4pt}
\begin{multicols}{2}\small
\begin{itemize}
\item À la fin de cette leçon, je sais lire et comprendre un texte espagnol sur les conduites à risque des jeunes.
\item À la fin de cette leçon, je sais nommer en espagnol les principaux comportements à risque (droga, alcohol, tabaco, violencia, acoso escolar, sexting, absentismo, delincuencia, adicción).
\item À la fin de cette leçon, je sais distinguer les conduites qui apparaissent dans le milieu scolaire de celles du milieu social.
\item À la fin de cette leçon, je sais utiliser le présent de l'indicatif pour décrire une situation habituelle.
\item À la fin de cette leçon, je sais employer la périphrase « estar + gérondif » pour décrire une action en cours.
\item À la fin de cette leçon, je sais dégager l'idée principale et les idées secondaires d'un texte argumentatif.
\end{itemize}\end{multicols}}

\begin{sommaire}
\begin{multicols}{2}
\begin{enumerate}
\item Texte support : « Jóvenes en riesgo »
\item Compréhension du texte
\item Lexique thématique par champs
\item Grammaire : présent de l'indicatif et estar + gérondif
\item Méthode du commentaire de texte
\item Production guidée
\item Activité d'intégration
\item Méthodes et exercices résolus
\item Exercices
\item Corrigés des exercices
\item Fiche bilan
\end{enumerate}
\end{multicols}
\end{sommaire}

\begin{objectifs}
\begin{itemize}
\item À la fin de cette leçon, je sais lire et comprendre un texte espagnol sur les conduites à risque des jeunes.
\item À la fin de cette leçon, je sais nommer en espagnol les principaux comportements à risque (droga, alcohol, tabaco, violencia, acoso escolar, sexting, absentismo, delincuencia, adicción).
\item À la fin de cette leçon, je sais distinguer les conduites qui apparaissent dans le milieu scolaire de celles du milieu social.
\item À la fin de cette leçon, je sais utiliser le présent de l'indicatif pour décrire une situation habituelle.
\item À la fin de cette leçon, je sais employer la périphrase « estar + gérondif » pour décrire une action en cours.
\item À la fin de cette leçon, je sais dégager l'idée principale et les idées secondaires d'un texte argumentatif.
\item À la fin de cette leçon, je sais rédiger un court paragraphe descriptif sur un comportement à risque observé dans mon environnement.
\item À la fin de cette leçon, je sais réagir oralement à une situation de conduite à risque avec un vocabulaire approprié.
\end{itemize}
\end{objectifs}

\begin{prerequis}
\begin{itemize}
\item Le présent de l'indicatif des verbes réguliers et irréguliers courants (Seconde)
\item Le vocabulaire de l'école et de la famille (Seconde)
\item Les techniques de compréhension de texte : repérage du thème, des personnages, du lieu (Seconde)
\item Les adjectifs qualificatifs et leur accord (Seconde)
\item Les connecteurs simples : y, pero, porque, también (Seconde)
\end{itemize}
\end{prerequis}

\newpage

\coursec{Texte support : « Jóvenes en riesgo »}

Au lycée de Yaoundé, la conseillère d'orientation reçoit de plus en plus d'élèves en difficulté : absentéisme répété, bagarres dans la cour, téléphones portables utilisés pour humilier un camarade sur les réseaux sociaux. En Espagne comme au Cameroun, les médias parlent chaque semaine de \esp{conductas de riesgo} chez les adolescentes et \esp{los adolescentes}. Comment nommer ces comportements en espagnol ? Comment les décrire précisément dans un texte ? C'est tout l'enjeu de cette leçon : lire, comprendre et décrire les conduites à risque des jeunes dans le milieu scolaire et social.

\courssub{Situation d'accroche et entrée dans la leçon}

Avant d'ouvrir le manuel, l'enseignant projette au tableau deux photos contrastées : une cour de récréation animée, joyeuse, où des élèves rient ensemble, et une seconde où un élève isolé regarde son écran, l'air anxieux, tandis qu'un groupe murmure en le regardant. La question est posée en espagnol, puis reformulée en français : \esp{¿Qué pasa en esta escuela?} « Que se passe-t-il dans cet établissement ? ». Les élèves formulent des hypothèses : \esp{problemas, peleas, drogas, acoso}. L'enseignant note les mots-clés au tableau, les classe par catégories (substances, violence, école, numérique) et annonce le titre de la leçon : \textbf{Las conductas de riesgo de los jóvenes}. Le texte support que nous allons étudier est un article de presse court, de type reportage social, publié dans un journal éducatif hispanophone fictif \emph{« Escuela Abierta »}. Il décrit la réalité de nombreux instituts (lycées) en Espagne et, par extension, dans tout l'espace hispanophone, y compris en Guinée équatoriale.

\courssub{Lecture globale silencieuse puis lecture à voix haute}

\begin{methode}[Démarche de lecture en trois temps]
\begin{enumerate}
\item \textbf{Lecture silencieuse (3 min) :} L'élève lit seul, sans dictionnaire, pour saisir l'idée générale. Il souligne les mots transparents (cognats) : \esp{institutos, alcohol, drogas, consecuencias, adicción, delincuencia, familia}.
\item \textbf{Lecture expressive par l'enseignant :} Le professeur lit le texte à voix haute, en marquant les liaisons, l'intonation des phrases déclaratives et l'accent tonique (ex. \esp{es-tán}, \esp{au-men-tan-do}, \esp{pre-o-cu-pa}). Les élèves suivent du doigt.
\item \textbf{Lecture en écho (chorale) :} La classe relit ensemble à voix haute pour fixer la prosodie et la prononciation des mots nouveaux (\esp{sexting}, \esp{absentismo}, \esp{acoso}).
\end{enumerate}
\end{methode}

\courssub{Identification du type de texte, du thème, du lieu et des personnages}

\begin{exemplebox}[Analyse par l'observation]
\begin{itemize}
\item \textbf{Type de texte :} Article de presse / reportage social (\esp{artículo de opinión / reportaje}). Indices : ton informatif, vocabulaire factuel, structure paragraphe unique, absence de dialogues directs.
\item \textbf{Thème central :} Les conduites à risque (\esp{conductas de riesgo}) chez les adolescents.
\item \textbf{Lieux (espacio) :} \esp{institutos} (lycées), \esp{patio} (cour de récréation), \esp{fiestas} (soirées), \esp{calle} (rue), \esp{redes sociales} (réseaux sociaux).
\item \textbf{Personnages (actores) :} \esp{jóvenes / alumnos} (acteurs principaux), \esp{padres} (familles), \esp{profesores} (équipe éducative).
\item \textbf{Intention communicative :} Alerter, informer, appeler à l'action conjointe (\esp{trabajar juntas}).
\end{itemize}
\end{exemplebox}

\courssub{Texte support : « Jóvenes en riesgo »}

\begin{textebox}[Jóvenes en riesgo — Artículo de « Escuela Abierta »]
En muchos institutos, algunos jóvenes están adoptando conductas de riesgo. Fuman tabaco en el patio, beben alcohol durante las fiestas y, a veces, consumen drogas. El acoso escolar también está aumentando: un alumno insulta a otro en las redes sociales o practica el sexting sin pensar en las consecuencias. Además, el absentismo escolar preocupa a los profesores: hay chicos que no vienen a clase y pasan el día en la calle. Estas conductas pueden llevar a la adicción y a la delincuencia. La familia y la escuela deben trabajar juntas para ayudar a estos jóvenes.
\end{textebox}

\courssub{Glossaire des mots nouveaux (vocabulaire clé en marge)}

Les mots ci-dessous sont les termes « transparents » (proches du français) et les termes « opaques » (à mémoriser prioritairement). L'enseignant les écrit au tableau en deux colonnes et demande aux élèves de deviner le sens avant de donner la traduction.

\begin{center}
\begin{tabularx}{\linewidth}{|>{\bfseries}l|X|>{\itshape}X|}
\hline
\rowcolor{popblueL}
\textbf{Español} & \textbf{Français} & \textbf{Note phonétique / Contexte} \\
\hline
\esp{instituto} & lycée (enseignement secondaire) & Synonyme : \esp{colegio} (souvent privé), \esp{IES} (Instituto de Educación Secundaria). \\
\esp{conducta de riesgo} & comportement à risque & \esp{conducta} = conduite, comportement (fém.). \\
\esp{adoptar} & adopter, prendre (une habitude) & Verbe régulier en \esp{-ar}. \esp{Adoptan} (3e pl. présent). \\
\esp{fumar} & fumer & \esp{Fuman} (3e pl.). \esp{Tabaco} = tabac (masc.). \\
\esp{patio} & cour de récréation & Attention : \esp{patio} $\neq$ patio (français = terrasse). \\
\esp{beber} & boire (de l'alcool surtout) & \esp{Beben} (3e pl.). \esp{Alcohol} = alcool (masc., accent sur \emph{o}). \\
\esp{consumir} & consommer & \esp{Consumen} (3e pl.). Régulier en \esp{-ir}. \\
\esp{droga} & drogue, stupéfiant & \textbf{Féminin} : \esp{la droga}, \esp{las drogas}. \\
\esp{acoso escolar} & harcèlement scolaire / bullying & \esp{Acoso} (masc.) de \esp{acosar} (harceler). \\
\esp{aumentar} & augmenter, augmenter (intransitif) & \esp{Está aumentando} (présent progressif). \\
\esp{insultar} & insulter & \esp{Insulta} (3e sg.). \esp{Un alumno insulta a otro}. \\
\esp{redes sociales} & réseaux sociaux & \esp{Red} (fém.) = réseau, filet. \\
\esp{sexting} & sexting (anglicisme intégré) & Invariant. Pratique d'envoi de contenus sexuels par mobile. \\
\esp{consecuencia} & conséquence & Cognat transparent. \esp{Pensar en las consecuencias}. \\
\esp{absentismo escolar} & absentéisme scolaire & \esp{Absentismo} (masc.). Préoccupation majeure. \\
\esp{preocupar} & préoccuper, inquiéter & Verbe à construction inverse : \esp{Me preocupa} = « Il me préoccupe / Je m'inquiète ». \\
\esp{calle} & rue & \esp{En la calle} = dans la rue, dehors. \\
\esp{adicción} & addiction, dépendance & \esp{La adicción} (fém.). De \esp{adicto} (accro). \\
\esp{delincuencia} & délinquance & \esp{La delincuencia} (fém.). \esp{Delincuente} = délinquant. \\
\hline
\end{tabularx}
\end{center}

\courssub{Reformulation orale guidée (production orale)}

L'enseignant guide la reformulation par des questions fermées puis ouvertes. Les élèves répondent en phrases simples (niveau A2+/B1). On accepte l'espagnol approximatif mais on recaste (reformulation correcte par le prof) systématiquement.

\begin{exemplebox}[Modèle de reformulation professeur-élèves]
\begin{itemize}
\item \textbf{Prof :} \esp{¿Qué pasa en los institutos?} \\
\textbf{Élève :} \esp{Los jóvenes adoptan conductas de riesgo.} \\
\textbf{Prof (recast) :} \esp{Exacto. En muchos institutos, los jóvenes están adoptando conductas de riesgo.}
\item \textbf{Prof :} \esp{¿Qué hacen en el patio?} \\
\textbf{Élève :} \esp{Fuman tabaco.} \\
\textbf{Prof :} \esp{Bien. Fuman tabaco en el patio.}
\item \textbf{Prof :} \esp{¿Y en las fiestas?} \\
\textbf{Élève :} \esp{Beben alcohol.} \\
\textbf{Prof :} \esp{Correcto. Beben alcohol durante las fiestas.}
\item \textbf{Prof :} \esp{¿Qué problema hay en las redes sociales?} \\
\textbf{Élève :} \esp{Acoso y sexting.} \\
\textbf{Prof :} \esp{El acoso escolar está aumentando. Un alumno insulta a otro o practica el sexting.}
\item \textbf{Prof :} \esp{¿Qué pasa con las clases?} \\
\textbf{Élève :} \esp{No vienen. Absentismo.} \\
\textbf{Prof :} \esp{Hay absentismo escolar. Unos chicos no vienen a clase y pasan el día en la calle.}
\item \textbf{Prof :} \esp{¿Cuáles son las consecuencias graves?} \\
\textbf{Élève :} \esp{Adicción y delincuencia.} \\
\textbf{Prof :} \esp{Estas conductas pueden llevar a la adicción y a la delincuencia.}
\item \textbf{Prof :} \esp{¿Quién debe ayudar?} \\
\textbf{Élève :} \esp{Familia y escuela.} \\
\textbf{Prof :} \esp{La familia y la escuela deben trabajar juntas.}
\end{itemize}
\end{exemplebox}

\courssub{Notion clé : Définition de « conducta de riesgo »}

\begin{definition}[Conducta de riesgo (Comportement à risque)]
Une \cle{conducta de riesgo} est tout comportement volontaire ou répété qui met en danger la santé physique ou mentale, la sécurité, la scolarité ou l'avenir social d'un jeune. Dans le contexte scolaire hispanophone, on classe généralement ces conduites en quatre grands domaines :
\begin{enumerate}
\item \textbf{Consommation de substances :} tabac (\esp{tabaco}), alcool (\esp{alcohol}), drogues illicites (\esp{drogas}).
\item \textbf{Violence et harcèlement :} harcèlement scolaire (\esp{acoso escolar}), cyberharcèlement (\esp{ciberacoso}), bagarres (\esp{peleas}), insultes (\esp{insultos}).
\item \textbf{Rupture scolaire :} absentéisme (\esp{absentismo}), décrochage (\esp{fracaso escolar / abandono escolar}).
\item \textbf{Risques numériques :} \esp{sexting}, \esp{grooming} (cyber-prédation), exposition à des contenus violents, jeu pathologique (\esp{ludopatía}).
\end{enumerate}
\end{definition}

\courssub{Lexique thématique : Tableau récapitulatif par champs sémantiques}

Ce tableau sert de « fiche de vocabulaire » à coller dans le cahier. Il sera réutilisé en Section 3 (Lexique) et en Section 6 (Production).

\begin{center}
\begin{tabularx}{\linewidth}{|>{\bfseries}l|X|X|}
\hline
\rowcolor{popblueL}
\textbf{Champ lexical} & \textbf{Español (avec article)} & \textbf{Français} \\
\hline
\textbf{Sustancias} & \esp{el tabaco}, \esp{el alcohol}, \esp{la droga}, \esp{la adicción} & Tabac, alcool, drogue, addiction \\
\hline
\textbf{Violencia} & \esp{el acoso (escolar/cibernético)}, \esp{la pelea}, \esp{el insulto}, \esp{la violencia} & Harcèlement (scolaire/cyber), bagarre, insulte, violence \\
\hline
\textbf{Escuela} & \esp{el absentismo (escolar)}, \esp{el fracaso escolar}, \esp{el abandono}, \esp{la clase}, \esp{el instituto} & Absentéisme, échec scolaire, décrochage, classe, lycée \\
\hline
\textbf{Redes sociales} & \esp{el sexting}, \esp{el ciberacoso}, \esp{el grooming}, \esp{las redes sociales}, \esp{el móvil} & Sexting, cyberharcèlement, grooming, réseaux sociaux, portable \\
\hline
\textbf{Conséquences} & \esp{la delincuencia}, \esp{la calle}, \esp{las consecuencias}, \esp{el futuro} & Délinquance, rue, conséquences, avenir \\
\hline
\textbf{Acteurs / Solutions} & \esp{la familia}, \esp{la escuela}, \esp{los profesores}, \esp{los padres}, \esp{trabajar juntos}, \esp{ayudar} & Famille, école, profs, parents, travailler ensemble, aider \\
\hline
\end{tabularx}
\end{center}

\courssub{Point de vigilance : Pièges de genre (faux amis morphologiques)}

\begin{attention}[Attention : Genres pièges pour francophones !]
\begin{itemize}
\item \textbf{\esp{La droga} (féminin).} En français « drogue » est féminin, mais beaucoup d'élèves disent \esp{*el droga} par analogie avec \esp{el problema} ou parce que le mot finit par \esp{-a}. \textbf{Règle :} Les noms en \esp{-a} sont généralement féminins (\esp{la casa, la mesa, la droga}). Exception notable : \esp{el problema, el tema, el sistema, el programa, el clima, el mapa, el día, el idioma}.
\item \textbf{\esp{El problema} (masculin).} Malgré sa finale en \esp{-a}, c'est un nom masculin d'origine grecque (comme \esp{el tema, el sistema}). On dit \esp{un problema grave}, \esp{los problemas}.
\item \textbf{\esp{El absentismo} (masculin).} Finit par \esp{-o} $\rightarrow$ masculin. Facile.
\item \textbf{\esp{La adicción} (féminin).} Finit par \esp{-ión} $\rightarrow$ toujours féminin (\esp{la nación, la televisión, la lección}).
\item \textbf{\esp{La delincuencia} (féminin).} Finit par \esp{-cia} $\rightarrow$ féminin.
\end{itemize}
\textbf{Astuce mnémotechnique :} \emph{« La DRÔLE de DRogue est Féminine ; Le PROblème est Masculin comme un PROfesseur. »}
\end{attention}

\courssub{Le saviez-vous ? : Le harcèlement dans l'espace hispanophone}

\begin{savaistu}[Acoso escolar, bullying, ciberacoso : des mots pour une même réalité]
\begin{itemize}
\item En \textbf{Espagne}, le terme officiel dans l'éducation est \esp{acoso escolar}. L'anglicisme \esp{bullying} est très utilisé dans les médias et par les jeunes (\esp{« Sufro bullying »}). La loi organique 2/2006 (LOE) et la LOMLOE (2020) obligent les centres à avoir un \esp{plan de convivencia} et un \esp{coordinador de bienestar}.
\item Au \textbf{Mexique}, on parle couramment d'\esp{acoso escolar} et d'\esp{ciberacoso} (ou \esp{acoso cibernético}) pour le harcèlement en ligne. La \esp{Ley General de los Derechos de Niñas, Niños y Adolescentes} (2014) protège explicitement contre ces violences.
\item En \textbf{Guinée équatoriale} (seul pays africain hispanophone), le terme \esp{acoso escolar} s'impose dans les textes officiels (Ministère de l'Éducation), mais le français « harcèlement » reste très utilisé dans l'oral quotidien à Malabo et Bata, de même que les langues locales (fang, bubi).
\item Au \textbf{Cameroun}, pays bilingue FR/EN, le phénomène existe sous les termes « harcèlement scolaire » / « bullying ». L'apprentissage de l'espagnol permet aux élèves camerounais de nommer ces réalités dans une troisième langue et d'accéder à des ressources pédagogiques hispanophones (guides de l'UNICEF Espagne, campagnes \esp{« No al acoso»}).
\end{itemize}
\end{savaistu}

\courssub{Carte mentale : Classification visuelle des conduites à risque}

\begin{popfigure}
\begin{tikzpicture}[
  mindmap,
  concept color=popblue!80,
  level 1 concept/.append style={level distance=4.5cm, sibling angle=90, font=\bfseries\small},
  level 2 concept/.append style={level distance=3cm, sibling angle=45, font=\small},
  every node/.style={align=center, text width=2.8cm},
  grow cyclic
]
\node[concept, align=center]{Conductas\\ de riesgo}
  child[concept color=poporange!80] { node[concept] {Sustancias}
    child { node {Tabaco} }
    child { node {Alcohol} }
    child { node {Drogas} }
  }
  child[concept color=popgreen!80] { node[concept] {Violencia}
    child { node[align=center]{Acoso\\ escolar} }
    child { node {Ciberacoso} }
    child { node[align=center]{Peleas /\\ Insultos} }
  }
  child[concept color=poppurple!80] { node[concept] {Escuela}
    child { node {Absentismo} }
    child { node[align=center]{Fracaso\\ escolar} }
    child { node {Abandono} }
  }
  child[concept color=poppink!80] { node[concept, align=center]{Redes\\ sociales}
    child { node {Sexting} }
    child { node {Grooming} }
    child { node[align=center]{Adicción\\ a pantallas} }
  };
\end{tikzpicture}
\legende{Carte mentale des conduites à risque (vocabulaire clé de la leçon EA03).}
\end{popfigure}

\courssub{Grammaire : Le présent de l'indicatif et le gérondif (estar + gerundio)}

Le texte support utilise deux formes du présent pour décrire la réalité : le \textbf{présent de l'indicatif} (valeur habituelle ou gnomique) et le \textbf{présent progressif} (\esp{estar + gerundio}, valeur processuelle, action en cours). C'est un point clé du programme de Première (LV2) pour décrire une situation actuelle.

\begin{propriete}[Présent de l'indicatif : Formation et emplois principaux]
\textbf{1. Formation (réguliers) :} On ajoute les terminaisons au radical.
\begin{center}
\begin{tabular}{|l|c|c|c|}
\hline
\rowcolor{popblueL}
\textbf{Personne} & \textbf{-AR (hablar)} & \textbf{-ER (beber)} & \textbf{-IR (consumir)} \\
\hline
yo & habl\textbf{o} & beb\textbf{o} & consum\textbf{o} \\
tú & habl\textbf{as} & beb\textbf{es} & consum\textbf{es} \\
él/ella/usted & habl\textbf{a} & beb\textbf{e} & consum\textbf{e} \\
nosotros/as & habl\textbf{amos} & beb\textbf{emos} & consum\textbf{imos} \\
vosotros/as & habl\textbf{áis} & beb\textbf{éis} & consum\textbf{ís} \\
ellos/ellas/ustedes & habl\textbf{an} & beb\textbf{en} & consum\textbf{en} \\
\hline
\end{tabular}
\end{center}
\textbf{2. Verbes irréguliers fréquents dans le texte :}
\begin{itemize}
\item \esp{estar} (yo estoy, tú estás, él está, nosotros estamos, vosotros estáis, ellos \textbf{están}) $\rightarrow$ auxiliaire du progressif.
\item \esp{tener} (yo tengo, tú tienes, él tiene...) $\rightarrow$ \esp{tener que + inf.} (obligation).
\item \esp{poder} (yo puedo, tú puedes, él puede...) $\rightarrow$ \esp{pueden llevar} (capacité/possibilité).
\item \esp{deber} (yo debo, tú debes, él debe...) $\rightarrow$ \esp{deben trabajar} (devoir moral).
\item \esp{venir} (yo vengo, tú vienes, él viene, nosotros venimos, vosotros venís, ellos \textbf{vienen}) $\rightarrow$ \esp{no vienen a clase} (changement de radical \emph{e} $\rightarrow$ \emph{ie}).
\item \esp{preocupar} (verbe comme \esp{gustar}) : \esp{preocupa} (3e sg.), \esp{preocupan} (3e pl.).
\end{itemize}
\textbf{3. Emplois (niveau A2+/B1) :}
\begin{enumerate}
\item \textbf{Habitude, fait permanent :} \esp{Fuman tabaco en el patio.} (Ils fument / Ils ont l'habitude de fumer).
\item \textbf{Vérité générale, fait gnomique :} \esp{El alcohol daña la salud.} (L'alcool nuit à la santé).
\item \textbf{Action présente, narration :} \esp{Un alumno insulta a otro.} (Un élève insulte un autre — en ce moment / dans le récit).
\item \textbf{Futur proche (contexte clair) :} \esp{Mañana vienen los padres.} (Demain les parents viennent).
\end{enumerate}
\end{propriete}

\begin{propriete}[Le gérondif (Gerundio) et la périfrase \textbf{Estar + Gerundio}]
\textbf{1. Formation du gérondif (réguliers) :}
\begin{itemize}
\item Verbes en \esp{-AR} : radical + \textbf{-ando} $\rightarrow$ \esp{hablar $\rightarrow$ hablando}, \esp{fumar $\rightarrow$ fumando}, \esp{estudiar $\rightarrow$ estudiando} (attention \emph{i} + \emph{ando} = \emph{yendo} pour les voyelles fortes).
\item Verbes en \esp{-ER / -IR} : radical + \textbf{-iendo} $\rightarrow$ \esp{beber $\rightarrow$ bebiendo}, \esp{consumir $\rightarrow$ consumiendo}, \esp{aumentar $\rightarrow$ aumentando}.
\end{itemize}
\textbf{2. Irrégularités majeures (changement vocalique radical) :}
\begin{itemize}
\item \esp{e $\rightarrow$ i} : \esp{pedir $\rightarrow$ pidiendo}, \esp{seguir $\rightarrow$ siguiendo}, \esp{servir $\rightarrow$ sirviendo}, \esp{repetir $\rightarrow$ repitiendo}.
\item \esp{o $\rightarrow$ u} : \esp{dormir $\rightarrow$ durmiendo}, \esp{morir $\rightarrow$ muriendo}, \esp{poder $\rightarrow$ pudiendo}.
\item \esp{-yendo} (lecture forcée) : \esp{leer $\rightarrow$ leyendo}, \esp{oir $\rightarrow$ oyendo}, \esp{construir $\rightarrow$ construyendo}, \esp{huir $\rightarrow$ huyendo}, \esp{influir $\rightarrow$ influyendo}.
\end{itemize}
\textbf{3. Structure \esp{Estar + Gerundio} (Présent progressif) :}
\begin{center}
\begin{tabular}{|l|l|}
\hline
\rowcolor{popblueL}
\textbf{Sujet + Estar (présent) + Gérondif} & \textbf{Sens} \\
\hline
\esp{Estoy estudiando.} & Je suis en train d'étudier (maintenant). \\
\esp{Los jóvenes están adoptando conductas de riesgo.} & Les jeunes sont en train d'adopter / adoptent actuellement des conduites à risque. \\
\esp{El acoso está aumentando.} & Le harcèlement est en train d'augmenter / augmente actuellement. \\
\hline
\end{tabular}
\end{center}
\textbf{4. Différence clé \textbf{Presente} vs \textbf{Estar + Gerundio} :}
\begin{itemize}
\item \esp{Fuman en el patio} = Ils fument (habitude, fait constant, description générale).
\item \esp{Están fumando en el patio} = Ils sont en train de fumer (action en cours, visible maintenant, insistance sur la durée).
\item \textbf{Dans le texte :} \esp{están adoptando} (processus en cours, phénomène actuel qui se développe), \esp{está aumentando} (évolution en cours). Les autres verbes (\esp{fuman, beben, consumen, insulta, practica, vienen, pasan}) sont à l'indicatif simple pour décrire des faits constatés, des habitudes ou des situations permanentes.
\end{itemize}
\end{propriete}

\begin{attention}[Piège fréquent : \textbf{Estar + Gerundio} $\neq$ « Être en train de » pour tous les verbes d'état]
On n'utilise généralement \textbf{PAS} \esp{estar + gerundio} avec les verbes d'état, de perception, de sentiment ou de possession : \esp{saber, conocer, querer, preferir, tener (possession), parecer, gustar}.
\begin{itemize}
\item \textbf{Incorrect :} \esp{*Estoy sabiendo la lección.} $\rightarrow$ \textbf{Correct :} \esp{Sé la lección.}
\item \textbf{Incorrect :} \esp{*Está teniendo un coche.} $\rightarrow$ \textbf{Correct :} \esp{Tiene un coche.}
\item \textbf{Correct (changement de sens) :} \esp{Está teniendo problemas} (Il est en train d'avoir / traverse des problèmes — événementiel).
\end{itemize}
\end{attention}

\courssub{Exercice résolu : Questions de compréhension globale et détaillée}

\begin{exoresolu}[Compréhension du texte « Jóvenes en riesgo »]
\textbf{Énoncé :} Relis le texte et réponds en espagnol par des phrases complètes aux questions suivantes.
\begin{enumerate}
\item ¿Qué tipo de texto es? ¿Dónde se publicaría?
\item ¿Cuáles son las tres sustancias que consumen los jóvenes según el texto?
\item ¿Dónde fuman tabaco los estudiantes?
\item ¿Qué dos fenómenos vinculados a las redes sociales menciona el artículo?
\item ¿Por qué preocupa el absentismo a los profesores?
\item ¿A qué pueden llevar estas conductas según la penúltima frase?
\item ¿Quiénes deben trabajar juntos para ayudar a los jóvenes?
\end{enumerate}
\tcblower
\textbf{Solution détaillée :}
\begin{enumerate}
\item \esp{Es un artículo de prensa (reportaje social) sobre la juventud. Se publicaría en una revista educativa o un periódico escolar.} (Type : article de presse / reportage. Lieu : revue éducative ou journal scolaire.)
\item \esp{Las tres sustancias son: el tabaco, el alcohol y las drogas.} (Tabac, alcool, drogues.)
\item \esp{Fuman tabaco en el patio (del instituto).} (Dans la cour du lycée.)
\item \esp{Menciona el acoso escolar (insultar en redes sociales) y el sexting.} (Harcèlement scolaire — insulter sur les réseaux — et sexting.)
\item \esp{Porque hay chicos que no vienen a clase y pasan el día en la calle.} (Parce que des garçons ne viennent pas en cours et passent la journée dans la rue.)
\item \esp{Pueden llevar a la adicción y a la delincuencia.} (Peuvent mener à l'addiction et à la délinquance.)
\item \esp{La familia y la escuela deben trabajar juntas.} (La famille et l'école doivent travailler ensemble.)
\end{enumerate}
\textbf{Note méthodologique :} Les réponses reprennent le vocabulaire du texte (reprise lexicale). Pour la question 4, on distingue bien \esp{acoso escolar} (le phénomène global) et sa manifestation numérique \esp{insultar en redes sociales}, ainsi que \esp{sexting} (pratique spécifique).
\end{exoresolu}

\courssub{Mini-tâche intermédiaire : Annoter le texte (devoirs / travail en classe)}

\begin{exercice}[Annotation guidée du texte support]
Sur votre exemplaire papier ou sur la version numérique (OnBuch+), effectuez les annotations suivantes au crayon ou via l'outil surligneur :
\begin{enumerate}
\item Entourez en \textbf{bleu} tous les verbes conjugués au présent de l'indicatif (ex. \esp{están adoptando, fuman, beben, consumen, está aumentando, insulta, practica, preocupa, vienen, pasan, pueden llevar, deben trabajar}).
\item Entourez en \textbf{vert} tous les noms de \textbf{conductas de riesgo} (\esp{tabaco, alcohol, drogas, acoso escolar, sexting, absentismo}).
\item Soulignez en \textbf{rouge} les \textbf{lieux} (\esp{institutos, patio, fiestas, redes sociales, clase, calle}).
\item Mettez une \textbf{vaguelette} sous les \textbf{acteurs} (\esp{jóvenes, alumno, otro, profesores, chicos, familia, escuela}).
\item Écrivez en marge, en français, l'idée principale de chaque phrase (il y a 7 phrases).
\end{enumerate}
\end{exercice}

\begin{corrige}[Exercice : Annotation guidée]
\begin{enumerate}
\item \textbf{Verbes (présent indicatif / présent progressif) :} \esp{están adoptando} (estar+ger.), \esp{fuman}, \esp{beben}, \esp{consumen}, \esp{está aumentando} (estar+ger.), \esp{insulta}, \esp{practica}, \esp{preocupa}, \esp{vienen}, \esp{pasan}, \esp{pueden llevar}, \esp{deben trabajar}.
\item \textbf{Conductas (noms) :} \esp{tabaco, alcohol, drogas, acoso escolar, sexting, absentismo (escolar), adicción, delincuencia}.
\item \textbf{Lieux :} \esp{institutos, patio, fiestas, redes sociales, clase, calle}.
\item \textbf{Acteurs :} \esp{jóvenes, alumno, otro, profesores, chicos, familia, escuela}.
\item \textbf{Idées principales (résumé séquentiel) :}
      \begin{enumerate}
      \item Constat : adoption de conduites à risque dans les lycées.
      \item Exemples substances : tabac (cour), alcool (fêtes), drogues.
      \item Hausse du harcèlement : insultes en ligne, sexting.
      \item Problème absentéisme : élèves dans la rue au lieu de la classe.
      \item Conséquences graves : addiction, délinquance.
      \item Appel à l'action : alliance famille-école.
      \end{enumerate}
\end{enumerate}
\end{corrige}

\courssub{Méthode : Commentaire de texte (niveau Première)}

Le commentaire de texte est une épreuve type du Baccalauréat (épreuve écrite). En Première, on en apprend la structure rigoureuse en 4 étapes. On s'entraîne d'abord sur des textes courts comme celui-ci.

\begin{methode}[Étapes du commentaire de texte (Texte argumentatif / Article de presse)]
\begin{enumerate}
\item \textbf{Introduction (1 paragraphe) :}
      \begin{itemize}
      \item \textbf{Présenter le document :} Nature (article de presse), titre (\emph{« Jóvenes en riesgo »}), source fictive (\emph{Escuela Abierta}), thème général (conduites à risque des jeunes).
      \item \textbf{Problématique :} Une question ouverte qui guide l'analyse. Ex. : \emph{« ¿De qué manera el texto presenta las conductas de riesgo de los jóvenes y qué soluciones propone? »} (De quelle manière le texte présente-t-il les conduites à risque des jeunes et quelles solutions propose-t-il ?).
      \item \textbf{Annonce du plan :} \emph{« En primer lugar, analizaremos las conductas descritas; en segundo lugar, las causas y consecuencias; finalmente, la solución propuesta. »}
      \end{itemize}
\item \textbf{Développement (2 à 3 paragraphes structurés) :} Chaque paragraphe = une idée force + citation courte + analyse lexicale/grammaticale.
      \begin{itemize}
      \item \textbf{Axe 1 : L'inventaire des conduites à risque (Le « Quoi »).} Relevé des verbes d'action (\esp{fuman, beben, consumen, insulta, practica}) et du lexique des substances/violence. Distinction lieu/acte (\esp{patio/fiestas/redes}).
      \item \textbf{Axe 2 : L'aggravation et les conséquences (Le « Pourquoi » et le « Résultat »).} Analyse du progressif \esp{están adoptando, está aumentando} (processus en cours). Conséquences : \esp{adicción, delincuencia}. Verbes modaux : \esp{pueden llevar, deben trabajar}.
      \item \textbf{Axe 3 : Les acteurs et la solution (Le « Qui » et le « Comment »).} Opposition \esp{chicos/profesores/familia/escuela}. Solution collective : \esp{trabajar juntas}.
      \end{itemize}
\item \textbf{Conclusion (1 paragraphe) :}
      \begin{itemize}
      \item Réponse synthétique à la problématique.
      \item Ouverture : lien avec le contexte camerounais / Guinée équatoriale / rôle de l'éducation aux médias.
      \end{itemize}
\item \textbf{Relecture :} Vérifier l'orthographe (accents, \esp{ñ}, \esp{¡¿}), la cohérence des temps (présent de l'indicatif pour commenter), les connecteurs logiques (\esp{por tanto, además, sin embargo, en conclusión}).
\end{enumerate}
\end{methode}

\courssub{Production écrite guidée : Décrire une situation de risque}

\begin{exercice}[Rédaction d'un paragraphe descriptif (80-100 mots)]
\textbf{Consigne :} Vous êtes délégué(e) de classe dans un lycée de Douala ou Yaoundé. Vous écrivez un court rapport pour le journal du lycée \emph{« La Voz del Liceo »} afin d'alerter sur une conduite à risque que vous observez (choisir : absentéisme, harcèlement sur WhatsApp, consommation d'alcool aux abords de l'école, usage excessif du téléphone en classe).

\textbf{Contraintes obligatoires :}
\begin{enumerate}
\item Utiliser au moins \textbf{5 mots} du tableau de lexique (Section 8).
\item Utiliser \textbf{2 verbes au présent progressif (estar + gerundio)} pour décrire des actions en cours.
\item Utiliser \textbf{3 verbes au présent de l'indicatif} (habitude/constat).
\item Utiliser \textbf{un verbe modal (poder, deber, tener que, necesitar) + infinitif}.
\item Structurer : Situation $\rightarrow$ Exemples précis $\rightarrow$ Conséquence $\rightarrow$ Appel à l'action.
\end{enumerate}
\end{exercice}

\begin{corrige}[Exercice : Production écrite guidée (Proposition de modèle)]
\textbf{Titre :} \esp{El absentismo y el móvil: riesgos en nuestro patio}

\esp{En nuestro instituto de Douala, \textbf{estamos observando} un aumento del absentismo escolar. Muchos \textbf{chicos no vienen} a clase y \textbf{pasan} la mañana en la \textbf{calle} o en los cibercafés cercanos. \textbf{Están usando} el \textbf{móvil} para jugar o chatear en lugar de estudiar. El \textbf{profesor} \textbf{está preocupado} porque el \textbf{fracaso escolar} \textbf{puede llevar} a la \textbf{delincuencia}. \textbf{Debemos trabajar} juntos: \textbf{familia}, \textbf{escuela} y alumnos para cambiar esta \textbf{conducta de riesgo}.}

\textbf{Analyse du modèle :}
\begin{itemize}
\item \textbf{Lexique (7 items) :} absentismo escolar, chicos, calle, móvil, profesor, fracaso escolar, delincuencia, conducta de riesgo, familia, escuela.
\item \textbf{Estar + Gerundio (2) :} \esp{estamos observando, están usando} (aussi \esp{está preocupado} participe/adjectif, mais \esp{está aumentando} serait progressif).
\item \textbf{Présent indicatif (3+) :} \esp{no vienen, pasan, puede llevar, debemos trabajar}.
\item \textbf{Modal + Infinitif :} \esp{puede llevar, debemos trabajar}.
\item \textbf{Structure :} Respectée.
\end{itemize}
\end{corrige}

\coursec{Compréhension du texte}

Après la lecture du texte support \emph{« Jóvenes en riesgo »} (section précédente), nous entrons dans la phase active de compréhension. L'objectif n'est pas seulement de « comprendre le sens général », mais de \cle{maîtriser les stratégies de lecture} exigées au baccalauréat : repérage d'informations explicites, inférence, analyse lexicale et synthèse des idées. Chaque exercice ci-dessous est calibré sur le niveau A2+/B1 du CECRL et suit la grille d'évaluation officielle MINESEC.

\begin{textebox}[Texte support : Jóvenes en riesgo]
En los alrededores del liceo de Ngoa-Ekellé, \textbf{en} Yaoundé, como en los barrios populares de Douala o de Madrid, la realidad de los jóvenes presenta caras preocupantes. Cada mañana, profesores y orientadores observan conductas que ponen en peligro el futuro de los alumnos. El \cle{consumo} de \cle{alcohol} y de \cle{tabaco} comienza a edades muy tempranas, a veces desde los trece años, durante los recreos o a la salida de las clases. Más grave aún, la \cle{droga} —especialmente el tramadol y el cannabis— circula discretamente cerca de los establecimientos escolares; algunos estudiantes \cle{consumen} estas sustancias para « olvidar los problemas » o « tener valor ».

El \cle{acoso escolar} y el \cle{ciberacoso} (el \cle{sexting} incluido) destruyen la autoestima de las víctimas. Los agresores actúan en grupo, graban las humillaciones con sus teléfonos y difunden los vídeos en las redes sociales. El \cle{absentismo} escolar aumenta de forma alarmante: muchos chicos y chicas faltan a clase para pasar horas en los cibercafés, en las salas de juegos o simplemente deambulando por las calles. La \cle{violencia} física y verbal se vuelve cotidiana; las peleas entre bandas rivales y los robos a mano armada —la \cle{delincuencia} juvenil— son ya una realidad en varios barrios.

Ante este panorama, la \cle{adicción} no debe verse como un vicio, sino como una enfermedad que requiere atención médica y psicológica. Las familias, a menudo desbordadas por la precariedad económica, no siempre logran supervisar a sus hijos. La escuela, primer espacio de socialización, debe reforzar la \cle{prevención} y la escucha. Los profesores, los padres, los líderes comunitarios y las autoridades sanitarias tienen la responsabilidad compartida de \cle{identificar} las señales de alarma —cambio brusco de comportamiento, notas en caída libre, aislamiento— y de actuar antes de que sea demasiado tarde. El futuro de la juventud camerunesa e hispana se juega hoy, en cada aula y en cada calle.
\end{textebox}

\begin{definition}[Glossaire des termes clés du texte]
\begin{itemize}
    \item \esp{los alrededores} : les alentours, les environs.
    \item \esp{ponen en peligro} : mettent en danger (verbe \esp{poner en peligro}).
    \item \esp{el recreo} : la récréation.
    \item \esp{el tramadol} : le tramadol (antalgique opioïde détourné comme drogue, très répandu en Afrique centrale).
    \item \esp{la autoestima} : l'estime de soi.
    \item \esp{humillaciones} : humiliations.
    \item \esp{difunden} : ils diffusent (verbe \esp{difundir}, présent indicatif).
    \item \esp{deambulando} : errant, flânant (gérondif de \esp{deambular}).
    \item \esp{a mano armada} : à main armée.
    \item \esp{desbordadas} : débordées, dépassées (participe passé/adjectif de \esp{desbordar}).
    \item \esp{señales de alarma} : signaux d'alarme, signaux d'alerte.
    \item \esp{en caída libre} : en chute libre.
\end{itemize}
\end{definition}

\courssub{Compréhension globale — Questions ouvertes}

Répondez en français, par des phrases complètes, en vous appuyant sur le texte.

\begin{enumerate}
    \item \textbf{¿De qué habla el texto?} Quel est le thème général du texte ? Résumez en une phrase.
    \item \textbf{¿Dónde ocurren estas conductas?} Citez au moins trois lieux précis mentionnés dans le texte (milieu scolaire et milieu social).
    \item \textbf{¿Quiénes son los actores principales?} Identifiez les victimes, les auteurs et les adultes responsables cités.
    \item \textbf{¿Cuál es la tesis del autor?} Que l'auteur dénonce-t-il et quelle solution propose-t-il implicitement ?
\end{enumerate}

\begin{exoresolu}[Modèle de réponse à la question 1]
\textbf{Énoncé :} ¿De qué habla el texto? \\
\textbf{Solution détaillée :}
\begin{enumerate}
    \item \textbf{Repérage :} Le premier paragraphe introduit le sujet (\esp{la realidad de los jóvenes presenta caras preocupantes}), les paragraphes 2 et 3 détaillent les conduites (consommation, harcèlement, absentisme, violence, délinquance), le dernier paragraphe appelle à l'action (\esp{actuar antes de que sea demasiado tarde}).
    \item \textbf{Synthèse :} Le texte décrit les différents comportements à risque observés chez les jeunes (consommation de substances, violence, harcèlement, absentisme, délinquance) dans leur environnement scolaire et social, et appelle à une mobilisation collective (famille, école, autorités) pour la prévention et la prise en charge.
    \item \textbf{Réponse finale :} Le texte parle des conduites à risque des jeunes (drogue, alcool, violence, harcèlement, absentisme) à l'école et dans la rue, et de la nécessité d'une action concertée des adultes pour les prévenir.
\end{enumerate}
\end{exoresolu}

\courssub{Compréhension détaillée — Questions ciblées}

Les réponses doivent être justifiées par une \cle{citation textuelle} (copier-coller le segment exact du texte entre guillemets espagnols « »).

\begin{enumerate}
    \item \textbf{¿Qué consumen los jóvenes y a qué edad comienzan?}
    \item \textbf{¿Qué es el absentismo según el texto?} Expliquez avec vos mots et citez le passage.
    \item \textbf{¿Cómo actúan los agresores en el acoso escolar?} (Trois actions précises).
    \item \textbf{¿Quiénes deben ayudar a los jóvenes?} Listez les quatre groupes d'acteurs cités au dernier paragraphe.
    \item \textbf{¿Por qué las familias no siempre logran supervisar a sus hijos?}
    \item \textbf{¿Qué señales de alarma deben identificar los adultos?} (Citez-en trois).
\end{enumerate}

\begin{exoresolu}[Modèle de réponse à la question 3]
\textbf{Énoncé :} ¿Cómo actúan los agresores en el acoso escolar? \\
\textbf{Solution détaillée :}
\begin{enumerate}
    \item \textbf{Repérage :} Paragraphe 2, phrase 2 : \esp{« Los agresores actúan en grupo, graban las humillaciones con sus teléfonos y difunden los vídeos en las redes sociales. »}
    \item \textbf{Extraction des trois actions :}
        \begin{enumerate}
            \item Actúan en grupo (ils agissent en groupe).
            \item Graban las humillaciones con sus teléfonos (ils filment les humiliations avec leurs téléphones).
            \item Difunden los vídeos en las redes sociales (ils diffusent les vidéos sur les réseaux sociaux).
        \end{enumerate}
    \item \textbf{Réponse rédigée :} Les agresseurs agissent en groupe, filment les humiliations avec leurs téléphones et diffusent les vidéos sur les réseaux sociaux. \emph{(Citation : « actúan en grupo, graban las humillaciones con sus teléfonos y difunden los vídeos en las redes sociales »)}.
\end{enumerate}
\end{exoresolu}

\courssub{Vocabulaire en contexte — Recherche et déduction}

\begin{enumerate}
    \item Retrouvez dans le texte le nom masculin qui signifie « \textbf{harcèlement} » (paragraphe 2). \rightarrow \esp{\_\_\_\_\_\_\_\_\_\_}
    \item Retrouvez le nom féminin qui signifie « \textbf{addiction / dépendance} » (paragraphe 3). \rightarrow \esp{\_\_\_\_\_\_\_\_\_\_}
    \item Quel verbe à l'infinitif signifie « \textbf{identifier / repérer} » ? (Dernier paragraphe). \rightarrow \esp{\_\_\_\_\_\_\_\_\_\_}
    \item \textbf{Champ lexical de la consommation :} Relevez 4 mots/expressions du texte liés à la prise de substances (paragraphe 1).
    \item \textbf{Famille de mots :} À partir de \esp{adicción}, écrivez l'adjectif correspondant (\esp{adicto / adicta}) et le verbe (\esp{adictar} n'existe pas ; on utilise \esp{ser adicto a} ou \esp{engancharse a}). Donnez un exemple : \esp{Es adicto al tramadol} = « Il est accro au tramadol ».
\end{enumerate}

\begin{attention}[Piège de traduction : « Consumir » ne se traduit pas toujours par « consommer »]
En français, « consommer » est souvent neutre (consommer des fruits, de l'électricité). En espagnol, \esp{consumir} appliqué aux substances (drogue, alcool, tabac) implique une notion d'abus, de prise régulière et nocive.
\begin{itemize}
    \item \esp{Consumen drogas} $\neq$ « Ils consomment de la drogue » (trop neutre).
    \item \esp{Consumen drogas} = \textbf{« Ils se droguent » / « Ils prennent de la drogue » / « Ils font usage de stupéfiants »}.
    \item \esp{El consumo de alcohol} = \textbf{« La consommation d'alcool »} (terme statistique/sociologique) OU \textbf{« L'alcoolisation »} (contexte médical).
\end{itemize}
\emph{Conseil :} Au bac, en version, privilégiez « se droguer », « boire de l'alcool », « fumer » selon le contexte, plutôt que le mot générique « consommer ».
\end{attention}

\courssub{Vrai ou Faux ? Justifiez par une citation du texte}

Pour chaque affirmation, cochez V (Vrai) ou F (Faux) et recopiez \textbf{la phrase exacte du texte} qui justifie votre choix.

\begin{center}
\begin{tabularx}{\linewidth}{|X|c|X|}
\hline
\rowcolor{popblueL}
\textbf{Affirmation} & \textbf{V / F} & \textbf{Justification (Citation textuelle)} \\
\hline
1. La consommation de drogue commence généralement à l'âge adulte. & & \\
\hline
2. Le sexting est une forme de cyberharcèlement mentionnée dans le texte. & & \\
\hline
3. L'absentisme scolaire diminue grâce aux cibercafés. & & \\
\hline
4. La délinquance juvénile se limite aux bagarres sans armes. & & \\
\hline
5. L'auteur considère l'addiction comme un vice moral. & & \\
\hline
6. Les professeurs font partie des acteurs qui doivent agir. & & \\
\hline
\end{tabularx}
\end{center}

\begin{corrige}[Exercice Vrai/Faux]
1. \textbf{Faux.} Citation : \esp{« El consumo de alcohol y de tabaco comienza a edades muy tempranas, a veces desde los trece años »} / \esp{« la droga [...] circula discretamente cerca de los establecimientos escolares »}.
2. \textbf{Vrai.} Citation : \esp{« El acoso escolar y el ciberacoso (el sexting incluido) destruyen la autoestima de las víctimas »}.
3. \textbf{Faux.} Citation : \esp{« El absentismo escolar aumenta de forma alarmante »} (il augmente, ne diminue pas) + \esp{« faltan a clase para pasar horas en los cibercafés »} (cause de l'absentisme).
4. \textbf{Faux.} Citation : \esp{« los robos a mano armada —la delincuencia juvenil— son ya una realidad »}.
5. \textbf{Faux.} Citation : \esp{« la adicción no debe verse como un vicio, sino como una enfermedad »}.
6. \textbf{Vrai.} Citation : \esp{« Los profesores, los padres, los líderes comunitarios y las autoridades sanitarias tienen la responsabilidad compartida »}.
\end{corrige}

\courssub{Dégagement de l'idée principale et des idées secondaires}

\begin{methode}[Comment dégager l'idée principale d'un texte argumentatif / de société]
\begin{enumerate}
    \item \textbf{Lisez le titre et le premier paragraphe :} Ils annoncent souvent le thèse (\esp{la tesis}).
    \item \textbf{Identifiez le sujet (de qué habla) et la thèse (qué dice el autor) :} Sujet = conduites à risque. Thèse = c'est un danger grave qui exige une action collective urgente.
    \item \textbf{Repérez les connecteurs logiques :} \esp{Más grave aún} (aggravation), \emph{Ante este panorama} (conséquence/solution), \esp{no debe verse como..., sino como...} (rectification/définition), \esp{deben actuar} (obligation).
    \item \textbf{Formulez l'idée principale en une phrase nominale ou verbale :} « L'auteur dénonce l'ampleur des conduites à risque chez les jeunes et appelle à une mobilisation collective pour la prévention et la prise en charge. »
    \item \textbf{Classez les idées secondaires par paragraphe :} Chaque paragraphe = une facette du problème ou une partie de la solution.
\end{enumerate}
\end{methode}

\begin{exemplebox}[Application au texte « Jóvenes en riesgo »]
\begin{itemize}
    \item \textbf{Idée principale :} \emph{La multiplicité des conduites à risque (consommation, violence, harcèlement, absentisme, délinquance) menace la jeunesse camerounaise et hispanophone, nécessitant une réponse urgente et concertée de tous les acteurs éducatifs et sociaux.}
    \item \textbf{Idée secondaire 1 (§1) :} \emph{La consommation précoce et diversifiée de substances licites (alcool, tabac) et illicites (tramadol, cannabis) dans l'environnement scolaire immédiat.}
    \item \textbf{Idée secondaire 2 (§2) :} \emph{La violence sous toutes ses formes (harcèlement physique et numérique, absentisme, délinquance armée) gangrène le quotidien des jeunes.}
    \item \textbf{Idée secondaire 3 (§3) :} \emph{La nécessité de reconsidérer l'addiction comme une pathologie et d'engager une responsabilité partagée (famille, école, communauté, santé) pour la détection précoce et l'action.}
\end{itemize}
\end{exemplebox}

\courssub{Champ lexical de la dangerosité et de la protection}

Relevez dans le texte les mots appartenant au champ lexical du \textbf{risque, du danger, de la conséquence négative} (en \textcolor{popred}{rouge}) et ceux du \textbf{remède, de la protection, de l'action positive} (en \textcolor{popgreen}{vert}).

\begin{center}
\begin{tabularx}{\linewidth}{|>{\columncolor{popcream}}X|>{\columncolor{popred!15}}X|>{\columncolor{popgreen!15}}X|}
\hline
\rowcolor{popblueL}
\textbf{Paragraphe} & \textbf{Champ lexical : DANGER / RISQUE / NÉGATIF} & \textbf{Champ lexical : PROTECTION / SOLUTION / POSITIF} \\
\hline
§ 1 & \esp{preocupantes, peligro, consumo, tempranas, grave, droga, circula discretamente, consumen, problemas} & \esp{observan (vigilance)} \\
\hline
§ 2 & \esp{acoso, ciberacoso, destruyen, humillaciones, agresores, absentismo, aumenta, alarmante, violencia, peleas, robos, delincuencia} & \esp{—} \\
\hline
§ 3 & \esp{adicción, enfermedad, precariedad, desbordadas, señales de alarma, cambio brusco, caída libre, aislamiento, demasiado tarde} & \esp{atención médica, psicológica, prevención, escucha, responsabilidad compartida, identificar, actuar, futuro} \\
\hline
\end{tabularx}
\end{center}

\begin{aretenir}[Observation lexicale]
Le texte construit une \cle{progressivité dramatique} : le vocabulaire du danger s'accumule (§1 et §2) pour créer l'urgence, puis le vocabulaire de la solution apparaît massivement au §3 (\esp{prevención, escucha, responsabilidad, identificar, actuar}). C'est une structure classique du discours de plaidoyer (\emph{advocacy}) : \textbf{Constat alarmant $\rightarrow$ Analyse des causes $\rightarrow$ Appel à l'action}.
\end{aretenir}

% --- FIGURE TIKZ : ORGANIGRAMME DE COMPRÉHENSION ---
\begin{popfigure}
\begin{tikzpicture}[
    node distance=1.2cm and 2cm,
    main/.style={draw=popblue, thick, rounded corners, fill=popblue!10, text width=3.2cm, align=center, font=\small\bfseries, minimum height=1.2cm},
    sub/.style={draw=poporange, thick, rounded corners, fill=poporange!10, text width=2.8cm, align=center, font=\small, minimum height=1cm},
    conc/.style={draw=popgreen, thick, rounded corners, fill=popgreen!10, text width=3.5cm, align=center, font=\small\bfseries, minimum height=1.2cm},
    arrow/.style={-Stealth, thick, popdark}
]
% Nœuds
\node[main] (titre) {Texte : « Jóvenes en riesgo »};
\node[main, below=of titre] (theme) {Thème : Conduites à risque des jeunes (milieu scolaire \& social)};
\node[main, below=of theme] (ideeP) {Idée principale : Dénonciation de l'ampleur du phénomène \& appel à la mobilisation collective};
\node[sub, below left=of ideeP] (id1) {Idée secondaire 1 : Consommation précoce substances (licites/illicites)};
\node[sub, below=of ideeP] (id2) {Idée secondaire 2 : Violence, harcèlement, absentisme, délinquance};
\node[sub, below right=of ideeP] (id3) {Idée secondaire 3 : Addiction = maladie ; responsabilité partagée \& prévention};
\node[conc, below=1.5cm of id2] (concl) {Conclusion / Synthèse : Urgence éducative et sanitaire — Agir maintenant};

% Flèches
\draw[arrow] (titre) -- (theme);
\draw[arrow] (theme) -- (ideeP);
\draw[arrow] (ideeP) -- (id1);
\draw[arrow] (ideeP) -- (id2);
\draw[arrow] (ideeP) -- (id3);
\draw[arrow] (id1) -- (concl);
\draw[arrow] (id2) -- (concl);
\draw[arrow] (id3) -- (concl);

% Accolade décorative
\draw[popmuted, thick, decorate, decoration={brace, amplitude=5pt}] (id1.north west) -- (id3.north east) node[midway, above=6pt, font=\tiny\bfseries, popdark] {Corps du texte (3 paragraphes)};
\end{tikzpicture}
\legende{Organigramme de la structure logique du texte : du titre à la synthèse.}
\end{popfigure}

\courssub{Entraînement guidé — Production écrite courte}

\begin{exercice}[Rédiger un paragraphe descriptif (80-100 mots)]
\textbf{Consigne :} En vous appuyant sur le texte, le lexique acquis et la structure « Idée principale / Idées secondaires », rédigez en espagnol un paragraphe cohérent pour présenter la situation des jeunes à risque à un correspondant espagnol. Utilisez le \textbf{présent de l'indicatif} et au moins \textbf{deux formes en \esp{estar + gerundio}} (ex. \esp{están consumiendo, está aumentando}).

\textbf{Mots-outils obligatoires :} \esp{actualmente, además, por eso, deben, urgente}.
\textbf{Contraintes :} 3 idées secondaires distinctes, 1 phrase de conclusion appelant à l'action.
\end{exercice}

\begin{corrige}[Exercice production écrite — Proposition de correction]
\textbf{Texte modèle :}
\esp{Actualmente, la situación de los jóvenes en Yaoundé y en otras ciudades es muy preocupante. En primer lugar, el consumo de alcohol, tabaco y drogas como el tramadol comienza muy temprano, a veces a los trece años. Además, el acoso escolar y el ciberacoso están destruyendo la autoestima de muchos estudiantes; los agresores graban y difunden vídeos en redes sociales. Por otro lado, el absentismo está aumentando porque los alumnos faltan a clase para ir a cibercafés. Por eso, la adicción no es un vicio, es una enfermedad. Los profesores, los padres y las autoridades deben identificar las señales de alarma y actuar ya. Es urgente proteger el futuro de la juventud.}

\textbf{Analyse grammaticale (pour l'enseignant) :}
\begin{itemize}
    \item \esp{es, comienza, destruyen, graban, difunden, faltan, es, deben, identificar, actuar, proteger} = Présent indicatif (valeur descriptive, vérité générale, actualité).
    \item \esp{están destruyendo, está aumentando} = \esp{Estar + gerundio} (valeur aspectuelle : processus en cours, situation dynamique actuelle).
    \item Connecteurs : \esp{Actualmente, En primer lugar, Además, Por otro lado, Por eso, Es urgente}.
\end{itemize}
\end{corrige}

\begin{methode}[Réussir le commentaire de texte (Bac LV2) — Étape 1 : La compréhension]
Avant d'écrire l'introduction ou le développement, \textbf{remplissez cette fiche de lecture} sur votre brouillon :
\begin{enumerate}
    \item \textbf{Nature du document :} Article de presse / Témoignage / Rapport / Texte d'opinion ? (Ici : Texte d'opinion / Plaidoyer).
    \item \textbf{Source / Auteur / Date :} (Si fournis). Contexte : Cameroun / Monde hispanophone.
    \item \textbf{Thème général :} (Un nom ou groupe nominal). Ex: \emph{Las conductas de riesgo juveniles}.
    \item \textbf{Problématique implicite :} Pourquoi les jeunes sont-ils en danger et que faire ? (\esp{¿Por qué los jóvenes están en riesgo y qué hay que hacer?}).
    \item \textbf{Thèse de l'auteur :} (Une phrase affirmative). Ex: \emph{El autor denuncia la gravedad de las conductas de riesgo y exige una acción colectiva inmediata.}
    \item \textbf{Plan du texte (2 ou 3 parties) :} 1. Constat : consommation \& violence. 2. Conséquences : absentisme \& délinquance. 3. Solutions : prévention \& responsabilité partagée.
    \item \textbf{Champs lexicaux majeurs :} Danger (2 cols) / Solution (1 col). Repérez les \emph{champs lexicaux antagonistes}.
\end{enumerate}
\emph{Cette fiche vaut 4 à 6 points sur 20 à la compréhension écrite au Bac. Ne la négligez pas.}
\end{methode}

\begin{savaistu}[Le Cameroun et la Guinée équatoriale : un espace hispanophone africain]
La Guinée équatoriale (capital : Malabo) est le seul pays africain où l'espagnol est langue officielle (avec le français et le portugais). Elle fait partie de la CEMAC (Communauté Économique et Monétaire de l'Afrique Centrale) comme le Cameroun. De nombreux jeunes Camerounais vont y étudier ou y travailler ; inversement, des Guinéens équatoriaux étudient à Yaoundé ou Douala. Les \cle{conductas de riesgo} décrites dans le texte (tramadol, absentisme, bandes, sexting) sont des réalités partagées des deux côtés de la frontière. La coopération éducative et sanitaire entre les deux pays est un enjeu majeur pour la \cle{prevención} en zone frontalière (Kye-Ossi, Ebebiyín).
\end{savaistu}

\coursec{Lexique thématique par champs}

Après avoir lu et compris le texte \emph{« Jóvenes en riesgo »}, nous allons maintenant organiser le vocabulaire essentiel par champs sémantiques. Cette méthode permet de mémoriser les mots par réseaux de sens, de les réutiliser à l'oral comme à l'écrit, et de repérer les nuances entre l'espagnol et le français. Chaque champ sera illustré par un tableau comparatif, des exemples de phrases authentiques et des points de vigilance pour les francophones.

\courssub{Champ 1 — Les substances et la consommation}

\begin{definition}[Les substances : vocabulaire de base]
Les noms des substances sont généralement des noms communs \textbf{masculins} en espagnol (sauf \esp{la droga}, féminin). Le verbe générique pour « consommer » est \esp{consumir} ; les verbes spécifiques sont \esp{fumar} (tabac) et \esp{beber} (alcool).
\end{definition}

\begin{center}
\begin{tabularx}{\linewidth}{|X|X|X|}
\hline
\rowcolor{popblueL}
\textbf{Español} & \textbf{Français} & \textbf{Remarque / Prononciation} \\
\hline
\esp{la droga} & la drogue & Féminin. Pluriel : \esp{las drogas}. Pron. « dro-ga ». \\
\esp{el alcohol} & l'alcool & Masculin. Accent tonique sur le \emph{o} final : « al-co-\textbf{ol} ». \\
\esp{el tabaco} & le tabac & Masculin. Pron. « ta-\textbf{ba}-ko ». \\
\esp{fumar} & fumer & Verbe régulier en \emph{-ar}. \esp{Yo fumo, tú fumas…} \\
\esp{beber} & boire & Verbe régulier en \emph{-er}. \esp{Yo bebo, tú bebes…} \\
\esp{consumir} & consommer & Verbe régulier en \emph{-ir}. Plus formel ; englobe drogue, alcool, médicaments. \\
\hline
\end{tabularx}
\end{center}

\begin{exemplebox}[Phrases contextuelles]
\begin{itemize}
\item \esp{Muchos jóvenes \textbf{consumen} alcohol los fines de semana.} « Beaucoup de jeunes consomment de l'alcool le week-end. »
\item \esp{Está prohibido \textbf{fumar} en el recinto del instituto.} « Il est interdit de fumer dans l'enceinte du lycée. »
\item \esp{El \textbf{tabaco} y el \textbf{alcohol} son drogas legales.} « Le tabac et l'alcool sont des drogues légales. »
\item \esp{No \textbf{bebas} alcohol si vas a conducir.} « Ne bois pas d'alcool si tu vas conduire. » (Impératif négatif = subjonctif présent, 2\up{e} pers. du singulier.)
\end{itemize}
\end{exemplebox}

\begin{attention}[Pièges pour francophones]
\begin{itemize}
\item \textbf{Faux ami :} \esp{droga} désigne couramment la « substance illicite / le stupéfiant ». Pour « médicament », on dit \esp{medicamento} ou \esp{fármaco}, jamais \esp{droga} dans ce sens.
\item \textbf{Genre :} on dit \esp{el alcohol} (masculin), jamais *\esp{la alcohol}. Attention à ne pas confondre avec le cas de \esp{el agua} / \esp{el águila} : ces noms-là sont \textbf{féminins} et prennent \esp{el} uniquement parce qu'ils commencent par un \emph{a} tonique (\esp{el agua fría}). \esp{Alcohol}, lui, est masculin par nature : \esp{el alcohol puro}.
\item \textbf{\emph{Beber} vs \emph{tomar} :} \esp{tomar} (« prendre ») s'emploie beaucoup pour « boire un verre » : \esp{tomar una cerveza} « boire une bière ». \esp{beber} est plus neutre ou plus clinique.
\end{itemize}
\end{attention}

\begin{popfigure}
\begin{tikzpicture}[
    node distance=1.2cm and 1.5cm,
    main/.style={draw=popblue, thick, rounded corners, fill=popblueL, text width=3.5cm, align=center, font=\small},
    sub/.style={draw=poporange, rounded corners, fill=poporangeL, text width=2.8cm, align=center, font=\small\itshape},
    arrow/.style={-Stealth, popblue, thick}
]
\node[main, align=center] (centre){\textbf{SUSTANCIAS}\\\emph{Consumo}};
\node[sub, align=center] (droga) [above left=of centre]{la droga\\(drogues illicites)};
\node[sub, align=center] (alcohol) [above=of centre]{el alcohol\\(boissons)};
\node[sub, align=center] (tabaco) [above right=of centre]{el tabaco\\(cigarettes)};
\node[sub] (fumar) [below left=of centre] {fumar};
\node[sub] (beber) [below=of centre] {beber / tomar};
\node[sub, align=center] (consumir) [below right=of centre]{consumir\\(terme général)};
\draw[arrow] (centre) -- (droga);
\draw[arrow] (centre) -- (alcohol);
\draw[arrow] (centre) -- (tabaco);
\draw[arrow] (fumar) -- (centre);
\draw[arrow] (beber) -- (centre);
\draw[arrow] (consumir) -- (centre);
\end{tikzpicture}
\legende{Carte mentale : substances et verbes de consommation}
\end{popfigure}

\courssub{Champ 2 — La violence, le harcèlement et le sexting}

\begin{definition}[Violence et harcèlement : termes clés]
La \esp{violencia} est le terme générique. L'\esp{acoso escolar} (harcèlement scolaire) implique une répétition et un déséquilibre de pouvoir. Le \esp{sexting} est un anglicisme lexicalisé en espagnol (masculin : \esp{el sexting}). Les verbes d'action sont \esp{insultar}, \esp{pegar}, \esp{acosar}.
\end{definition}

\begin{center}
\begin{tabularx}{\linewidth}{|X|X|X|}
\hline
\rowcolor{poporangeL}
\textbf{Español} & \textbf{Français} & \textbf{Remarque} \\
\hline
\esp{la violencia} & la violence & Féminin. \esp{Violencia física / psicológica / de género}. \\
\esp{el acoso escolar} & le harcèlement scolaire & Masculin. \esp{acoso} = harcèlement ; \esp{escolar} = scolaire. \\
\esp{el ciberacoso} & le cyberharcèlement & Aussi \esp{ciberbullying} (anglicisme). \\
\esp{el sexting} & le sexting & \textbf{Anglicisme}, genre masculin. Voir \emph{Attention} ci-dessous. \\
\esp{insultar} & insulter & Régulier en \emph{-ar}. \esp{Me insulta constantemente.} \\
\esp{pegar} & frapper, donner des coups & Régulier en \emph{-ar}. Aussi « coller ». \esp{Le pegó una bofetada.} « Il lui a donné une gifle. » \\
\esp{la pelea} & la bagarre, la dispute & Féminin. \esp{Tener una pelea.} « Se disputer. » \\
\esp{acosar} & harceler & Verbe souche de \esp{acoso}. \esp{No acoses a tus compañeros.} \\
\hline
\end{tabularx}
\end{center}

\begin{attention}[Anglicisme : « sexting »]
\esp{El sexting} s'utilise tel quel dans la presse, les campagnes officielles (ex. \esp{campaña contra el sexting}) et le langage courant des jeunes. \textbf{N'inventez pas le verbe *\esp{sextear} ni le nom *\esp{sexteo} dans un écrit formel ou scolaire.} Pour dire « faire du sexting », on utilise une périphrase : \esp{enviar mensajes o imágenes de contenido sexual} ou \esp{practicar sexting}.
\end{attention}

\begin{exemplebox}[Situations réelles]
\begin{itemize}
\item \esp{El \textbf{acoso escolar} puede ser físico, verbal o \textbf{ciberacoso}.} « Le harcèlement scolaire peut être physique, verbal ou prendre la forme de cyberharcèlement. »
\item \esp{Los profesores detectaron una \textbf{pelea} en el patio.} « Les professeurs ont repéré une bagarre dans la cour. »
\item \esp{El \textbf{sexting} entre menores es un delito en muchos países.} « Le sexting entre mineurs est un délit dans beaucoup de pays. »
\item \esp{No debes \textbf{insultar} a nadie por su origen o religión.} « Tu ne dois insulter personne en raison de son origine ou de sa religion. »
\end{itemize}
\end{exemplebox}

\begin{savaistu}[Cyberharcèlement en Espagne, en Amérique latine et au Cameroun]
Le terme \esp{ciberacoso} (ou \esp{ciberbullying}) est très utilisé. En Espagne, la loi organique 8/2021 de protection de l'enfance (LOPIVI) protège les mineurs contre la violence numérique. Au Cameroun, la loi n°2010/012 du 21 décembre 2010 relative à la cybersécurité et à la cybercriminalité sanctionne aussi le harcèlement en ligne. En Guinée équatoriale, seul pays hispanophone d'Afrique, les réalités scolaires sont proches de celles du Cameroun (classes nombreuses, forte présence du téléphone portable).
\end{savaistu}

\courssub{Champ 3 — Le milieu scolaire : absentéisme et échec}

\begin{definition}[Vocabulaire de l'institution scolaire]
En Espagne et en Amérique latine, le lycée s'appelle \esp{el instituto} (ou \esp{instituto de educación secundaria}). L'élève est \esp{el alumno} (ou \esp{la alumna}). \esp{Suspender} signifie « échouer » (à un examen, à une matière) ; \esp{expulsar} signifie « exclure » (mesure disciplinaire).
\end{definition}

\begin{center}
\begin{tabularx}{\linewidth}{|X|X|X|}
\hline
\rowcolor{popgreenL}
\textbf{Español} & \textbf{Français} & \textbf{Contexte / Nuance} \\
\hline
\esp{el absentismo (escolar)} & l'absentéisme (scolaire) & Masculin. Souvent précédé de l'article : \esp{el absentismo escolar}. \\
\esp{el instituto} & le lycée, le collège public & \esp{Colegio} = école (souvent primaire ou privée) ; \esp{instituto} = secondaire public. \\
\esp{el alumno / la alumna} & l'élève & Synonyme : \esp{el/la estudiante} (plus formel, aussi niveau supérieur). \\
\esp{suspender} & échouer, être recalé & \esp{Suspendí matemáticas.} « J'ai raté les maths. » (En Espagne : note inférieure à 5 sur 10.) \\
\esp{expulsar} & exclure (sanction) & \esp{Le expulsaron del instituto por un mes.} « Il a été exclu du lycée pour un mois. » \\
\esp{repetir (curso)} & redoubler (l'année) & \esp{Repetir curso} = redoubler. \esp{Repetir una asignatura} = refaire une matière. \\
\hline
\end{tabularx}
\end{center}

\begin{exemplebox}[Parcours scolaires]
\begin{itemize}
\item \esp{El \textbf{absentismo escolar} aumenta en los últimos cursos de la ESO.} « L'absentéisme scolaire augmente dans les dernières années du collège (ESO). »
\item \esp{Mi hermano \textbf{repitió} segundo de la ESO.} « Mon frère a redoublé la deuxième année de l'ESO. »
\item \esp{Si \textbf{suspendes} tres asignaturas, tienes que \textbf{repetir curso}.} « Si tu rates trois matières, tu dois redoubler. »
\item \esp{El \textbf{consejero} del instituto habla con las familias de los alumnos con absentismo.} « Le conseiller d'orientation du lycée parle avec les familles des élèves en situation d'absentéisme. »
\end{itemize}
\end{exemplebox}

\begin{attention}[Faux amis : \emph{suspender} vs « suspendre »]
En français, « suspendre » un élève = l'exclure temporairement.
En espagnol, \esp{suspender} (un examen, une matière) = \textbf{échouer}. Pour « exclure temporairement », on utilise \esp{expulsar} (souvent \esp{expulsión temporal}). \textbf{Ne traduisez jamais « il a été suspendu (du lycée) » par *\esp{ha sido suspendido} (qui voudrait dire « il a échoué / été recalé »), mais par \esp{ha sido expulsado}.}
\end{attention}

\courssub{Champ 4 — Conséquences : addiction, délinquance, danger}

\begin{definition}[Noms abstraits des conséquences]
Ces noms sont essentiels pour analyser, commenter un texte ou rédiger un paragraphe argumentatif. Beaucoup sont féminins en \emph{-ción} ou \emph{-ncia} (équivalents français en \emph{-tion} / \emph{-nce}) : \esp{adicción}, \esp{delincuencia}, \esp{consecuencia}. \esp{Riesgo} et \esp{peligro} sont masculins.
\end{definition}

\begin{center}
\begin{tabularx}{\linewidth}{|X|X|X|}
\hline
\rowcolor{poppurpleL}
\textbf{Español} & \textbf{Français} & \textbf{Collocations fréquentes} \\
\hline
\esp{la adicción} & l'addiction, la dépendance & \esp{Adicción a las drogas / al juego / al móvil}. \esp{Caer en la adicción}. \\
\esp{la delincuencia} & la délinquance & \esp{Delincuencia juvenil}. \esp{Luchar contra la delincuencia}. \\
\esp{el riesgo} & le risque & \esp{Correr un riesgo}. \esp{Estar en riesgo}. \esp{Factor de riesgo}. \\
\esp{el peligro} & le danger & \esp{Estar en peligro}. \esp{Peligro de muerte}. Plus fort que \esp{riesgo}. \\
\esp{las consecuencias} & les conséquences & \esp{Sufrir las consecuencias}. \esp{Consecuencias graves / legales / psicológicas}. \\
\esp{el daño} & le dommage, le tort & \esp{Causar daño}. \esp{Daño cerebral / moral / físico}. \\
\hline
\end{tabularx}
\end{center}

\begin{exemplebox}[Phrases d'analyse]
\begin{itemize}
\item \esp{La \textbf{adicción} al móvil es un nuevo \textbf{riesgo} para la salud mental.} « L'addiction au portable est un nouveau risque pour la santé mentale. »
\item \esp{El \textbf{absentismo} escolar es a menudo el primer paso hacia la \textbf{delincuencia} juvenil.} « L'absentéisme scolaire est souvent le premier pas vers la délinquance juvénile. »
\item \esp{Los jóvenes no miden las \textbf{consecuencias} de sus actos.} « Les jeunes ne mesurent pas les conséquences de leurs actes. »
\item \esp{El \textbf{peligro} no es solo la droga, sino la \textbf{normalización} de su consumo.} « Le danger n'est pas seulement la drogue, mais la normalisation de sa consommation. »
\end{itemize}
\end{exemplebox}

\begin{aretenir}[Nuance : \emph{riesgo} vs \emph{peligro}]
\esp{Riesgo} = probabilité d'un mal (être exposé au risque). \esp{Peligro} = menace imminente, gravité plus grande.
\begin{itemize}
\item \esp{Conducir rápido es un \textbf{riesgo}.} « Conduire vite est un risque. » (Probabilité d'accident.)
\item \esp{¡\textbf{Peligro}! Alta tensión.} « Danger ! Haute tension. » (Menace immédiate, panneau d'avertissement.)
\end{itemize}
Dans les textes de prévention, on parle de \esp{factores de riesgo} (facteurs de risque) et de \esp{situaciones de peligro} (situations de danger).
\end{aretenir}

\courssub{Champ 5 — L'aide, la prévention et les acteurs}

\begin{definition}[Acteurs et actions de prévention]
La prévention (\esp{la prevención}) implique la famille (\esp{la familia}), l'école (\esp{el instituto}, \esp{el consejero}) et les institutions. Les verbes clés sont \esp{ayudar} (aider), \esp{prevenir} (prévenir), \esp{orientar} (orienter), \esp{detectar} (dépister), \esp{denunciar} (signaler).
\end{definition}

\begin{center}
\begin{tabularx}{\linewidth}{|X|X|X|}
\hline
\rowcolor{poppinkL}
\textbf{Español} & \textbf{Français} & \textbf{Exemple d'usage} \\
\hline
\esp{ayudar (a + inf.)} & aider (à + inf.) & \esp{Ayudo a mi hermano a estudiar.} « J'aide mon frère à étudier. » \\
\esp{prevenir} & prévenir & \esp{Prevenir es mejor que curar.} « Mieux vaut prévenir que guérir. » \\
\esp{la familia} & la famille & \esp{El papel de la familia es fundamental.} \\
\esp{el consejero / la consejera} & le/la conseiller(-ère) d'orientation & En Espagne, on dit aussi \esp{orientador/orientadora}. \\
\esp{la campaña de prevención} & la campagne de prévention & \esp{Campaña contra el acoso / el alcohol}. \\
\esp{detectar} & détecter, repérer & \esp{Los profesores detectan signos de acoso.} \\
\esp{denunciar} & signaler, porter plainte & \esp{Denunciar un caso de acoso.} (Démarche officielle.) \\
\esp{el teléfono de ayuda} & le numéro d'aide, la ligne d'écoute & Ex. Espagne : 016 (violences de genre). \\
\hline
\end{tabularx}
\end{center}

\begin{exemplebox}[Discours de prévention]
\begin{itemize}
\item \esp{La \textbf{familia} debe \textbf{ayudar} a los jóvenes a \textbf{prevenir} el consumo.} « La famille doit aider les jeunes à prévenir la consommation. »
\item \esp{El \textbf{consejero} del instituto organiza \textbf{campañas de prevención} cada trimestre.} « Le conseiller d'orientation organise des campagnes de prévention chaque trimestre. »
\item \esp{Si ves \textbf{acoso}, \textbf{denúncialo}. No te calles.} « Si tu vois du harcèlement, signale-le. Ne te tais pas. »
\item \esp{Los padres deben \textbf{detectar} cambios de conducta en sus hijos.} « Les parents doivent détecter les changements de comportement chez leurs enfants. »
\end{itemize}
\end{exemplebox}

\begin{savaistu}[Le conseiller d'orientation : \emph{orientador} vs \emph{consejero}]
En Espagne, le terme officiel est \esp{orientador/orientadora} (psychopédagogue). Au Cameroun et en Guinée équatoriale, on parle de \esp{consejero/consejera de orientación} ou simplement \esp{consejero}. Les deux sont corrects ; adaptez-vous au contexte du texte. Au Cameroun, le conseiller d'orientation joue un rôle clé dans le choix des séries (A, C, D, TI…) et le suivi des élèves en difficulté.
\end{savaistu}

\courssub{Verbes associés : locutions verbales clés}

Ces locutions figent le sens et sont indispensables pour décrire une trajectoire (entrée dans le risque, sortie de l'addiction).

\begin{propriete}[Locutions verbales : forme et emploi]
\begin{itemize}
\item \esp{Adoptar una conducta (de riesgo)} : « adopter un comportement (à risque) ». \esp{Adoptar} + nom. Registre formel, d'analyse.
\item \esp{Caer en la droga / en la adicción / en el vicio} : « tomber dans la drogue / l'addiction / le vice ». \esp{Caer en} + nom. Image de la chute, souvent involontaire, progressive.
\item \esp{Salir de la adicción / de la droga / del túnel} : « sortir de l'addiction / de la drogue / du tunnel ». \esp{Salir de} + nom. Processus actif, positif, long.
\item \esp{Entrar en el mundo de la droga} : « entrer dans le monde de la drogue ».
\item \esp{Dejar de + infinitif} (\esp{dejar de fumar}, \esp{dejar de beber}) : « arrêter de… ». \textbf{Attention :} \esp{dejar} seul = « quitter, laisser ». \esp{Dejar la droga} = « arrêter la drogue ».
\end{itemize}
\end{propriete}

\begin{exoresolu}[Emploi des locutions verbales]
\textbf{Énoncé :} Traduisez en espagnol les phrases suivantes en utilisant les locutions ci-dessus.
\begin{enumerate}
\item « Beaucoup de jeunes adoptent des comportements à risque à cause de la pression du groupe. »
\item « Il est tombé dans la drogue à 15 ans. »
\item « Elle a réussi à sortir de l'addiction grâce à sa famille. »
\item « Il a arrêté de fumer l'an dernier. »
\end{enumerate}
\tcblower
\textbf{Solution détaillée :}
\begin{enumerate}
\item \esp{Muchos jóvenes \textbf{adoptan conductas de riesgo} por la presión del grupo.} (\esp{conductas} au pluriel car plusieurs comportements possibles ; \esp{presión del grupo} = pression des pairs ; \esp{por} exprime la cause.)
\item \esp{\textbf{Cayó en la droga} a los 15 años.} (\esp{cayó} = prétérit indéfini, 3\up{e} pers. du singulier de \esp{caer} ; \esp{a los 15 años} = à 15 ans.)
\item \esp{Logró \textbf{salir de la adicción} gracias a su familia.} (\esp{logró} = elle a réussi, elle parvint ; \esp{salir de} + nom ; \esp{gracias a} = grâce à.)
\item \esp{\textbf{Dejó de fumar} el año pasado.} (\esp{dejó de} + infinitif = il arrêta de ; \esp{el año pasado} = l'an dernier.)
\end{enumerate}
\end{exoresolu}

\begin{attention}[Erreurs fréquentes : \emph{dejar} vs \emph{dejar de}]
\begin{itemize}
\item \esp{Dejar la droga} = arrêter la drogue (objet direct = la chose qu'on abandonne).
\item \esp{Dejar de fumar} = arrêter de fumer (\textbf{de} + infinitif = cesser une action).
\item \textbf{Ne dites jamais :} *\esp{dejar fumar} (cela signifierait « laisser fumer, permettre de fumer »).
\end{itemize}
De même : \esp{dejar de beber} (arrêter de boire) vs \esp{dejar el alcohol} (renoncer à l'alcool).
\end{attention}

\courssub{Synthèse visuelle : milieu scolaire vs milieu social}

Les conduites à risque ne se manifestent pas au même endroit ni de la même façon. Ce schéma vous aide à classer le vocabulaire pour l'expression orale ou écrite (description d'image, commentaire).

\begin{popfigure}
\begin{tikzpicture}[
    zone/.style={draw=popdark, thick, rounded corners, fill=popcream, text width=6.5cm, align=left, minimum height=7cm},
    title/.style={font=\bfseries\large, text=popblue, anchor=north},
    item/.style={font=\small, anchor=west},
    arrow/.style={-Stealth, poporange, thick, dashed}
]
% Deux colonnes
\node[zone] (scolaire) at (0,0) {};
\node[zone] (social) at (8,0) {};

% Titres
\node[title, align=center] at (scolaire.north){MILIEU SCOLAIRE\\(instituto, aula, patio)};
\node[title, align=center] at (social.north){MILIEU SOCIAL\\(barrio, parque, redes, casa)};

% Items Scolaire
\begin{scope}[shift={(scolaire.center)}]
\node[item] (s1) at (-2.8, 2.5) {\textcolor{popblue}{\textbullet} \esp{Absentismo escolar}};
\node[item] (s2) at (-2.8, 1.7) {\textcolor{popblue}{\textbullet} \esp{Acoso escolar / ciberacoso}};
\node[item] (s3) at (-2.8, 0.9) {\textcolor{popblue}{\textbullet} \esp{Peleas en el patio / baños}};
\node[item] (s4) at (-2.8, 0.1) {\textcolor{popblue}{\textbullet} \esp{Suspender / repetir curso}};
\node[item] (s5) at (-2.8, -0.7) {\textcolor{popblue}{\textbullet} \esp{Expulsión temporal}};
\node[item] (s6) at (-2.8, -1.5) {\textcolor{popblue}{\textbullet} \esp{Consumo en el recreo (tabaco)}};
\node[item] (s7) at (-2.8, -2.3) {\textcolor{popblue}{\textbullet} \esp{Rol del consejero / orientador}};
\end{scope}

% Items Social
\begin{scope}[shift={(social.center)}]
\node[item] (so1) at (-2.8, 2.5) {\textcolor{poporange}{\textbullet} \esp{Consumo alcohol / drogas (botellón)}};
\node[item] (so2) at (-2.8, 1.7) {\textcolor{poporange}{\textbullet} \esp{Violencia de grupo / bandas}};
\node[item] (so3) at (-2.8, 0.9) {\textcolor{poporange}{\textbullet} \esp{Sexting / ciberacoso (redes)}};
\node[item] (so4) at (-2.8, 0.1) {\textcolor{poporange}{\textbullet} \esp{Delincuencia juvenil (robos)}};
\node[item] (so5) at (-2.8, -0.7) {\textcolor{poporange}{\textbullet} \esp{Adicción (sin control escolar)}};
\node[item] (so6) at (-2.8, -1.5) {\textcolor{poporange}{\textbullet} \esp{Familia: apoyo o factor de riesgo}};
\node[item] (so7) at (-2.8, -2.3) {\textcolor{poporange}{\textbullet} \esp{Campañas de prevención / tel. ayuda}};
\end{scope}

% Flèches de liaison
\draw[arrow] ($(scolaire.east)+(0,1.5)$) to[bend left=15] node[midway, above, font=\tiny, text=popdark] {Continuidad} ($(social.west)+(0,1.5)$);
\draw[arrow] ($(social.west)+(0,-1.5)$) to[bend left=15] node[midway, below, font=\tiny, text=popdark] {Repercusión} ($(scolaire.east)+(0,-1.5)$);
\end{tikzpicture}
\legende{Classification des conduites à risque : milieu scolaire vs milieu social (flèches = interactions)}
\end{popfigure}

\courssub{Texte d'entraînement : « La historia de Kevin »}

Ce texte original d'environ 230 mots réutilise l'essentiel du lexique des cinq champs. Lisez-le, repérez les mots étudiés, puis répondez aux questions.

\begin{textebox}[La historia de Kevin — Texte original pour l'entraînement]
Kevin tiene dieciséis años y estudia en un instituto público de las afueras de Yaoundé. Desde hace unos meses, su conducta ha cambiado mucho. Antes era un alumno aplicado, pero ahora presenta un \textbf{absentismo escolar} preocupante: falta a clase dos o tres veces por semana. Sus profesores han \textbf{detectado} que llega tarde, cansado, y a veces huele a \textbf{tabaco}.

En el \textbf{patio}, Kevin ya no juega con sus antiguos amigos. Se junta con un grupo de chicos mayores que él, que \textbf{fuman} y \textbf{beben} alcohol escondidos detrás del gimnasio. La semana pasada, hubo una \textbf{pelea} grave: Kevin \textbf{insultó} a un compañero y le \textbf{pegó} un puñetazo. El \textbf{consejero} del instituto ha citado a sus padres para una reunión urgente. Ellos están muy preocupados: temen que Kevin \textbf{caiga en la droga} y en la \textbf{delincuencia}.

Además, han descubierto mensajes de \textbf{sexting} en su móvil: imágenes íntimas intercambiadas con una chica de su barrio. Es un caso de \textbf{ciberacoso}, porque la chica lo amenaza con difundir las fotos si no le da dinero. Kevin siente \textbf{miedo} y \textbf{vergüenza}. No ve ninguna \textbf{salida}. El \textbf{consejero} les explica que sí hay \textbf{salida}: \textbf{denunciar} el chantaje, pedir \textbf{ayuda} psicológica y \textbf{prevenir} que la situación empeore. Kevin quiere \textbf{salir de la adicción} al tabaco y al alcohol, y volver a ser el chico de antes. Pero sabe que el \textbf{camino} es largo y que necesita el apoyo de su \textbf{familia} y de la escuela.
\end{textebox}

\begin{definition}[Glossaire des mots difficiles du texte]
\begin{itemize}
\item \esp{las afueras} : la périphérie, les faubourgs.
\item \esp{aplicado} : appliqué, sérieux, studieux.
\item \esp{preocupante} : inquiétant.
\item \esp{faltar a clase} : manquer les cours (être absent).
\item \esp{huele a tabaco} : il sent le tabac (verbe \esp{oler}, radical changeant \emph{o} $\to$ \emph{ue} : \esp{huelo, hueles, huele…}).
\item \esp{se junta con} : il traîne avec, il fréquente (verbe pronominal \esp{juntarse}).
\item \esp{escondidos} : cachés (participe passé employé comme adjectif, accord avec \esp{chicos} sous-entendu : « eux, cachés »).
\item \esp{un puñetazo} : un coup de poing.
\item \esp{citar} : convoquer.
\item \esp{temen que + subjonctif} : ils craignent que + subjonctif (\esp{caiga}, subjonctif de \esp{caer}).
\item \esp{chantaje} : chantage.
\item \esp{amenaza con + infinitif} : elle menace de + infinitif.
\item \esp{difundir} : diffuser, propager.
\item \esp{vergüenza} : honte.
\item \esp{salida} : issue, sortie (ici au sens figuré : solution).
\item \esp{empeorar} : empirer, s'aggraver.
\item \esp{camino} : chemin, parcours (au figuré : processus).
\end{itemize}
\end{definition}

\begin{exercice}[Questions de compréhension et de lexique]
\begin{enumerate}
\item Relevez dans le texte 4 noms du champ « Conséquences » (Champ 4) et 3 verbes du champ « Aide/Prévention » (Champ 5).
\item Trouvez dans le texte les deux locutions verbales étudiées (\esp{caer en…}, \esp{salir de…}) et traduisez les phrases où elles apparaissent.
\item Le texte mentionne \esp{absentismo escolar} et \esp{ciberacoso}. Expliquez en français la différence entre le milieu où ils se manifestent principalement (scolaire vs social/numérique) et leur lien possible.
\item \textbf{Production écrite (30-40 mots) :} imaginez la réponse de Kevin au conseiller. Écrivez 3 phrases en espagnol en utilisant : \esp{quiero dejar de…}, \esp{necesito ayuda para…}, \esp{voy a denunciar…}.
\end{enumerate}
\end{exercice}

\begin{corrige}[Exercice Lexique — La historia de Kevin]
\begin{enumerate}
\item \textbf{Noms du champ « Conséquences » :} \esp{la delincuencia}, \esp{la adicción}, \esp{el miedo}, \esp{la vergüenza} (on peut ajouter \esp{el chantaje}, forme de délinquance). \textbf{Verbes du champ « Aide/Prévention » :} \esp{detectar} (\esp{han detectado}), \esp{denunciar}, \esp{prevenir} (on peut ajouter \esp{pedir ayuda} et \esp{salir de}).
\item \textbf{Locutions :}
\begin{itemize}
\item \esp{« …temen que Kevin \textbf{caiga en la droga} y en la \textbf{delincuencia} ».} $\to$ « …ils craignent que Kevin tombe dans la drogue et la délinquance. » (Subjonctif \esp{caiga} après \esp{temer que} + changement de sujet.)
\item \esp{« Kevin quiere \textbf{salir de la adicción} al tabaco y al alcohol… »} $\to$ « Kevin veut sortir de l'addiction au tabac et à l'alcool… »
\end{itemize}
\item \textbf{Analyse :} l'\esp{absentismo escolar} se manifeste \textbf{dans l'enceinte de l'école} (absence physique aux cours). Le \esp{ciberacoso} (sexting, chantage) se déroule \textbf{sur les réseaux sociaux et le téléphone}, donc hors des murs, dans le \textbf{milieu social et numérique}. \textbf{Lien :} l'absentéisme isole l'élève, le fragilise et le pousse vers un groupe de pairs déviants du quartier ; le téléphone devient alors le lieu de prolongation du harcèlement et du chantage. L'école (le \esp{consejero}) est le point de détection et de sortie.
\item \textbf{Production écrite (exemple de corrigé) :}

\esp{Quiero \textbf{dejar de fumar} y \textbf{dejar de beber}. \textbf{Necesito ayuda para} salir de este problema y recuperar mi vida de antes. \textbf{Voy a denunciar} el chantaje con mis padres y con el consejero del instituto.}
\end{enumerate}
\end{corrige}

\begin{methode}[Comment réviser ce lexique pour l'évaluation ?]
\begin{enumerate}
\item \textbf{Par champs :} recopiez les 5 tableaux (Espagnol / Français / Collocation) sur des fiches.
\item \textbf{Par paires :} associez chaque nom à son verbe typique (\esp{adicción} $\leftrightarrow$ \esp{caer en / salir de} ; \esp{acoso} $\leftrightarrow$ \esp{detectar / denunciar / sufrir} ; \esp{absentismo} $\leftrightarrow$ \esp{presentar / combatir}).
\item \textbf{Contexte camerounais :} écrivez 5 phrases vraies sur votre lycée ou votre quartier en utilisant 5 mots différents de la leçon (ex. \esp{En mi instituto hay una campaña contra el acoso escolar.}).
\item \textbf{Pièges :} relisez les boîtes \textbf{Attention} (\esp{el alcohol}, \esp{suspender} vs \esp{expulsar}, \esp{dejar de} + infinitif). Ce sont des points faciles à gagner à l'examen.
\end{enumerate}
\end{methode}

\coursec{Grammaire : présent de l'indicatif et estar + gérondif}

Dans la section précédente, nous avons classé le vocabulaire des conduites à risque par champs lexicaux : \esp{droga}, \esp{alcohol}, \esp{tabaco}, \esp{violencia}, \esp{acoso escolar}, \esp{sexting}, \esp{absentismo}, \esp{delincuencia}, \esp{adicción}. Mais nommer un comportement ne suffit pas : il faut aussi savoir le \cle{décrire}. Or, pour décrire, l'espagnol distingue deux situations très différentes : ce qu'un jeune fait \emph{habituellement} et ce qu'il est \emph{en train de faire} au moment où l'on parle. C'est toute la différence entre \esp{fuma} et \esp{está fumando}. Cette section grammaticale va nous donner les deux outils indispensables : le \cle{présent de l'indicatif} et la périphrase \esp{estar + gerundio}.

\courssub{Observation : deux situations, deux formes verbales}

Lisons d'abord un court texte original qui reprend le lexique de la leçon et mélange volontairement les deux types de verbes. Soulignez mentalement les verbes conjugués.

\begin{textebox}[En el patio del instituto]
Son las diez de la mañana y suena el timbre del recreo. En el patio del instituto, los alumnos hablan, ríen y juegan al fútbol. Pero detrás del gimnasio, unos chicos fuman tabaco todos los días : es una costumbre peligrosa que preocupa a los profesores. Hoy, la situación es diferente : un alumno mayor está ofreciendo cigarrillos a los más jóvenes y otro está grabando la escena con su teléfono móvil. Cerca de la puerta, dos chicas están insultando a una compañera : es un caso claro de acoso escolar. La víctima llora en silencio mientras las agresoras se ríen. En la sala de informática, varios estudiantes visitan páginas web prohibidas ; normalmente estudian allí, pero hoy están compartiendo fotos comprometedoras, una práctica conocida como sexting. El director del centro camina por el pasillo en este momento : está observando todo y está tomando notas para hablar con las familias. Los especialistas dicen que los jóvenes adoptan conductas de riesgo cuando se sienten solos o presionados por el grupo. Por eso, el centro organiza cada semana charlas de prevención sobre la droga, el alcohol y la violencia. Mientras tanto, en el aula de tutoría, la profesora de español está explicando a sus alumnos cómo identificar estos problemas y cómo pedir ayuda. La prevención empieza siempre con la información y el diálogo.
\end{textebox}

\textbf{Glossaire :} \esp{el timbre} : la sonnerie ; \esp{el recreo} : la récréation ; \esp{una costumbre} : une habitude ; \esp{grabando} : en train de filmer ; \esp{comprometedoras} : compromettantes ; \esp{el pasillo} : le couloir ; \esp{tomando notas} : en train de prendre des notes ; \esp{charlas} : causeries, exposés ; \esp{pedir ayuda} : demander de l'aide ; \esp{el centro} : l'établissement (scolaire) ; \esp{un alumno mayor} : un élève plus âgé.

\textbf{Questions de compréhension :}
\begin{enumerate}
\item ¿Qué hacen normalmente los chicos detrás del gimnasio ?
\item ¿Qué está haciendo hoy un alumno mayor ?
\item ¿Qué está pasando cerca de la puerta ?
\item ¿Qué está haciendo el director en este momento ?
\end{enumerate}

Comparons maintenant deux phrases tirées du texte :

\begin{exemplebox}[Deux verbes, deux réalités]
\esp{Unos chicos \textbf{fuman} tabaco todos los días.} « Des garçons fument du tabac tous les jours. » $\rightarrow$ habitude, fait répété.\par
\esp{Un alumno \textbf{está ofreciendo} cigarrillos a los más jóvenes.} « Un élève est en train d'offrir des cigarettes aux plus jeunes. » $\rightarrow$ action en cours, en ce moment précis.
\end{exemplebox}

L'espagnol, comme le français, choisit sa forme verbale selon le sens : le \cle{présent simple} pour l'habitude et les faits généraux, la forme \esp{estar + gerundio} pour l'action en train de se dérouler.

\courssub{Le présent de l'indicatif : rappel des formes}

\begin{propriete}[Formation du présent de l'indicatif]
On prend le radical de l'infinitif et on ajoute les terminaisons propres à chaque groupe :
\begin{itemize}
\item verbes en \cle{-ar} (hablar) : \textbf{-o, -as, -a, -amos, -áis, -an}
\item verbes en \cle{-er} (comer) : \textbf{-o, -es, -e, -emos, -éis, -en}
\item verbes en \cle{-ir} (vivir) : \textbf{-o, -es, -e, -imos, -ís, -en}
\end{itemize}
\end{propriete}

\begin{center}
\begin{tabular}{|l|l|l|l|}
\hline
\rowcolor{popblueL} Personne & hablar (parler) & comer (manger) & vivir (vivre) \\ \hline
yo & habl\textbf{o} & com\textbf{o} & viv\textbf{o} \\ \hline
tú & habl\textbf{as} & com\textbf{es} & viv\textbf{es} \\ \hline
él / ella / usted & habl\textbf{a} & com\textbf{e} & viv\textbf{e} \\ \hline
nosotros & habl\textbf{amos} & com\textbf{emos} & viv\textbf{imos} \\ \hline
vosotros & habl\textbf{áis} & com\textbf{éis} & viv\textbf{ís} \\ \hline
ellos / ustedes & habl\textbf{an} & com\textbf{en} & viv\textbf{en} \\ \hline
\end{tabular}
\end{center}

\begin{attention}[L'accent de la 2e personne du pluriel]
Les formes \esp{habláis}, \esp{coméis}, \esp{vivís} portent toujours un accent écrit. N'oubliez pas non plus que le sujet est souvent sous-entendu en espagnol : \esp{Fuman} suffit pour dire « Ils fument », car la terminaison \esp{-an} indique déjà la personne.
\end{attention}

\begin{propriete}[Emplois du présent de l'indicatif]
Le présent simple exprime :
\begin{enumerate}
\item une \cle{habitude} ou une action répétée : \esp{Los jóvenes fuman en el patio.} « Les jeunes fument dans la cour. »
\item un \cle{fait général}, une vérité : \esp{El tabaco daña la salud.} « Le tabac nuit à la santé. »
\item la \cle{description} d'une situation durable : \esp{El instituto organiza charlas de prevención.} « Le lycée organise des causeries de prévention. »
\end{enumerate}
Il est souvent accompagné de marqueurs de fréquence : \esp{siempre} (toujours), \esp{todos los días} (tous les jours), \esp{normalmente} (normalement), \esp{a veces} (parfois), \esp{nunca} (jamais).
\end{propriete}

\courssub{Le gérondif : formation régulière et irrégulière}

Pour exprimer l'action en cours, il faut d'abord savoir former le \cle{gérondif} (\esp{gerundio}), l'équivalent du français « -ant » (fumant, mangeant).

\begin{propriete}[Formation du gérondif régulier]
\begin{itemize}
\item verbes en \cle{-ar} : radical + \textbf{-ando} $\rightarrow$ \esp{hablar} $\rightarrow$ \esp{hablando} (en train de parler)
\item verbes en \cle{-er} et \cle{-ir} : radical + \textbf{-iendo} $\rightarrow$ \esp{comer} $\rightarrow$ \esp{comiendo} ; \esp{vivir} $\rightarrow$ \esp{viviendo}
\end{itemize}
\end{propriete}

\begin{propriete}[Gérondifs irréguliers à connaître par cœur]
\begin{itemize}
\item Les verbes en \cle{-ir} à changement de voyelle : la voyelle du radical change (\textbf{e $\rightarrow$ i}, \textbf{o $\rightarrow$ u}) : \esp{decir} $\rightarrow$ \esp{diciendo}, \esp{pedir} $\rightarrow$ \esp{pidiendo}, \esp{dormir} $\rightarrow$ \esp{durmiendo}, \esp{mourir} $\rightarrow$ \esp{muriendo}, \esp{sentir} $\rightarrow$ \esp{sintiendo}, \esp{venir} $\rightarrow$ \esp{viniendo}, \esp{tenir} $\rightarrow$ \esp{teniendo}.
\item Les verbes dont le radical finit par une voyelle : \textbf{-iendo} devient \textbf{-yendo} : \esp{leer} $\rightarrow$ \esp{leyendo}, \esp{construir} $\rightarrow$ \esp{construyendo}, \esp{oír} $\rightarrow$ \esp{oyendo}, \esp{huir} $\rightarrow$ \esp{huyendo}, \esp{incluir} $\rightarrow$ \esp{incluyendo}.
\item Cas particuliers : \esp{ir} $\rightarrow$ \esp{yendo}, \esp{poder} $\rightarrow$ \esp{pudiendo}.
\end{itemize}
\end{propriete}

\begin{exemplebox}[Exemples traduits]
\esp{Está durmiendo en clase.} « Il est en train de dormir en classe. »\par
\esp{Están leyendo el artículo.} « Ils sont en train de lire l'article. »\par
\esp{La alumna está pidiendo ayuda.} « L'élève est en train de demander de l'aide. »\par
\esp{Estamos construyendo un futuro mejor.} « Nous sommes en train de construire un avenir meilleur. »
\end{exemplebox}

\courssub{La périphrase estar + gérondif : l'action en cours}

\begin{propriete}[estar + gérondif]
La périphrase \esp{estar + gerundio} exprime une \cle{action en train de se dérouler au moment où l'on parle}. Elle correspond exactement au français « être en train de + infinitif ».\par
\esp{Está fumando.} = « Il est en train de fumer. »\par
Le verbe \esp{estar} se conjugue au présent : \textbf{estoy, estás, está, estamos, estáis, están}. Le gérondif, lui, reste invariable.
\end{propriete}

\begin{popfigure}
\begin{tikzpicture}
% Frise du temps
\draw[-{Stealth}, thick, popdark] (-0.5,0) -- (11,0) node[right, font=\small] {tiempo};
\draw[thick, popdark] (5.2,-0.15) -- (5.2,0.15) node[above=2pt, font=\small] {ahora};
% Présent : habitudes répétées
\foreach \x in {0.8, 2.2, 3.6, 6.8, 8.2, 9.6} {
  \draw[popblue, thick] (\x,0.5) circle (0.09);
}
\node[popblue, font=\small, align=center] at (2.2,1.1) {presente : hábito\\ \esp{Fuma todos los días.}};
% Estar + gérondif : action en cours
\draw[poporange, very thick] (4.4,-0.7) -- (6.0,-0.7);
\draw[poporange, very thick, -{Stealth}] (6.0,-0.7) -- (6.4,-0.7);
\fill[poporange] (5.2,-0.7) circle (0.12);
\node[poporange, font=\small, align=center] at (8.4,-1.2) {estar + gerundio : acción en curso\\ \esp{Está fumando ahora.}};
\end{tikzpicture}
\legende{Le présent exprime des actions répétées (points bleus) ; \esp{estar + gerundio} désigne l'action en cours au moment où l'on parle (segment orange).}
\end{popfigure}

\begin{popfigure}
\begin{tikzpicture}
\node[draw=popgreen, fill=popgreenL, rounded corners, font=\bfseries, align=center] (estar) at (0,0) {ESTAR\\ (présent)};
\node[draw=popblue, fill=popblueL, rounded corners, align=center, font=\small] (c1) at (-4.2,-2) {estoy\\ estás\\ está};
\node[draw=popblue, fill=popblueL, rounded corners, align=center, font=\small] (c2) at (4.2,-2) {estamos\\ estáis\\ están};
\node[draw=poporange, fill=poporangeL, rounded corners, align=center, font=\small] (g1) at (-4.2,-4.2) {gérondifs réguliers\\ fum\textbf{ando} -- com\textbf{iendo}\\ beb\textbf{iendo} -- escrib\textbf{iendo}};
\node[draw=poppurple, fill=poppurpleL, rounded corners, align=center, font=\small] (g2) at (4.2,-4.2) {gérondifs irréguliers\\ d\textbf{u}rmiendo -- p\textbf{i}diendo\\ le\textbf{yendo} -- \textbf{yendo}};
\draw[-{Stealth}, thick, popdark] (estar) -- (c1);
\draw[-{Stealth}, thick, popdark] (estar) -- (c2);
\draw[-{Stealth}, thick, popdark] (c1) -- (g1);
\draw[-{Stealth}, thick, popdark] (c2) -- (g2);
\end{tikzpicture}
\legende{Schéma de construction : on conjugue \esp{estar} selon la personne, puis on ajoute le gérondif invariable.}
\end{popfigure}

\begin{center}
\begin{tabularx}{\linewidth}{|X|X|X|}
\hline
\rowcolor{popblueL} Français & Español & Remarque \\ \hline
Les jeunes fument dans la cour. (habitude) & Los jóvenes fuman en el patio. & présent simple \\ \hline
Un élève est en train d'insulter un autre. & Un alumno está insultando a otro. & estar + gérondif \\ \hline
Elle dort en classe en ce moment. & Está durmiendo en clase ahora. & gérondif irrégulier (dormir) \\ \hline
Nous sommes en train de lire l'article. & Estamos leyendo el artículo. & leer $\rightarrow$ leyendo \\ \hline
\end{tabularx}
\end{center}

\begin{attention}[Le piège n° 1 des francophones : jamais « ser » avec le gérondif]
En français, on dit « il \emph{est} en train de fumer », et le verbe « être » se traduit souvent par \esp{ser}. Mais avec le gérondif, c'est \textbf{toujours} \esp{estar} !\par
\esp{Es fumando.} est \textbf{FAUX}. On dit : \esp{Está fumando.}\par
Autre confusion à éviter : ne confondez pas le gérondif (\esp{fumando} = « en train de fumer ») et le participe passé (\esp{fumado} = « fumé »). \esp{Está fumado} signifierait « il est fumé » (comme un poisson fumé !), ce qui n'a rien à voir.
\end{attention}

\begin{savaistu}[Vocabulaire scolaire : Espagne vs Amérique latine / Guinée équatoriale]
En Espagne, le secondaire s'appelle \esp{instituto} (ou \esp{IES}). En Amérique latine et en Guinée équatoriale, on dit plutôt \esp{colegio}, \esp{liceo} ou \esp{secundaria}. Le terme \esp{aula de tutoría} (salle de tutorat / vie scolaire) existe partout, mais l'enseignant responsable est le \esp{tutor} (Espagne) ou \esp{profesor jefe / tutor} (Amérique). Pour le \esp{sexting}, l'anglicisme s'impose dans tout l'espace hispanophone, bien que l'Académie recommande \esp{sexting} (italique) ou la périphrase \esp{envío de contenido sexual}.
\end{savaistu}

\begin{exoresolu}[Habitude ou action en cours ?]
Énoncé : conjuguez le verbe entre parenthèses au présent simple ou avec \esp{estar + gerundio}, selon le sens.\par
1. Normalmente, los alumnos \_\_\_\_\_\_ (llegar) a clase a tiempo.\par
2. ¡Mira! Un chico \_\_\_\_\_\_ (fumar) detrás del gimnasio.\par
3. En este momento, la profesora \_\_\_\_\_\_ (explicar) la lección.\par
4. El acoso escolar \_\_\_\_\_\_ (destruir) la vida de muchos jóvenes.
\tcblower
Solution détaillée :\par
1. \esp{llegan} : l'adverbe \esp{normalmente} indique une habitude $\rightarrow$ présent simple, 3e personne du pluriel.\par
2. \esp{está fumando} : \esp{¡Mira!} attire l'attention sur une action en cours $\rightarrow$ \esp{estar} (3e pers. sing.) + gérondif régulier de \esp{fumar}.\par
3. \esp{está explicando} : \esp{en este momento} signale l'action en train de se dérouler $\rightarrow$ \esp{estar + explicando}.\par
4. \esp{destruye} : il s'agit d'un fait général, d'une vérité $\rightarrow$ présent simple, 3e personne du singulier.
\end{exoresolu}

\begin{exercice}[Entraînement : choisir la bonne forme]
Conjuguez le verbe entre parenthèses au présent de l'indicatif ou à \esp{estar + gerundio}.
\begin{enumerate}
\item Los profesores \_\_\_\_\_\_ (organizar) charlas de prevención cada mes.
\item ¡Silencio! El director \_\_\_\_\_\_ (hablar) por el megáfono ahora mismo.
\item Mi amigo \_\_\_\_\_\_ (no / fumar) nunca, dice que es veneno.
\item En este instante, los alumnos \_\_\_\_\_\_ (hacer) el examen de español.
\item La violencia \_\_\_\_\_\_ (causar) mucho dolor a las víctimas.
\end{enumerate}
\end{exercice}

\begin{corrige}[Exercice n°1]
\begin{enumerate}
\item \esp{organizan} (habitude, \emph{cada mes}) $\rightarrow$ présent simple.
\item \esp{está hablando} (action en cours, \emph{ahora mismo}) $\rightarrow$ \esp{estar + gerundio}.
\item \esp{no fuma} (habitude négative, \emph{nunca}) $\rightarrow$ présent simple.
\item \esp{están haciendo} (action en cours, \emph{en este instante}) $\rightarrow$ \esp{estar + gerundio} (irrégulier : \esp{hacer} $\rightarrow$ \esp{haciendo}).
\item \esp{causa} (fait général, vérité) $\rightarrow$ présent simple.
\end{enumerate}
\end{corrige}

\begin{aretenir}[L'essentiel de la section]
\begin{itemize}
\item \cle{Présent simple} = habitude, fait général, description durable : \esp{Los jóvenes fuman.} « Les jeunes fument. »
\item \cle{estar + gérondif} = action en cours au moment de la parole : \esp{Están fumando ahora.} « Ils sont en train de fumer maintenant. »
\item Gérondif : \textbf{-ando} (-ar), \textbf{-iendo} (-er/-ir) ; irréguliers principaux : \esp{diciendo, durmiendo, pidiendo, viniendo, teniendo, leyendo, oyendo, construyendo, yendo, pudiendo}.
\item Jamais \esp{ser} avec le gérondif : \esp{está fumando}, pas \esp{es fumando}.
\item Distinguer gérondif (\esp{-ando/-iendo}) et participe passé (\esp{-ado/-ido}) : \esp{está fumando} vs \esp{ha fumado}.
\end{itemize}
\end{aretenir}

Ces deux outils grammaticaux maîtrisés, nous sommes prêts pour la section suivante : la méthode du commentaire de texte, où il faudra repérer et analyser ces formes verbales dans un document authentique.

\coursec{Méthode du commentaire de texte}

Dans la section précédente, nous avons appris à \emph{décrire} une situation grâce au présent de l'indicatif et à la périphrase \esp{estar + gerundio} : \esp{Los jóvenes están consumiendo sustancias peligrosas.} Savoir décrire ne suffit pas : au baccalauréat comme à l'examen de fin d'année, on vous demandera de \emph{commenter} un texte, c'est-à-dire de le présenter, d'en dégager les idées, d'analyser les procédés de l'auteur et de donner votre avis. Cette section vous donne une méthode simple, en cinq étapes, appliquée au texte support \esp{« Jóvenes en riesgo »} étudié au début de la leçon.

\courssub{Qu'est-ce qu'un commentaire de texte ?}

\begin{definition}[Le commentaire de texte]
Le \cle{commentaire de texte} (\esp{comentario de texto}) est un exercice écrit qui consiste à :
\begin{itemize}
\item \textbf{présenter} le texte (sa nature, son thème, sa source supposée) ;
\item \textbf{expliquer} son contenu (idée principale et idées secondaires, organisées en plan) ;
\item \textbf{analyser} la manière dont l'auteur écrit (lexique, procédés, ton) ;
\item \textbf{donner son avis} personnel et argumenté.
\end{itemize}
En espagnol, on dit : \esp{presentar, resumir, analizar y opinar}.
\end{definition}

\begin{attention}[Commentaire n'est pas résumé]
Le piège classique du candidat francophone est de \emph{raconter} le texte avec ses propres mots. Un commentaire n'est pas un résumé : il faut \textbf{expliquer comment le texte est construit} et \textbf{justifier chaque idée par une citation} entre guillemets. Dire \esp{« El autor habla de las drogas »} ne suffit pas ; il faut écrire : \esp{« El autor denuncia el consumo de drogas cuando afirma: '...' »}.
\end{attention}

\courssub{Les cinq étapes du commentaire}

\textbf{Étape 1 : présenter le texte}\par
Avant toute analyse, situez le document en répondant à trois questions :
\begin{itemize}
\item \textbf{Nature} : de quel type de texte s'agit-il ? Un article de presse (\esp{artículo periodístico}), un récit (\esp{relato}), une lettre, un discours ?
\item \textbf{Thème} : de quoi parle le texte ? Ici : \esp{las conductas de riesgo de los jóvenes en el medio escolar y social}.
\item \textbf{Source supposée} : où ce texte a-t-il pu paraître ? Un journal, une revue scolaire, un site Internet de santé ?
\end{itemize}
Formule modèle : \esp{« El texto es un artículo periodístico que trata de las conductas de riesgo de los jóvenes. Probablemente fue publicado en un periódico o en una revista escolar. »} « Le texte est un article de presse qui traite des comportements à risque des jeunes. Il a probablement été publié dans un journal ou une revue scolaire. »

\textbf{Étape 2 : dégager l'idée principale en une phrase}\par
L'\cle{idée principale} (\esp{idea principal}) est le message essentiel de l'auteur, résumé en \textbf{une seule phrase}, sans détail. Technique : demandez-vous « que veut dire l'auteur, en somme ? ». Pour notre texte : \esp{« Los jóvenes adoptan hoy conductas de riesgo — drogas, alcohol, violencia, absentismo — y la escuela y la familia deben prevenirlas. »} « Les jeunes adoptent aujourd'hui des comportements à risque — drogues, alcool, violence, absentéisme — et l'école et la famille doivent les prévenir. »

\textbf{Étape 3 : construire un plan en deux ou trois parties}\par
Regroupez les \cle{idées secondaires} en deux ou trois grandes parties. Chaque partie correspond à un aspect du thème. Pour \esp{« Jóvenes en riesgo »}, un plan en deux parties s'impose naturellement :
\begin{enumerate}
\item \esp{Las sustancias peligrosas} : \esp{la droga, el alcohol, el tabaco, el sexting} (les substances dangereuses) ;
\item \esp{Los problemas de convivencia} : \esp{la violencia, el acoso escolar, el absentismo} (les problèmes de vie en communauté : violence, harcèlement, absentéisme).
\end{enumerate}
Annoncez ce plan à la fin de l'introduction : \esp{« Estudiaremos primero las sustancias peligrosas y después los problemas de convivencia en la escuela. »}

\textbf{Étape 4 : relever les procédés de l'auteur}\par
C'est le cœur de l'analyse. Observez \emph{comment} l'auteur écrit :
\begin{itemize}
\item \textbf{Le champ lexical du danger} : \esp{riesgo, peligro, droga, adicción, violencia, muerte, daño, sexting, acoso}. Tous ces mots appartiennent à la même famille de sens et créent une atmosphère inquiétante.
\item \textbf{L'énumération} : \esp{« la droga, el alcohol, el tabaco, el sexting, el acoso »} — la liste accumule les dangers et impressionne le lecteur.
\item \textbf{Le ton alarmiste} : l'auteur veut \emph{alerter} ; il utilise des verbes forts (\esp{denunciar, destruir, amenazar}), des adjectifs graves (\esp{peligroso, grave, preocupante}) et parfois des questions rhétoriques : \esp{« ¿Qué futuro tendrán nuestros jóvenes? »}
\end{itemize}
Pour chaque procédé, citez le texte entre guillemets et dites quel \textbf{effet} il produit sur le lecteur.

\textbf{Étape 5 : donner son avis argumenté}\par
Terminez par votre opinion personnelle, introduite par une formule d'opinion, suivie d'un \textbf{argument} (une raison) et si possible d'un \textbf{exemple} (votre école, votre quartier, le Cameroun).

\begin{exemplebox}[Formules d'opinion : exemples traduits]
\begin{itemize}
\item \esp{En mi opinión, la escuela debe prevenir estas conductas.} « À mon avis, l'école doit prévenir ces comportements. »
\item \esp{Pienso que la familia tiene un papel muy importante.} « Je pense que la famille a un rôle très important. »
\item \esp{Me parece que el acoso escolar es un problema grave en mi barrio.} « Il me semble que le harcèlement scolaire est un problème grave dans mon quartier. »
\item \esp{Estoy de acuerdo con el autor porque veo estos problemas en mi colegio.} « Je suis d'accord avec l'auteur parce que je vois ces problèmes dans mon collège. »
\item \esp{Sin embargo, creo que los jóvenes también son responsables de sus actos.} « Cependant, je crois que les jeunes sont aussi responsables de leurs actes. »
\end{itemize}
\end{exemplebox}

\begin{attention}[Pièges des francophones dans l'expression de l'opinion]
\begin{itemize}
\item Ne dites pas \esp{*en mi opinión que...} : \esp{en mi opinión} seul suffit, suivi d'une virgule.
\item Après \esp{pienso que / creo que / me parece que} (opinion affirmative), on utilise l'\textbf{indicatif} : \esp{Pienso que la droga \textbf{es} peligrosa.} Mais à la forme négative, l'espagnol exige le \textbf{subjonctif} : \esp{No pienso que la droga \textbf{sea} una solución.} « Je ne pense pas que la drogue soit une solution. »
\item \esp{Opinar} ne se traduit pas par « opiner » au sens de « donner son avis » dans une rédaction scolaire ; préférez \esp{dar su opinión} ou \esp{expresar su punto de vista}.
\end{itemize}
\end{attention}

\courssub{Tableau récapitulatif : les outils linguistiques du commentaire}

\begin{center}
\begin{tabularx}{\linewidth}{|X|X|X|}
\hline
\rowcolor{popblueL} Étape & Español (expressions utiles) & Français \\
\hline
Présenter & \esp{El texto es un artículo que trata de...} & Le texte est un article qui traite de... \\
\hline
Idée principale & \esp{La idea principal es que...} & L'idée principale est que... \\
\hline
Annoncer le plan & \esp{Estudiaremos primero... y después...} & Nous étudierons d'abord... puis... \\
\hline
Citer & \esp{El autor afirma: "..." / dice que...} & L'auteur affirme : « ... » / dit que... \\
\hline
Analyser & \esp{El autor utiliza un tono alarmista / una enumeración} & L'auteur utilise un ton alarmiste / une énumération \\
\hline
Opiner & \esp{En mi opinión... / Pienso que... / Me parece que...} & À mon avis... / Je pense que... / Il me semble que... \\
\hline
Conclure & \esp{En conclusión... / Para terminar...} & En conclusion... / Pour terminer... \\
\hline
\end{tabularx}
\end{center}

\courssub{Schéma : la structure du commentaire}

\begin{popfigure}
\begin{tikzpicture}[
  node distance=1cm and 1.5cm,
  etape/.style={rectangle, rounded corners, draw=popblue, fill=popblueL, text width=6cm, align=center, font=\small, inner sep=4pt},
  parte/.style={rectangle, rounded corners, draw=poporange, fill=poporangeL, text width=4.2cm, align=center, font=\small, inner sep=4pt},
  flecha/.style={-{Stealth[length=2.5mm]}, thick, popdark}
]
\node[etape, align=center] (intro){\textbf{Introducción}\\ naturaleza + tema + idea principal + plan};
\node[etape, below=of intro, align=center] (idea){\textbf{Idea principal}\\ una sola frase};
\node[parte, below left=of idea, align=center] (p1){\textbf{Parte 1}\\ las sustancias:\\ droga, alcohol, tabaco, sexting};
\node[parte, below right=of idea, align=center] (p2){\textbf{Parte 2}\\ la violencia,\\ el acoso, el absentismo};
\node[etape, below=3.5cm of idea, align=center] (op){\textbf{Opinión personal}\\ en mi opinión, pienso que...};
\node[etape, below=of op, align=center] (conc){\textbf{Conclusión}\\ balance + apertura};
\draw[flecha] (intro) -- (idea);
\draw[flecha] (idea) -- (p1);
\draw[flecha] (idea) -- (p2);
\draw[flecha] (p1) |- (op);
\draw[flecha] (p2) |- (op);
\draw[flecha] (op) -- (conc);
\end{tikzpicture}
\legende{Organigramme du commentaire de texte : de l'introduction à la conclusion.}
\end{popfigure}

\courssub{Méthode : rédiger l'introduction}

\begin{methode}[Rédiger l'introduction du commentaire en quatre phrases]
\begin{enumerate}
\item \textbf{Type de texte} : \esp{El texto es un artículo periodístico...} (un article, un récit, une lettre...).
\item \textbf{Thème} : \esp{...que trata de las conductas de riesgo de los jóvenes.}
\item \textbf{Idée principale} : \esp{El autor denuncia el consumo de drogas, alcohol, tabaco y sexting, así como la violencia y el absentismo escolar.}
\item \textbf{Annonce du plan} : \esp{Estudiaremos primero las sustancias peligrosas y después los problemas de convivencia, antes de dar nuestra opinión.}
\end{enumerate}
Quatre phrases suffisent : une introduction courte et claire vaut mieux qu'un long paragraphe confus.
\end{methode}

\begin{attention}[Au Bac : la citation obligatoire]
Au baccalauréat, le commentaire exige des \textbf{citations du texte entre guillemets} pour justifier chaque idée. Une idée non justifiée par le texte perd des points. Modèle de phrase : \esp{« El autor muestra la gravedad del problema cuando escribe: "muchos jóvenes están consumiendo drogas en los alrededores de la escuela". »} Recopiez la citation \emph{mot pour mot}, avec ses accents, et refermez toujours les guillemets.
\end{attention}

\courssub{Un commentaire modèle}

Lisez ce commentaire complet du texte \esp{« Jóvenes en riesgo »}. Repérez les quatre phrases de l'introduction, les deux parties du développement, les citations entre guillemets et l'opinion personnelle finale.

\begin{textebox}[Comentario modelo : « Jóvenes en riesgo »]
El texto es un artículo periodístico que trata de las conductas de riesgo de los jóvenes en el medio escolar y social. El autor denuncia el consumo de sustancias peligrosas y los problemas de convivencia. Estudiaremos primero las sustancias y después la violencia y el absentismo.

En la primera parte, el autor presenta las sustancias peligrosas. Utiliza una enumeración impresionante: "la droga, el alcohol, el tabaco y el sexting entran en la vida de los adolescentes". El campo léxico del peligro — "adicción", "daño", "riesgo" — crea un tono alarmista que quiere despertar a las familias.

En la segunda parte, el texto describe los problemas de convivencia: la violencia, el acoso escolar y el absentismo. El autor escribe: "muchos alumnos están faltando a clase y algunos sufren acoso en los pasillos". La pregunta final, "¿qué futuro tendrán nuestros jóvenes?", muestra su inquietud.

En mi opinión, la escuela debe prevenir estas conductas con campañas de información. Pienso que la familia también tiene un papel esencial, porque veo estos problemas en mi propio barrio. En conclusión, este texto es una llamada urgente a proteger a la juventud.
\end{textebox}

\textbf{Glossaire} : \textbf{convivencia} = vie en communauté ; \textbf{despertar} = réveiller, alerter ; \textbf{faltar a clase} = manquer les cours ; \textbf{pasillos} = couloirs ; \textbf{inquietud} = inquiétude ; \textbf{llamada} = appel ; \textbf{propio} = propre.

\courssub{Entraînement guidé}

\begin{exoresolu}[Identifier les cinq étapes dans le commentaire modèle]
Énoncé. Dans le commentaire modèle ci-dessus, retrouvez : 1) la phrase qui donne la nature du texte ; 2) la phrase qui annonce le plan ; 3) deux citations du texte support ; 4) deux formules d'opinion ; 5) la formule de conclusion.
\tcblower
Solution détaillée.
\begin{enumerate}
\item Nature du texte : \esp{« El texto es un artículo periodístico que trata de las conductas de riesgo de los jóvenes... »}
\item Annonce du plan : \esp{« Estudiaremos primero las sustancias y después la violencia y el absentismo. »}
\item Citations : \esp{« la droga, el alcohol, el tabaco y el sexting entran en la vida de los adolescentes »} et \esp{« muchos alumnos están faltando a clase y algunos sufren acoso en los pasillos »}.
\item Formules d'opinion : \esp{« En mi opinión, la escuela debe prevenir estas conductas »} et \esp{« Pienso que la familia también tiene un papel esencial »}.
\item Conclusion : \esp{« En conclusión, este texto es una llamada urgente a proteger a la juventud. »}
\end{enumerate}
\end{exoresolu}

\begin{exercice}[À vous de commenter]
Reprenez le texte \esp{« Jóvenes en riesgo »} de la section 1. Rédigez un commentaire de 10 à 12 lignes en suivant l'organigramme : introduction en quatre phrases, deux parties avec au moins une citation chacune, opinion personnelle avec \esp{en mi opinión} ou \esp{pienso que}, et conclusion commençant par \esp{en conclusión}. Soulignez vos citations et vos formules d'opinion.
\end{exercice}

\begin{aretenir}[L'essentiel de la méthode]
Un bon commentaire de texte suit cinq étapes : \textbf{présenter} (nature, thème, source), dégager l'\textbf{idée principale} en une phrase, construire un \textbf{plan} en deux ou trois parties, \textbf{analyser} les procédés (champ lexical du danger, énumération, ton alarmiste) avec des \textbf{citations entre guillemets}, et donner son \textbf{avis} avec \esp{en mi opinión}, \esp{pienso que}, \esp{me parece que}. L'introduction tient en quatre phrases : type de texte, thème, idée principale, annonce du plan.
\end{aretenir}

\coursec{Production guidée}

\courssub{Objectifs et consignes de la production}

Nous arrivons à l'aboutissement de cette leçon : après avoir lu, compris, analysé le lexique et révisé la grammaire nécessaire (\emph{présent de l'indicatif} et \emph{estar + gérondif}), vous allez maintenant \textbf{produire votre propre texte}. L'objectif est de mobiliser vos connaissances pour décrire une réalité observée dans votre environnement scolaire camerounais.

\begin{definition}[Consignes de l'exercice final]
Rédigez un court paragraphe descriptif (environ \textbf{5 à 8 lignes}) sur une conduite à risque observée dans votre lycée ou votre quartier.
\begin{itemize}
    \item \textbf{Lexique obligatoire :} utilisez \textbf{au moins 5 mots} de la liste : \cle{droga}, \cle{alcohol}, \cle{tabaco}, \cle{violencia}, \cle{acoso escolar}, \cle{sexting}, \cle{absentismo}, \cle{delincuencia}, \cle{adicción}.
    \item \textbf{Grammaire obligatoire :}
    \begin{enumerate}
        \item Au moins \textbf{2 verbes conjugués au présent de l'indicatif} (ex. \esp{consumen}, \esp{provoca}, \esp{observamos}).
        \item Au moins \textbf{1 périphrase \emph{estar + gérondif}} pour décrire une action en cours (ex. \esp{están fumando}, \esp{está sufriendo}).
    \end{enumerate}
    \item \textbf{Grille d'aide (à répondre dans le texte) :}
    \begin{enumerate}
        \item \esp{¿Qué conducta?} (Quelle conduite ?)
        \item \esp{¿Dónde?} (Où ?)
        \item \esp{¿Quién?} (Qui ? / Quels acteurs ?)
        \item \esp{¿Qué consecuencias?} (Quelles conséquences ?)
        \item \esp{¿Qué solución?} (Quelle solution ?)
    \end{enumerate}
    \item \textbf{Expression de l'opinion :} intégrez au moins une structure d'opinion : \esp{pienso que}, \esp{creo que}, \esp{es peligroso porque}, \esp{hay que + infinitif}.
\end{itemize}
\end{definition}

\begin{popfigure}
\centering
\begin{tikzpicture}[
    node distance=1.8cm and 1.2cm,
    question/.style={ellipse, draw=popblue, thick, fill=popblueL, text width=2.2cm, align=center, font=\small},
    final/.style={rectangle, rounded corners, draw=popgreen, thick, fill=popgreenL, text width=3.5cm, align=center, font=\small\bfseries},
    arrow/.style={-Stealth, popblue, thick}
]
% Central node
\node[final, align=center] (final){Paragraphe final\\ \emph{(5--8 lignes)}};

% Question nodes
\node[question] (q1) [above left=of final] {¿Qué conducta?};
\node[question] (q2) [above=of final] {¿Dónde?};
\node[question] (q3) [above right=of final] {¿Quién?};
\node[question] (q4) [below right=of final] {¿Qué consecuencias?};
\node[question] (q5) [below left=of final] {¿Qué solución?};

% Arrows
\draw[arrow] (q1) -- (final);
\draw[arrow] (q2) -- (final);
\draw[arrow] (q3) -- (final);
\draw[arrow] (q4) -- (final);
\draw[arrow] (q5) -- (final);
\end{tikzpicture}
\legende{Schéma de la grille d'aide : les cinq questions guident la construction du paragraphe descriptif.}
\end{popfigure}

\courssub{Outils linguistiques : Exprimer l'opinion et proposer des solutions}

Avant d'écrire, rappelons deux structures indispensables pour la partie « opinion et solution » de votre paragraphe. Elles sont très valorisées au Baccalauréat.

\begin{propriete}[La structure impersonnelle \texorpdfstring{\emph{hay que + infinitif}}{hay que + infinitif}]
\emph{Hay que} exprime une \textbf{obligation générale, impersonnelle} (équivalent à « il faut » en français). Il n'y a \textbf{pas de sujet} : la forme \emph{hay} (verbe \emph{haber} à la 3e personne du singulier, présent d'indicatif) reste invariée.
\begin{center}
\begin{tabularx}{\linewidth}{|X|X|X|}
\hline
\rowcolor{popblueL} \textbf{Espagnol} & \textbf{Français} & \textbf{Contexte} \\
\hline
\esp{Hay que prevenir el acoso escolar.} & Il faut prévenir le harcèlement scolaire. & Solution générale \\
\esp{Hay que ayudar a los jóvenes.} & Il faut aider les jeunes. & Action nécessaire \\
\esp{No hay que ignorar las señales.} & Il ne faut pas ignorer les signaux. & Interdiction / Conseil négatif \\
\hline
\end{tabularx}
\end{center}
\emph{Comparaison :} En français, « il faut » peut être suivi du subjonctif (\emph{il faut que tu fasses}) ou de l'infinitif (\emph{il faut faire}). En espagnol, \emph{hay que} \textbf{ne s'utilise qu'avec l'infinitif}. Pour exprimer « il faut que quelqu'un fasse quelque chose » avec un sujet précis, on utilise \emph{tener que + infinitif} (\esp{Tienes que estudiar}) ou \emph{deber + infinitif} (\esp{Debes estudiar}).
\end{propriete}

\begin{attention}[Piège majeur : \texorpdfstring{\emph{Hay que} vs \emph{Il faut que}}{Hay que vs Il faut que}]
Ne traduisez \textbf{jamais} « il faut que » par \emph{hay que} si vous introduisez un sujet différent.
\begin{itemize}
    \item \textbf{Faux :} \esp{Hay que los profesores vigilan.} (Incorrect : \emph{hay que} n'accepte pas de sujet).
    \item \textbf{Vrai :} \esp{Los profesores \textbf{tienen que} vigilar.} ou \esp{\textbf{Hay que} vigilar los recreos.}
    \item \textbf{Faux ami :} \emph{Hay} vient de \emph{haber} (il y a), pas de \emph{hacer} (faire). \esp{Hay que} = « Il y a [besoin] que » $\rightarrow$ « Il faut ».
\end{itemize}
\end{attention}

\begin{propriete}[Structures pour donner son opinion (niveau A2+/B1)]
Ces tournures permettent d'introduire votre point de vue personnel dans le paragraphe. Après \esp{pienso que / creo que / es \dots que}, on utilise l'\textbf{indicatif} (certitude/opinion).
\begin{center}
\begin{tabularx}{\linewidth}{|l|X|X|}
\hline
\rowcolor{poppurpleL} \textbf{Structure} & \textbf{Exemple espagnol} & \textbf{Traduction} \\
\hline
\esp{Pienso que} + indicatif & \esp{Pienso que el \textbf{sexting} \textbf{es} muy grave.} & Je pense que le sexting est très grave. \\
\esp{Creo que} + indicatif & \esp{Creo que \textbf{hay} mucha \textbf{violencia}.} & Je crois qu'il y a beaucoup de violence. \\
\esp{Es peligroso} + \esp{porque} + indicatif & \esp{Es peligroso \textbf{porque} \textbf{provoca} \textbf{adicción}.} & C'est dangereux car cela provoque une addiction. \\
\esp{Me parece} + \esp{que} + indicatif & \esp{Me parece \textbf{que es} un problema grave.} & Il me semble que c'est un problème grave. \\
\hline
\end{tabularx}
\end{center}
\emph{Note :} Le subjonctif n'apparaît qu'avec des expressions de doute, émotion ou négation (niveau B1/B2), pas nécessaire ici pour exprimer une opinion simple.
\end{propriete}

\courssub{Méthode : Rédiger un paragraphe descriptif cohérent}

\begin{methode}[Étapes pour rédiger votre paragraphe (5--8 lignes)]
\begin{enumerate}
    \item \textbf{Analyse de la grille :} Répondez mentalement aux 5 questions (\esp{¿Qué? ¿Dónde? ¿Quién? ¿Consecuencias? ¿Solución?}) en choisissant \textbf{une seule conduite précise} (ex. : la consommation de \emph{tabaco} et \emph{alcool} derrière les salles de classe pendant la récréation).
    \item \textbf{Phrase d'accroche (Introduction) :} Présentez le fait de manière générale. \textit{Ex. : \esp{En mi instituto de Yaoundé, observamos un problema creciente.}}
    \item \textbf{Description vivante (Corps du texte) :} Utilisez \textbf{\emph{estar + gérondif}} pour « filmer » la scène + \textbf{présent de l'indicatif} pour les faits habituels. Intégrez \textbf{3 à 4 mots de lexique} ici. \textit{Ex. : \esp{Algunos alumnos \textbf{están fumando} \textbf{tabaco} y \textbf{bebiendo alcohol} detrás de los bloques. Esto \textbf{provoca} \textbf{absentismo} y \textbf{violencia}.}}
    \item \textbf{Conséquences et Opinion :} Exprimez l'impact avec \esp{pienso que / es peligroso porque}. \textit{Ex. : \esp{\textbf{Pienso que} \textbf{es} muy \textbf{peligroso} \textbf{porque} \textbf{afecta} su salud y sus notas.}}
    \item \textbf{Solution (Conclusion) :} Utilisez \textbf{\emph{hay que + infinitif}}. \textit{Ex. : \esp{\textbf{Hay que} \textbf{organizar} charlas y \textbf{vigilar} los patios.}}
    \item \textbf{Relecture (Auto-correction) :} Vérifiez : 5 mots lexique ? 2 présents ? 1 \emph{estar+gerundio} ? Accords ? Accents ? Ponctuation ¿ ¡ ?
\end{enumerate}
\end{methode}

\begin{exemplebox}[Modèle de paragraphe réussi (niveau visé)]
\textbf{Texte espagnol :}
\esp{En mi \textbf{instituto} de Douala, hay un problema con el \textbf{alcohol} y el \textbf{tabaco}. Durante el recreo, un grupo de estudiantes \textbf{están fumando} y \textbf{bebiendo} cerveza cerca del muro. Ellos \textbf{creen} que es divertido, pero \textbf{pienso que} \textbf{es} muy \textbf{peligroso} \textbf{porque} \textbf{provoca} \textbf{adicción} y \textbf{absentismo}. Las notas \textbf{bajan} y hay más \textbf{violencia}. \textbf{Hay que} \textbf{informar} a los padres y \textbf{hay que} \textbf{controlar} las entradas.}

\textbf{Traduction :}
« Dans mon lycée de Douala, il y a un problème avec l'alcool et le tabac. Pendant la récréation, un groupe d'élèves \textbf{sont en train de fumer} et \textbf{boivent} de la bière près du mur. Ils \textbf{croient} que c'est amusant, mais \textbf{je pense que} \textbf{c'est} très \textbf{dangereux} \textbf{car} cela \textbf{provoque} de l'\textbf{addiction} et de l'\textbf{absentéisme}. Les notes \textbf{baissent} et il y a plus de \textbf{violence}. \textbf{Il faut} \textbf{informer} les parents et \textbf{il faut} \textbf{contrôler} les entrées. »

\textbf{Analyse du modèle :}
\begin{itemize}
    \item \textbf{Lexique (6) :} \emph{alcohol, tabaco, adicción, absentismo, violencia} (5+ atteints).
    \item \textbf{Présent indicatif (5+) :} \emph{hay, creen, pienso, es, provoca, bajan}.
    \item \textbf{Estar + gérondif (1) :} \emph{están fumando} (et \emph{están bebiendo} coordonné).
    \item \textbf{Opinion :} \emph{pienso que, es peligroso porque}.
    \item \textbf{Solution :} \emph{Hay que informar, hay que controlar}.
\end{itemize}
\end{exemplebox}

\courssub{Mise en œuvre : Activité d'écriture guidée}

\begin{experience}[Atelier d'écriture : « Mon lycée, ma réalité »]
\textbf{Durée :} 30 minutes. \textbf{Matériel :} Brouillon, grille d'aide, liste de lexique, grille de relecture (ci-dessous).

\textbf{Phase 1 : Remue-méninges individuel (5 min)}
Remplissez le tableau ci-dessous sur votre brouillon pour votre conduite choisie.

\begin{center}
\begin{tabularx}{\linewidth}{|X|c|}
\hline
\rowcolor{poporangeL} \textbf{Question guide} & \textbf{Mes notes / Mots-clés (en espagnol)} \\
\hline
\esp{¿Qué conducta?} & \textit{(ex: consumo de tabaco, sexting, acoso...)} \\
\hline
\esp{¿Dónde?} & \textit{(ex: detrás de las aulas, en el baño, en WhatsApp...)} \\
\hline
\esp{¿Quién?} & \textit{(ex: alumnos de 1ère, un grupo, las chicas...)} \\
\hline
\esp{¿Qué consecuencias?} & \textit{(ex: adicción, bajas notas, violencia, depresión...)} \\
\hline
\esp{¿Qué solución?} & \textit{(ex: hay que vigilar, hay que hablar, hay que castigar...)} \\
\hline
\end{tabularx}
\end{center}

\textbf{Phase 2 : Rédaction du brouillon (15 min)}
Rédigez votre paragraphe en suivant la \textbf{Méthode} (encadré ci-dessus). Ne cherchez pas la perfection du premier coup, écrivez des phrases simples et correctes.

\textbf{Phase 3 : Relecture par les pairs (10 min) — voir grille détaillée section suivante.}
Échangez votre brouillon avec un camarade. Il utilise la grille de correction. Vous corrigez ensuite votre version finale sur la copie à rendre.
\end{experience}

\begin{textebox}[Texte modèle d'élève (Simulation de production B1) — Pour analyse en classe]
\textbf{Contexto :} Un élève de 1ère A au Lycée Bilingue d'Essos (Yaoundé) décrit le \emph{sexting}.
\esp{En mi instituto, el \textbf{sexting} es una conducta de riesgo muy frecuente entre los chicos de mi edad. Muchos estudiantes \textbf{están enviando} fotos \textbf{íntimas} por WhatsApp sin pensar en las consecuencias. \textbf{Creo que} \textbf{es} un error grave \textbf{porque} la imagen \textbf{puede} circular muy \textbf{rápido} y causar \textbf{acoso escolar} y \textbf{depresión}. La \textbf{víctima} \textbf{siente} \textbf{vergüenza} y a veces \textbf{falta} a clase: hay mucho \textbf{absentismo}. Los profesores a veces no \textbf{saben} \textbf{qué} hacer. \textbf{Pienso que} \textbf{hay que} \textbf{educar} sobre el uso del \textbf{móvil} y \textbf{hay que} \textbf{denunciar} el \textbf{acoso} inmediatamente. No \textbf{hay que} callar.}

\textbf{Glossaire :}
\begin{itemize}
    \item \esp{enviando} (gérondif de \emph{enviar}) : en train d'envoyer.
    \item \esp{íntimas} (accord fém. pl. de \emph{íntimo}) : intimes.
    \item \esp{sin pensar en las consecuencias} : sans penser aux conséquences.
    \item \esp{circular} : circuler (se propager).
    \item \esp{rápido} (accent sur \emph{á}) : vite, rapidement.
    \item \esp{la víctima} (accent sur \emph{í}) : la victime.
    \item \esp{siente} (présent \emph{sentir}, diphtongue \emph{ie}) : ressent/éprouve.
    \item \esp{vergüenza} (tréma \emph{ü} pour \emph{gu} = \emph{güe}) : honte.
    \item \esp{falta a clase} (verbe \emph{faltar}, 3e pers. sg) : manque à la classe / est absent.
    \item \esp{saben} (présent \emph{saber}, irrégulier) : savent.
    \item \esp{qué} (accent interrogatif) : quoi / que.
    \item \esp{móvil} (accent sur \emph{ó}) : portable, mobile.
    \item \esp{denunciar} : signaler, porter plainte.
    \item \esp{callar} : se taire.
\end{itemize}

\textbf{Questions d'analyse (travail oral ou écrit) :}
\begin{enumerate}
    \item Combien de mots du lexique thématique trouvez-vous ? Citez-les.
    \item Identifiez la forme \emph{estar + gérondif}. Quelle action décrit-elle ?
    \item Relevez 3 verbes au présent de l'indicatif (autres que \emph{estar}).
    \item Quelles structures d'opinion et de solution sont utilisées ?
\end{enumerate}
\end{textebox}

\courssub{Relecture et correction par les pairs : Grille d'évaluation}

La relecture par les pairs (co-évaluation) est une compétence clé du programme. Elle développe votre sens critique et vous aide à repérer vos propres erreurs en les cherchant chez l'autre.

\begin{methode}[Grille de relecture par les pairs (Co-évaluation)]
Échangez votre texte. Cochez \textbf{OUI} ou \textbf{NON} et notez une correction si NON.
\begin{center}
\small
\begin{tabularx}{\linewidth}{|X|c|c|X|}
\hline
\rowcolor{poppinkL} \textbf{Critère} & \textbf{OUI} & \textbf{NON} & \textbf{Correction / Commentaire} \\
\hline
\textbf{1. LEXIQUE :} Au moins 5 mots de la liste (\emph{droga, alcohol, tabaco, violencia, acoso, sexting, absentismo, delincuencia, adicción}) ? & & & \\
\hline
\textbf{2. GRAMMAIRE :} Au moins 2 verbes au \textbf{présent de l'indicatif} (bien conjugués, accords sujet/verbe) ? & & & \\
\hline
\textbf{3. GRAMMAIRE :} Au moins 1 \textbf{\emph{estar + gérondif}} (bien formé : \emph{estar} conjugué + gérondif en \emph{-ando/-iendo}) ? & & & \\
\hline
\textbf{4. COHÉRENCE :} Les 5 questions de la grille (\esp{¿Qué? ¿Dónde? ¿Quién? ¿Consecuencias? ¿Solución?}) ont-elles une réponse dans le texte ? & & & \\
\hline
\textbf{5. OPINION / SOLUTION :} Présence de \esp{pienso que / creo que / es peligroso porque / hay que + infinitif} ? & & & \\
\hline
\textbf{6. ORTHOGRAPHE / ACCENTS :} Accents toniques corrects (\emph{adicción, están, jóvenes, solución, problema, rápido, móvil, qué, víctima, vergüenza}) ? Ponctuation \esp{¿ ? ¡ !} ? & & & \\
\hline
\textbf{7. LONGUEUR :} Entre 5 et 8 lignes (environ 60--100 mots) ? & & & \\
\hline
\end{tabularx}
\end{center}
\textbf{Consigne :} Si vous cochez « NON », écrivez la correction proposée dans la dernière colonne. Remettez la grille à l'auteur.
\end{methode}

\begin{exoresolu}[Exercice de relecture guidée : Corrigez ce paragraphe d'élève]
\textbf{Énoncé :} Voici un brouillon d'élève contenant des erreurs typiques (lexique, grammaire, cohérence). Appliquez la grille ci-dessus et proposez une version corrigée.

\esp{En mi escuela hay mucho problema con la droga. Los jovenes estan fumando marihuana en el baño. Ellos piensa que no pasa nada. Es peligroso por que hacen daño a la salud. Hay que hablar con los padres y los profesores debe vigilar mas. Tambien hay acoso y violentia.}

\tcblower

\textbf{Solution détaillée et corrections :}

\begin{enumerate}
    \item \textbf{Orthographe / Accents :} \esp{jovenes} $\rightarrow$ \esp{jóvenes} (accent sur \emph{ó} : oxytonne finissant par \emph{n/s/vowel}). \esp{estan} $\rightarrow$ \esp{están} (accent sur \emph{á} : même raison). \esp{por que} $\rightarrow$ \esp{porque} (conjonction = un mot, accent sur \emph{é}). \esp{violentia} $\rightarrow$ \esp{violencia} (\emph{c} devant \emph{i/e}, pas \emph{t}). \esp{Tambien} $\rightarrow$ \esp{También} (accent sur \emph{é}). \esp{mas} $\rightarrow$ \esp{más} (accent sur \emph{á} : monosyllabe tonique).
    \item \textbf{Grammaire (Accord) :} \esp{mucho problema} $\rightarrow$ \esp{muchos problemas} (accord pluriel). \esp{Ellos piensa} $\rightarrow$ \esp{Ellos piensan} (3e pers. pl. \emph{piensan}, diphtongue \emph{ie}). \esp{los profesores debe} $\rightarrow$ \esp{los profesores deben} (3e pers. pl.).
    \item \textbf{Grammaire (Gérondif) :} \esp{estan fumando} $\rightarrow$ \esp{están fumando} (corrigé accent). C'est le \emph{estar+gerundio} requis.
    \item \textbf{Lexique :} \esp{droga, acoso, violencia} (3 mots de la liste officielle ; \emph{marihuana} n'en fait pas partie). \textbf{Manque 2 mots.} \textit{Proposition : ajouter \emph{adicción} et \emph{absentismo}.}
    \item \textbf{Structure solution :} \esp{hay que hablar} (bien, impersonnel). \esp{profesores deben vigilar} (correct mais personnel). \textit{Mieux : garder \emph{hay que vigilar} pour l'impersonnel.}
    \item \textbf{Cohérence :} Répond aux 5 questions globalement.
\end{enumerate}

\textbf{Version corrigée (Modèle) :}
\esp{En mi escuela hay \textbf{muchos problemas} con la \textbf{droga}. Los \textbf{jóvenes} \textbf{están fumando} marihuana en el baño. Ellos \textbf{piensan} que no pasa nada. \textbf{Pienso que} \textbf{es peligroso} \textbf{porque} \textbf{provoca} \textbf{adicción} y \textbf{absentismo} y daña la salud. \textbf{Hay que} \textbf{hablar} con los padres y \textbf{hay que} \textbf{vigilar} más. \textbf{También} hay \textbf{acoso} y \textbf{violencia}.}
\end{exoresolu}

\courssub{Complément Bac : La production écrite à l'examen}

\begin{savaistu}[Le saviez-vous ? L'épreuve de production écrite au Baccalauréat (LV2 Espagnol)]
Au Cameroun, l'épreuve d'espagnol LV2 au Baccalauréat (série A) comporte une partie « Expression écrite » (souvent 6 à 8 points sur 20). Le correcteur utilise une grille analytique proche de notre grille de relecture. Les \textbf{trois piliers} de la notation sont :
\begin{enumerate}
    \item \textbf{Richesse lexicale (Vocabulaire) :} Utilisation précise du vocabulaire du module (\emph{conductas de riesgo, acoso, adicción, absentismo...}). Éviter le trop général (\emph{cosa, malo, bueno}). \textbf{Astuce :} Préparez des « îlots lexicaux » par thème (santé, école, famille, technologie).
    \item \textbf{Correction grammaticale (Morphosyntaxe) :} Maîtrise des temps de base (Présent, Passé composé/Indéfini, Imparfait, Futur, Conditionnel) + structures clés (\emph{estar+gerundio, hay que, tener que, ser/estar}). Les erreurs d'accord (genre/nombre) et d'accents pénalisent fortement.
    \item \textbf{Cohérence et Cohésion (Discours) :} Le texte doit avoir une structure logique (Introduction $\rightarrow$ Développement $\rightarrow$ Conclusion). Utilisez des \textbf{connecteurs} : \esp{por eso, además, sin embargo, por tanto, en conclusión}.
\end{enumerate}
\textbf{Conseil de pro :} Un texte court (80--100 mots) \textbf{bien structuré et sans fautes graves} vaut mieux qu'un long texte (150 mots) plein d'erreurs et incohérent. La \textbf{qualité prime sur la quantité}.
\end{savaistu}

\begin{aretenir}[Check-list finale avant de rendre votre copie (Bac / Devoir)]
\begin{itemize}
    \item \textbf{✓ Sujet respecté :} Une conduite à risque \textbf{camerounaise} (contexte local).
    \item \textbf{✓ Longueur :} 5--8 lignes (pas 3, pas 15).
    \item \textbf{✓ Contraintes linguistiques :} 5 mots lexique $\checkmark$ | 2 Présents $\checkmark$ | 1 \emph{Estar+Gerundio} $\checkmark$ | 1 Opinion/Solution $\checkmark$.
    \item \textbf{✓ Structure :} Intro (Contexte) -- Description (Action en cours + Faits) -- Conséquences -- Opinion -- Solution (\emph{Hay que}).
    \item \textbf{✓ Relecture :} Accents (\emph{adicción, jóvenes, están, solución, problema, rápido, móvil, qué, víctima, vergüenza}), \esp{¿ ? ¡ !}, accords (adj/nom, sujet/verbe), \emph{ser/estar} (description vs action en cours).
    \item \textbf{✓ Connecteurs :} \esp{y, pero, porque, por eso, también, en conclusión}.
\end{itemize}
\end{aretenir}

\begin{exercice}[Entraînement autonome : « Un autre regard »]
Choisissez une \textbf{autre} conduite à risque de la liste (différente de votre production principale) et rédigez un \textbf{deuxième paragraphe} (5 lignes) en respectant les mêmes contraintes (5 mots lexique, 2 présents, 1 \emph{estar+gerundio}, \emph{hay que}). Utilisez la grille d'aide.
\begin{itemize}
    \item \textbf{Sujet suggéré :} \emph{El absentismo escolar en mi barrio} / \emph{La violencia en los partidos de fútbol inter-clases} / \emph{El acoso escolar en las redes sociales}.
\end{itemize}
\textit{(Cet exercice peut être donné en devoir maison ou en contrôle continu).}
\end{exercice}

\begin{corrige}[Exercice « Un autre regard » — Pistes de correction type]
\textbf{Sujet :} \emph{El absentismo escolar en mi barrio.}
\esp{En mi barrio de Yaoundé, el \textbf{absentismo} escolar aumenta. Muchos chicos \textbf{están trabajando} en el mercado en vez de ir a clase. \textbf{Venden} cacao o \textbf{ayudan} a sus padres. \textbf{Creo que} \textbf{es} triste \textbf{porque} pierden su futuro y caen en la \textbf{delincuencia}. \textbf{Hay que} \textbf{apoyar} a las familias y \textbf{hay que} \textbf{controlar} la asistencia.}

\textbf{Vérification :}
\begin{itemize}
    \item Lexique (5) : \emph{absentismo, delincuencia, droga, violencia, adicción} (intégrés dans la version améliorée).
    \item Présents (3+) : \emph{venden, ayudan, pierden, caen, creo, es, aumenta}.
    \item \emph{Estar+gerundio} (1) : \emph{están trabajando}.
    \item Opinion/Solution : \emph{Creo que es triste porque / Hay que apoyar / Hay que controlar}.
\end{itemize}
\textbf{Version améliorée (intégrant 5 mots de la liste officielle) :}
\esp{En mi barrio de Yaoundé, el \textbf{absentismo} escolar aumenta. Muchos chicos \textbf{están trabajando} en el mercado en vez de ir a clase. \textbf{Venden} cacao o \textbf{ayudan} a sus padres. \textbf{Creo que} \textbf{es} triste \textbf{porque} pierden su futuro y caen en la \textbf{droga}, la \textbf{violencia} y la \textbf{adicción}. \textbf{Hay que} \textbf{apoyar} a las familias, \textbf{hay que} \textbf{prevenir} la \textbf{adicción} y \textbf{hay que} \textbf{controlar} la asistencia.}
\end{corrige}

\newpage

\coursec{Activité d'intégration}

\courssub{Objectifs de l'activité}
Cette activité d'intégration vise à mobiliser l'ensemble des savoirs et savoir-faire travaillés lors de la leçon : le lexique des conduites à risque (\cle{droga}, \cle{alcohol}, \cle{tabaco}, \cle{violencia}, \cle{acoso escolar}, \cle{sexting}, \cle{absentismo}, \cle{delincuencia}, \cle{adicción}), la conjugaison du présent de l'indicatif (réguliers et irréguliers fréquents) et la structure \cle{estar + gerundio} pour décrire des situations en cours. Les élèves devront produire un texte écrit (affiche de prévention) et une présentation orale courte, en situation de communication authentique, tout en développant des compétences transversales : travail collaboratif, esprit critique, créativité et prise de parole en public.

\courssub{Situation}
\begin{textebox}[Contexte : Campagne de prévention au Lycée Bilingue de Yaoundé]
El Departamento de Orientación del Liceo Bilingüe de Yaoundé lanza una campaña de sensibilización titulada « \emph{¡Tu vida vale más!} » para combatir las conductas de riesgo entre los estudiantes. El concurso está abierto a todos los alumnos de Première. Cada grupo debe diseñar un cartel (affiche) en español que será expuesto en los pasillos del establecimiento y presentado oralmente ante la clase. El mejor cartel, elegido por votación de los estudiantes y validado por el profesor, será impreso en gran formato y distribuido en los colegios asociados de la red UNESCO de Camerún.
\end{textebox}

\begin{definition}[Traduction et glossaire du contexte]
Le Département d'Orientation du Lycée Bilingue de Yaoundé lance une campagne de sensibilisation intitulée « Ta vie vaut plus ! » pour lutter contre les comportements à risque chez les élèves. Le concours est ouvert à tous les élèves de Première. Chaque groupe doit concevoir une affiche en espagnol qui sera exposée dans les couloirs de l'établissement et présentée oralement devant la classe. La meilleure affiche, choisie par vote des élèves et validée par le professeur, sera imprimée en grand format et distribuée dans les collèges partenaires du réseau UNESCO du Cameroun.
\par
\emph{Mots clés :} \esp{lanza} (lance), \esp{concurso} (concours), \esp{cartel} (affiche), \esp{pasillos} (couloirs), \esp{impreso} (imprimé), \esp{red} (réseau).
\end{definition}

\courssub{Consignes}
\begin{enumerate}
    \item \textbf{Formation des groupes (5 min) :} Constituez des groupes de 4 élèves. Désignez un \emph{coordinador} (chef de groupe), un \emph{secretario} (rédacteur), un \emph{diseñador} (illustrateur/mise en page) et un \emph{portavoz} (orateur principal).
    \item \textbf{Analyse du milieu (10 min) :} Sur une feuille de brouillon, dressez la liste de 3 à 4 conduites à risque que vous observez \emph{réellement} dans votre environnement scolaire (enceinte du lycée, abords, réseaux sociaux entre élèves) ou social (quartier, famille). Utilisez le lexique de la leçon. Exemple : \esp{El acoso escolar en los baños}, \esp{El consumo de alcohol en las fiestas de fin de semana}, \esp{El absentismo de los lunes}.
    \item \textbf{Rédaction du contenu textuel de l'affiche (25 min) :} Rédigez le texte de l'affiche en espagnol. Il doit comporter obligatoirement :
    \begin{itemize}
        \item Un \textbf{titre accrocheur} (impératif ou exclamation) : \esp{¡Di no a...!}, \esp{¡Detente!}, \esp{¡Piensa antes de actuar!}.
        \item Une \textbf{description de 3 situations observées} en utilisant \textbf{estar + gerundio} (ex. : \esp{Algunos estudiantes están fumando tabaco detrás del gimnasio.}, \esp{Muchos jóvenes están sufriendo acoso en las redes sociales.}).
        \item Les \textbf{conséquences} (présent indicatif) : \esp{La adicción destruye la salud.}, \esp{El absentismo provoca el fracaso escolar.}
        \item Des \textbf{solutions / appels à l'action} (impératif, \emph{hay que} + infinitif, \emph{deber} + infinitif) : \esp{Hay que hablar con el orientador.}, \esp{Denunciad el acoso.}, \esp{Debéis cuidar vuestra vida.}
    \end{itemize}
    \item \textbf{Conception visuelle (15 min) :} Le \emph{diseñador} réalise l'affiche sur papier A3 (fourni) ou sur tablette (Canva, PowerPoint). Intégrez des pictogrammes, des couleurs vives (rouge pour le danger, vert pour les solutions), le logo du lycée et le hashtag \esp{\#TuVidaValeMas}.
    \item \textbf{Répétition orale (10 min) :} Le \emph{portavoz} répète la présentation (2 minutes max). Les autres membres préparent des arguments pour répondre aux questions éventuelles de la classe.
    \item \textbf{Présentation et vote (20 min) :} Chaque groupe présente son affiche. La classe note sur une grille d'évaluation (lexique, grammaire, originalité, clarté orale). Vote final.
\end{enumerate}

\courssub{Ressources à mobiliser}
\begin{propriete}[Rappel : Présent de l'indicatif (usages descriptifs et habituels)]
On utilise le présent pour décrire des faits généraux, des habitudes ou des vérités actuelles.
\begin{center}
\begin{tabularx}{\linewidth}{|X|X|X|}
\hline
\rowcolor{popblueL} \textbf{Verbe} & \textbf{Conjugaison (exemple)} & \textbf{Exemple contextuel} \\
\hline
\esp{consumir} (régulier -ir) & \esp{consumo, consumes, consume, consumimos, consumís, consumen} & \esp{Los jóvenes consumen alcohol los fines de semana.} \\
\esp{fumar} (régulier -ar) & \esp{fumo, fumas, fuma, fumamos, fumáis, fuman} & \esp{Él fuma tabaco a escondidas.} \\
\esp{beber} (régulier -er) & \esp{bebo, bebes, bebe, bebemos, bebéis, beben} & \esp{Beber en exceso daña el hígado.} \\
\esp{tener} (irrégulier) & \esp{tengo, tienes, tiene, tenemos, tenéis, tienen} & \esp{Tienen problemas de adicción.} \\
\esp{poner} (irrégulier) & \esp{pongo, pones, pone, ponemos, ponéis, ponen} & \esp{Ponen en riesgo su futuro.} \\
\esp{hacer} (irrégulier) & \esp{hago, haces, hace, hacemos, hacéis, hacen} & \esp{Hacen sexting sin medir consecuencias.} \\
\hline
\end{tabularx}
\end{center}
\end{propriete}

\begin{propriete}[Rappel : Estar + Gerundio (actions en cours / situations temporaires)]
Structure : \textbf{estar (présent) + gérondif} (-\emph{ando} pour -ar, -\emph{iendo} pour -er/-ir).
\begin{center}
\begin{tabular}{|l|l|l|}
\hline
\rowcolor{popblueL} \textbf{Sujet} & \textbf{Estar} & \textbf{Gérondif (exemples)} \\
\hline
\esp{Yo} & \esp{estoy} & \esp{fumando, bebiendo, consumiendo, sufriendo} \\
\esp{Tú} & \esp{estás} & \esp{acosando, faltando, robando, vendiendo} \\
\esp{Él / Ella} & \esp{está} & \esp{viendo, leyendo, oyendo (\esp{leer/oir} $\to$ \esp{leyendo/oyendo})} \\
\esp{Nosotros} & \esp{estamos} & \esp{durmiendo (\esp{dormir} $\to$ \esp{durmiendo}), pidiendo (\esp{pedir} $\to$ \esp{pidiendo})} \\
\esp{Vosotros} & \esp{estáis} & \esp{sintiendo (\esp{sentir} $\to$ \esp{sintiendo})} \\
\esp{Ellos} & \esp{están} & \esp{muriendo (\esp{morir} $\to$ \esp{muriendo})} \\
\hline
\end{tabular}
\end{center}
\emph{Attention :} Changement vocalique pour les verbes en -\emph{ir} à radical changeant (\esp{e $\to$ i}, \esp{o $\to$ u}) au gérondif. Rappel : au Cameroun comme en Amérique latine et en Guinée équatoriale, on utilise \esp{ustedes} (3e pl.) à l'oral plutôt que \esp{vosotros} ; la forme \esp{están} couvre donc les deux contextes.
\end{propriete}

\begin{aretenir}[Lexique clé à réemployer impérativement]
\begin{center}
\begin{tabularx}{\linewidth}{|X|X|}
\hline
\rowcolor{popblueL} \textbf{Noms / Concepts} & \textbf{Verbes / Expressions d'action} \\
\hline
\esp{la droga}, \esp{el alcohol}, \esp{el tabaco}, \esp{la violencia}, \esp{el acoso escolar}, \esp{el ciberacoso}, \esp{el sexting}, \esp{el absentismo}, \esp{la delincuencia}, \esp{la adicción}, \esp{la salud mental}, \esp{la presión de grupo} & \esp{consumir}, \esp{fumar}, \esp{beber}, \esp{sufrir}, \esp{padecer}, \esp{faltar (a clase)}, \esp{agredir}, \esp{amenazar}, \esp{chantajear}, \esp{denunciar}, \esp{pedir ayuda}, \esp{hablar con}, \esp{acudir a}, \esp{evitar}, \esp{prevenir} \\
\hline
\end{tabularx}
\end{center}
\end{aretenir}

\begin{methode}[Comment présenter son affiche oralement (2 minutes)]
\begin{enumerate}
    \item \textbf{Accroche (15 s) :} Donnez le titre et le slogan. \esp{« Buenas tardes. Nuestro cartel se titula "¡Rompe el silencio!" y denuncia el acoso escolar. »}
    \item \textbf{Description des faits (45 s) :} Citez les 3 situations avec \esp{estar + gerundio}. \esp{« En nuestro liceo, vemos que algunos alumnos están sufriendo acoso en los baños... Otros están consumiendo alcohol antes de entrar a clase... »}
    \item \textbf{Conséquences (30 s) :} Présent indicatif. \esp{« Esto provoca baja autoestima, absentismo y fracaso escolar. »}
    \item \textbf{Solutions et appel (30 s) :} \esp{Hay que / Debéis / Imperativo}. \esp{« Hay que denunciar. ¡Hablad con vuestro tutor! »}
    \item \textbf{Conclusion (15 s) :} \esp{« Gracias. Votad por nuestra campaña para un liceo sin violencia. »}
\end{enumerate}
\end{methode}

\courssub{Proposition de production et corrigé commenté}

\begin{textebox}[Affiche modèle : « ¡Rompe el silencio! » — Groupe 3]
¡ROMPE EL SILENCIO! \\
¡EL ACOSO NO ES UN JUEGO! \\[0.5em]
\textbf{¿Qué está pasando en nuestro liceo?} \\
\textbullet\ \esp{Algunos estudiantes están sufriendo acoso escolar en los pasillos y baños.} \\
\textbullet\ \esp{Varios compañeros están recibiendo amenazas por WhatsApp (ciberacoso).} \\
\textbullet\ \esp{Muchas víctimas están callando por miedo a las represalias.} \\[0.5em]
\textbf{Consecuencias graves :} \\
\esp{El acoso destruye la autoestima.} \\
\esp{El absentismo aumenta cada día.} \\
\esp{La depresión y la ansiedad afectan el rendimiento académico.} \\[0.5em]
\textbf{¡ACTÚA YA!} \\
\esp{Hay que denunciar: habla con tu orientador o profesor de confianza.} \\
\esp{No seas cómplice: si ves algo, di algo.} \\
\esp{Llama al 116 (Línea de ayuda al niño en Camerún) o acude a la Dirección.} \\
\esp{\#TuVidaValeMas \quad \#LiceoSeguro}
\end{textebox}

\begin{definition}[Traduction de l'affiche modèle]
BRISER LE SILENCE ! LE HARCÈLEMENT N'EST PAS UN JEU ! / Que se passe-t-il dans notre lycée ? / Certains élèves subissent du harcèlement scolaire dans les couloirs et les toilettes. / Plusieurs camarades reçoivent des menaces sur WhatsApp (cyberharcèlement). / Beaucoup de victimes se taisent par peur des représailles. / Conséquences graves : Le harcèlement détruit l'estime de soi. L'absentéisme augmente chaque jour. La dépression et l'anxiété affectent les résultats scolaires. / AGIS MAINTENANT ! / Il faut dénoncer : parle à ton conseiller d'orientation ou à un prof de confiance. / Ne sois pas complice : si tu vois quelque chose, dis-le. / Appelle le 116 (Ligne d'aide à l'enfant au Cameroun) ou va à la Direction. / \#TaVieVautPlus \#LyceeSur.
\end{definition}

\begin{exoresolu}[Analyse linguistique de l'affiche modèle]
\textbf{Énoncé :} À partir de l'affiche ci-dessus, justifiez les choix grammaticaux et lexicaux effectués par le groupe.

\tcblower
\textbf{Commentaire des choix de langue (Corrigé commenté) :}
\begin{enumerate}
    \item \textbf{Titre et accroche :} \esp{¡Rompe el silencio!} (Impératif affirmatif 2e pers. sg. de \esp{romper}. Le COD \esp{el silencio} est explicite, pas de pronom enclitique. \esp{Rómpelo} signifierait « Brise-le »). \esp{¡El acoso no es un juego!} (Présent indicatif \esp{es} de \esp{ser} pour définition/identité essentielle). L'usage de l'impératif et de l'exclamation capte l'attention immédiatement.
    \item \textbf{Description (\esp{estar + gerundio}) :} Les trois puces illustrent parfaitement la structure cible pour décrire des situations \emph{en cours, observées actuellement} dans l'enceinte scolaire.
    \begin{itemize}
        \item \esp{están sufriendo} (3e pl. \esp{estar} + gérondif \esp{sufrir} $\to$ \esp{sufriendo}) : action subie, continue.
        \item \esp{están recibiendo} (\esp{recibir} $\to$ \esp{recibiendo}) : réception continue de messages.
        \item \esp{están callando} (\esp{callar} $\to$ \esp{callando}) : attitude passive maintenue dans le temps.
    \end{itemize}
    \textbf{Point de vigilance :} Ne pas confondre avec le présent simple (\esp{sufren}, \esp{reciben}, \esp{callan}) qui exprimerait une habitude ou une vérité générale. Ici, le gérondif ancre l'action dans le « maintenant » de l'observation, ce qui renforce l'urgence.
    \item \textbf{Conséquences (Présent indicatif) :} \esp{destruye}, \esp{aumenta}, \esp{afectan}. Verbes réguliers (-ar, -er) et irréguliers (\esp{afectan} de \esp{afectar}). Le présent exprime des vérités factuelles, des causalités établies.
    \item \textbf{Solutions (Modalité obligation/conseil) :}
    \begin{itemize}
        \item \esp{Hay que denunciar} : Impersonnel, valeur de nécessité générale (très courant en affichage public).
        \item \esp{habla} (Impératif \esp{tú} de \esp{hablar} -- attention : \esp{habla}, pas \esp{hable} qui est subjonctif/\emph{usted}), \esp{di} (Impératif irrégulier \esp{tú} de \esp{decir} -- \textbf{piège fréquent : *dice}).
        \item \esp{No seas cómplice} (Impératif négatif \esp{tú} de \esp{ser} $\to$ \esp{no seas} + subjonctif présent). \textbf{Attention :} L'impératif négatif emprunte toujours la forme du subjonctif (\esp{no seas}, \esp{no hagas}, \esp{no digas}).
        \item \esp{Llama}, \esp{acude} : Impératifs réguliers positifs \esp{tú}.
    \end{itemize}
    \item \textbf{Cohésion et connecteurs :} \esp{por miedo a} (locution prépositive), \esp{si ves algo, di algo} (structure conditionnelle réelle \esp{si} + présent + impératif/présent), \esp{para un liceo sin violencia} (finalité).
    \item \textbf{Ancrage contextuel camerounais :} L'intégration du numéro \esp{116} (numéro vert réel pour la protection de l'enfance au Cameroun, géré par le MINAS) et la mention \esp{Dirección} (Proviseur/Administration) ancrent l'affiche dans la réalité vécue par les élèves de Yaoundé/Douala.
\end{enumerate}
\end{exoresolu}

\courssub{Bilan de l'activité}
Cette activité de synthèse a permis de vérifier l'appropriation des trois piliers de la leçon :
\begin{enumerate}
    \item \textbf{Lexique :} Les élèves ont dû sélectionner et orthographier correctement les termes techniques (\esp{acoso}, \esp{ciberacoso}, \esp{absentismo}, \esp{adicción}) dans un texte court et percutant. L'affichage public impose une rigueur orthographique (accents, \esp{ñ}, ponctuation \esp{¡! ¿?}) supérieure au brouillon.
    \item \textbf{Grammaire :} L'opposition \textbf{Présent indicatif} (faits, conséquences, vérités) / \textbf{Estar + Gérondif} (observation immédiate, processus en cours) a été mise en œuvre de manière fonctionnelle : on \emph{décrit} ce qu'on \emph{voit} (\esp{están sufriendo}) pour \emph{expliquer} ce que cela \emph{provoque} (\esp{destruye, aumenta}). L'usage des impératifs (positifs et négatifs) et de \esp{hay que} a été contextualisé dans la persuasion et le conseil.
    \item \textbf{Compétences langagières :} La compréhension de la consigne (écrite en français, ressources en espagnol), la production écrite collaborative (négociation du sens entre pairs), la production orale (prise de parole en continu, 2 min, sans notes idéalement) et l'interaction (réponses aux questions, vote argumenté) ont été évaluées.
    \item \textbf{Dimension citoyenne :} L'ancrage dans le réel (numéro 116, lieux précis du lycée, pression de groupe) transforme l'exercice scolaire en acte de prévention réel. Le vote par les pairs valorise la clarté du message et l'impact visuel, critères essentiels de la communication en santé publique.
\end{enumerate}

\begin{attention}[Erreurs fréquentes observées lors de cette activité (à corriger en grand groupe)]
\begin{itemize}
    \item \textbf{Confusion \esp{ser} / \esp{estar} :} \esp{*Los estudiantes son fumando} (incorrect) $\to$ \esp{Los estudiantes \textbf{están} fumando}. Rappel : \esp{estar} pour action en cours / localisation / état passager.
    \item \textbf{Gérondif mal formé (changement vocalique -ir) :} \esp{*están leendo} (correct : \esp{leyendo}), \esp{*están oiendo} (correct : \esp{oyendo}), \esp{*están durmiendo} (correct : \esp{durmiendo}), \esp{*están pidiendo} (correct : \esp{pidiendo}), \esp{*están sintiendo} (correct : \esp{sintiendo}), \esp{*están muriendo} (correct : \esp{muriendo}). \textbf{Règle :} Verbes en -\emph{ir} radicaux changeants (\esp{e>i}, \esp{o>u}) au gérondif ; \emph{i} entre voyelles $\to$ \emph{y} (\esp{leer/oir/construir}).
    \item \textbf{Impératif négatif :} \esp{*No eres cómplice} (indicatif) $\to$ \esp{\textbf{No seas} cómplice} (subjonctif/impératif négatif). \esp{*No haces eso} $\to$ \esp{\textbf{No hagas} eso}. \esp{*No digas eso} (correct).
    \item \textbf{Faux ami :} \esp{*La adicción es un problema de \textbf{realización}} (pour « réalité » ou « fait réel »). \esp{Realización} = accomplissement, exécution. Dire \esp{un problema de \textbf{realidad}} ou \esp{un \textbf{hecho} real}.
    \item \textbf{Accents manquants / surnuméraires :} \esp{*esta} (démonstratif/pronom) vs \esp{está} (verbe), \esp{*pasa} (présent) vs \esp{pasá} (impératif \emph{vosotros} non utilisé ici), \esp{*habla} (ok, impératif \emph{tú}) vs \esp{*hablá} (\emph{vosotros}).
\end{itemize}
\end{attention}

\begin{savaistu}[Le 116 au Cameroun : Allô Enfance en Danger]
Le numéro \textbf{116} est le numéro vert gratuit, confidentiel et accessible 24h/24 au Cameroun pour signaler toute situation de violence, d'exploitation, de négligence ou de danger concernant un enfant (violences scolaires, mariages précoces, travail des enfants, violences sexuelles). Il est géré par le \emph{Ministère des Affaires Sociales (MINAS)} en partenariat avec l'UNICEF. L'intégrer dans une affiche de prévention en classe d'espagnol, c'est faire du cours de langue un vecteur de protection réelle pour les élèves. En Espagne, le numéro équivalent pour l'aide à l'enfance est le \textbf{116 111} (harmonisé en Europe). En Guinée équatoriale, pays hispanophone voisin, des mécanismes similaires existent sous l'égide du \emph{MINASFE} (Ministère des Affaires Sociales et de la Promotion de la Femme).
\end{savaistu}

\newpage

\coursec{Méthodes et exercices résolus}

\courssub{Texte support : « Jóvenes en riesgo »}

\begin{textebox}[Jóvenes en riesgo — article de presse scolaire]
En muchos barrios de Yaoundé y de Douala, algunos jóvenes empiezan a fumar y a beber alcohol muy pronto, a veces a los trece o catorce años. Cerca de los colegios, unos vendedores ofrecen cigarrillos y bebidas baratas a los alumnos. En el patio, el acoso escolar es otro problema grave : un grupo de estudiantes insulta o golpea a un compañero más débil, y la víctima tiene miedo de ir a clase.

Además, con los teléfonos móviles y las redes sociales, aparece un nuevo peligro : el sexting. Algunos adolescentes envían fotos íntimas y después sufren burlas y vergüenza. El absentismo también preocupa a los profesores : cada día, varios alumnos faltan a clase y caminan por el mercado o juegan al fútbol en la calle.

Ahora mismo, los padres, los maestros y las autoridades están buscando soluciones. En muchos colegios, los profesores están organizando charlas para informar a los jóvenes sobre los peligros de la droga, del alcohol y del tabaco. La educación es la mejor arma para proteger a los adolescentes y construir un futuro sin violencia ni adicciones.
\end{textebox}

\textbf{Traduction} : « Dans beaucoup de quartiers de Yaoundé et de Douala, certains jeunes commencent à fumer et à boire de l'alcool très tôt, parfois à treize ou quatorze ans. Près des collèges, des vendeurs proposent des cigarettes et des boissons bon marché aux élèves. Dans la cour, le harcèlement scolaire est un autre problème grave : un groupe d'élèves insulte ou frappe un camarade plus faible, et la victime a peur d'aller en classe.

De plus, avec les téléphones portables et les réseaux sociaux, un nouveau danger apparaît : le sexting. Certains adolescents envoient des photos intimes et subissent ensuite moqueries et honte. L'absentéisme inquiète aussi les professeurs : chaque jour, plusieurs élèves manquent les cours et se promènent au marché ou jouent au football dans la rue.

En ce moment même, les parents, les enseignants et les autorités sont en train de chercher des solutions. Dans beaucoup de collèges, les professeurs organisent des causeries pour informer les jeunes sur les dangers de la drogue, de l'alcool et du tabac. L'éducation est la meilleure arme pour protéger les adolescents et construire un avenir sans violence ni addictions. »

\textbf{Glossaire} : \esp{barato} = bon marché ; \esp{más débil} = plus faible ; \esp{burlas} = moqueries ; \esp{charlas} = causeries, exposés de sensibilisation ; \esp{arma} = arme (féminin : \esp{el arma}) ; \esp{faltar a clase} = manquer les cours.

\courssub{Lire et comprendre un texte sur les conduites à risque}

\begin{methode}[Comprendre un texte : la lecture méthodique]
Pour réussir les questions de compréhension à l'examen, suis toujours la même démarche :
\begin{enumerate}
\item \textbf{Première lecture globale} : lis le texte une fois sans dictionnaire. Repère le thème général (ici : les comportements à risque des jeunes), le type de texte (article, récit, reportage) et le ton (informatif, alarmiste).
\item \textbf{Repérage des mots-clés} : souligne le lexique du champ : \esp{droga, alcohol, tabaco, violencia, acoso escolar, sexting, absentismo, adicción}. Ces mots te donnent le plan du texte.
\item \textbf{Lecture des questions} : lis les questions AVANT la deuxième lecture. Elles t'indiquent quels passages regarder.
\item \textbf{Deuxième lecture ciblée} : pour chaque question, localise la phrase du texte qui contient la réponse. Numérote-la mentalement.
\item \textbf{Rédaction de la réponse} : réponds par une phrase complète en espagnol, en reprenant les mots de la question, jamais par un simple « sí » ou « no ». Si on demande de justifier, cite le texte entre guillemets.
\item \textbf{Vérification} : relis ta réponse : verbe conjugué à la bonne personne, accords, accents.
\end{enumerate}
\textbf{Réflexe} : une réponse de compréhension = idée du texte + reformulation personnelle. Ne recopie jamais tout un paragraphe.
\end{methode}

\begin{exoresolu}[Compréhension de texte : « Jóvenes en riesgo »]
\textbf{Questions} sur le texte support : 1) ¿Qué comportamientos de riesgo aparecen en los barrios de Yaoundé y Douala? 2) ¿Qué problemas preocupan a los profesores? 3) ¿Qué están haciendo los padres, los maestros y las autoridades?
\tcblower
\textbf{Raisonnement} : la question 1 porte sur le milieu social (\esp{barrios}) : on localise le premier paragraphe. La question 2 porte sur le milieu scolaire : deuxième paragraphe. La question 3 exige de reconnaître la structure \esp{estar + gerundio} dans le dernier paragraphe : elle décrit une action en cours.

\textbf{Réponses rédigées} :
\begin{enumerate}
\item \esp{En los barrios, algunos jóvenes fuman y beben alcohol muy pronto.} « Dans les quartiers, certains jeunes fument et boivent de l'alcool très tôt. »
\item \esp{A los profesores les preocupan el acoso escolar, el sexting y el absentismo.} « Le harcèlement scolaire, le sexting et l'absentéisme inquiètent les professeurs. » Attention : \esp{preocupan} est au pluriel car le sujet est une énumération ; avec le verbe \esp{preocupar}, la personne qui ressent l'inquiétude est introduite par \esp{a} : \esp{a los profesores les preocupan...}
\item \esp{Están buscando soluciones y están organizando charlas para informar a los jóvenes.} « Ils sont en train de chercher des solutions et d'organiser des causeries pour informer les jeunes. » On justifie : \esp{están buscando} = \esp{estar} (3\textsuperscript{e} pers. plur.) + gérondif de \esp{buscar}, action en cours de réalisation.
\end{enumerate}
\textbf{Conclusion} : chaque réponse reprend les termes de la question, cite l'information exacte du texte et respecte la grammaire (accords, conjugaison).
\end{exoresolu}

\courssub{Lexique thématique par champs}

\begin{definition}[Les conduites à risque : le vocabulaire essentiel]
\begin{center}
\begin{tabularx}{\linewidth}{|X|X|X|}\hline
\rowcolor{popblueL} Español & Français & Exemple \\ \hline
\esp{la droga / las drogas} & la drogue & \esp{La droga destruye vidas.} \\ \hline
\esp{el alcohol} & l'alcool & \esp{El alcohol causa enfermedades.} \\ \hline
\esp{el tabaco / el tabaquismo} & le tabac / le tabagisme & \esp{El tabaco daña los pulmones.} \\ \hline
\esp{la violencia} & la violence & \esp{La violencia no es una solución.} \\ \hline
\esp{el acoso escolar} & le harcèlement scolaire & \esp{El acoso escolar asusta a las víctimas.} \\ \hline
\esp{el sexting} & le sexting & \esp{El sexting es un peligro nuevo.} \\ \hline
\esp{el absentismo} & l'absentéisme & \esp{El absentismo provoca el fracaso escolar.} \\ \hline
\esp{la delincuencia} & la délinquance & \esp{La delincuencia juvenil aumenta.} \\ \hline
\esp{la adicción} & l'addiction & \esp{La adicción al móvil es frecuente.} \\ \hline
\end{tabularx}
\end{center}
\textbf{Verbes associés} : \esp{fumar} (fumer), \esp{beber} (boire), \esp{consumir drogas} (consommer de la drogue), \esp{acosar} (harceler), \esp{insultar} (insulter), \esp{golpear} (frapper), \esp{faltar a clase} (manquer les cours), \esp{enviar fotos} (envoyer des photos), \esp{proteger} (protéger), \esp{prevenir} (prévenir), \esp{luchar contra} (lutter contre).
\end{definition}

\begin{attention}[Genre et accents : les pièges du lexique]
\begin{itemize}
\item \esp{la droga} est FÉMININ ; mais \esp{el alcohol}, \esp{el tabaco}, \esp{el acoso}, \esp{el absentismo} sont masculins.
\item \esp{la adicción, la violencia, la atención, la solución, la educación} : les mots en \esp{-ción} sont féminins et portent TOUJOURS un accent sur la ó.
\item « drogué / toxicomane » se dit \esp{drogadicto, -a} ; « un jeune drogué » = \esp{un joven drogadicto}.
\item \esp{el miedo} est masculin : « la peur » = \esp{el miedo}. Ne dis jamais « la miedo ».
\end{itemize}
\end{attention}

\begin{savaistu}[Le Cameroun, la Guinée équatoriale et la sensibilisation des jeunes]
La Guinée équatoriale, voisine du Cameroun, est le seul pays d'Afrique subsaharienne dont la langue officielle est l'espagnol : on y parle de \esp{los jóvenes}, \esp{la escuela} et \esp{la salud} comme à Madrid ! Au Cameroun comme dans toute l'Afrique hispanophone, les campagnes de sensibilisation contre \esp{la droga}, \esp{el alcohol} et \esp{el acoso escolar} utilisent des slogans simples au présent : \esp{Di no a las drogas} « Dis non à la drogue », \esp{La educación es el camino} « L'éducation est le chemin ». Retiens ces slogans : ils sont d'excellents modèles pour tes productions écrites.
\end{savaistu}

\courssub{Nommer et décrire les comportements à risque}

\begin{methode}[Décrire un comportement à risque en quatre points]
Pour décrire une conduite à risque à l'écrit comme à l'oral :
\begin{enumerate}
\item \textbf{Nommer} la conduite : \esp{El acoso escolar es...}, \esp{El consumo de drogas es...}
\item \textbf{Définir} avec une relative : \esp{...un comportamiento que...} / \esp{...una situación en la que...}
\item \textbf{Situer} le lieu : \esp{en el colegio} (milieu scolaire), \esp{en el barrio, en la calle, en casa} (milieu social), \esp{en las redes sociales} (milieu numérique).
\item \textbf{Montrer la conséquence} avec \esp{por eso, como resultado, y entonces} : \esp{...y por eso los alumnos tienen miedo.}
\end{enumerate}
\textbf{Structures à retenir} :
\begin{itemize}
\item \esp{consistir en + infinitivo} : \esp{El absentismo consiste en faltar a clase.} « L'absentéisme consiste à manquer les cours. »
\item \esp{provocar / causar + nom} : \esp{El alcohol causa enfermedades.} « L'alcool provoque des maladies. »
\item \esp{estar + gerundio} pour une scène : \esp{Un grupo de jóvenes está fumando detrás del colegio.} « Un groupe de jeunes est en train de fumer derrière le collège. »
\end{itemize}
\end{methode}

\begin{exoresolu}[Définitions : reconnaître la conduite]
\textbf{Énoncé} : Voici quatre définitions. Trouve la conduite correspondante et complète par une phrase personnelle qui montre une conséquence. Conduites : \esp{el sexting, el absentismo, el acoso escolar, el tabaquismo}.
\begin{enumerate}
\item \esp{Es un comportamiento que consiste en insultar o golpear a un compañero en el colegio.}
\item \esp{Es la costumbre de fumar cigarrillos todos los días.}
\item \esp{Consiste en enviar fotos íntimas por el teléfono móvil.}
\item \esp{Es la situación de un alumno que falta a clase con frecuencia.}
\end{enumerate}
\tcblower
\textbf{Raisonnement} : on repère dans chaque définition le mot-clé : \esp{insultar o golpear} (violence entre élèves), \esp{fumar cigarrillos} (tabac), \esp{fotos íntimas} (images), \esp{faltar a clase} (absences).

\textbf{Solution rédigée} :
\begin{enumerate}
\item \esp{Es el acoso escolar. Como resultado, la víctima tiene miedo y no quiere ir al colegio.} « C'est le harcèlement scolaire. En conséquence, la victime a peur et ne veut plus aller au collège. »
\item \esp{Es el tabaquismo. Por eso, muchos jóvenes tienen problemas de salud.} « C'est le tabagisme. C'est pourquoi beaucoup de jeunes ont des problèmes de santé. »
\item \esp{Es el sexting. Esta conducta provoca vergüenza y problemas psicológicos.} « C'est le sexting. Cette conduite provoque honte et problèmes psychologiques. »
\item \esp{Es el absentismo. El alumno ausente no aprende y puede repetir el año.} « C'est l'absentéisme. L'élève absent n'apprend pas et peut redoubler. »
\end{enumerate}
\textbf{Conclusion} : chaque réponse nomme la conduite, puis ajoute une conséquence introduite par un connecteur logique (\esp{como resultado, por eso, y entonces}). C'est exactement ce qu'attend le correcteur.
\end{exoresolu}

\courssub{Grammaire : présent de l'indicatif et estar + gérondif}

\begin{propriete}[Le présent de l'indicatif : rappel de conjugaison]
Le présent décrit une habitude, un fait général, une vérité permanente.
\begin{center}
\begin{tabular}{|l|l|l|l|}\hline
\rowcolor{popblueL} Personne & \esp{fumar} (-ar) & \esp{beber} (-er) & \esp{consumir} (-ir) \\ \hline
\esp{yo} & \esp{fumo} & \esp{bebo} & \esp{consumo} \\ \hline
\esp{tú} & \esp{fumas} & \esp{bebes} & \esp{consumes} \\ \hline
\esp{él/ella/usted} & \esp{fuma} & \esp{bebe} & \esp{consume} \\ \hline
\esp{nosotros/-as} & \esp{fumamos} & \esp{bebemos} & \esp{consumimos} \\ \hline
\esp{vosotros/-as} & \esp{fumáis} & \esp{bebéis} & \esp{consumís} \\ \hline
\esp{ellos/ellas/ustedes} & \esp{fuman} & \esp{beben} & \esp{consumen} \\ \hline
\end{tabular}
\end{center}
Verbes utiles irréguliers au présent : \esp{tener} (\esp{tengo, tienes, tiene, tenemos, tenéis, tienen}), \esp{poder} (\esp{puedo, puedes, puede...}), \esp{empezar} (\esp{empiezo, empiezas, empieza...}).
\end{propriete}

\begin{propriete}[Estar + gérondif : la formation]
La structure \esp{estar + gerundio} décrit une action EN TRAIN de se passer au moment où l'on parle.

\textbf{1. Le verbe} \esp{estar} \textbf{au présent} :
\begin{center}
\begin{tabular}{|l|l|}\hline
\rowcolor{popgreenL} Personne & \esp{estar} \\ \hline
\esp{yo} & \esp{estoy} \\ \hline
\esp{tú} & \esp{estás} \\ \hline
\esp{él/ella/usted} & \esp{está} \\ \hline
\esp{nosotros/-as} & \esp{estamos} \\ \hline
\esp{vosotros/-as} & \esp{estáis} \\ \hline
\esp{ellos/ellas/ustedes} & \esp{están} \\ \hline
\end{tabular}
\end{center}

\textbf{2. Le gérondif} : radical + \esp{-ando} (verbes en -ar) ; radical + \esp{-iendo} (verbes en -er/-ir).

\begin{center}
\begin{tabular}{|l|l|l|}\hline
\rowcolor{popgreenL} Infinitivo & Gerundio & Remarque \\ \hline
\esp{fumar} & \esp{fumando} & régulier \\ \hline
\esp{beber} & \esp{bebiendo} & régulier \\ \hline
\esp{escribir} & \esp{escribiendo} & régulier \\ \hline
\esp{leer} & \esp{leyendo} & i $\rightarrow$ y entre voyelles \\ \hline
\esp{dormir} & \esp{durmiendo} & o $\rightarrow$ u \\ \hline
\esp{decir} & \esp{diciendo} & e $\rightarrow$ i \\ \hline
\esp{pedir} & \esp{pidiendo} & e $\rightarrow$ i \\ \hline
\end{tabular}
\end{center}
Le gérondif est \textbf{invariable} : il ne s'accorde jamais avec le sujet. \esp{Las chicas están fumando.} (et non « fumandas »).
\end{propriete}

\begin{methode}[Choisir entre présent simple et estar + gérondif]
\begin{enumerate}
\item \textbf{Présent de l'indicatif} : pour une habitude, un fait général, une vérité permanente. \esp{Muchos jóvenes beben alcohol.} « Beaucoup de jeunes boivent de l'alcool » (en général).
\item \textbf{Estar + gérondif} : pour une action EN TRAIN de se passer au moment où l'on parle. \esp{En este momento, un alumno está fumando en el patio.} « En ce moment, un élève est en train de fumer dans la cour. »
\item \textbf{Indices de temps} : si la phrase contient \esp{ahora, ahora mismo, en este momento, mientras}, utilise presque toujours \esp{estar + gerundio}. Si elle contient \esp{siempre, todos los días, normalmente, a veces}, utilise le présent simple.
\end{enumerate}
\end{methode}

\begin{popfigure}
\begin{tikzpicture}[node distance=0.6cm and 1.6cm]
\node[draw=popblue, fill=popblueL, rounded corners, align=center, text width=4.2cm] (q) {\textbf{Qué describes ?}\\ ¿Qué describes?};
\node[draw=popgreen, fill=popgreenL, rounded corners, align=center, text width=4.6cm, below left=0.7cm and -0.4cm of q] (hab) {\textbf{Habitude / fait général}\\ \esp{todos los días, siempre}\\ présent simple\\ \esp{Los jóvenes beben.}};
\node[draw=poporange, fill=poporangeL, rounded corners, align=center, text width=4.6cm, below right=0.7cm and -0.4cm of q] (acc) {\textbf{Action en cours}\\ \esp{ahora, en este momento}\\ \esp{estar + gerundio}\\ \esp{Están bebiendo.}};
\draw[-{Stealth}, thick, popgreen] (q) -- (hab);
\draw[-{Stealth}, thick, poporange] (q) -- (acc);
\end{tikzpicture}
\legende{Choisir entre le présent simple et \esp{estar + gerundio}}
\end{popfigure}

\begin{exoresolu}[Transformation : de l'habitude à l'action en cours]
\textbf{Énoncé} : Transforme les phrases suivantes en utilisant \esp{estar + gerundio} pour décrire une scène qui se passe maintenant.
\begin{enumerate}
\item \esp{Los jóvenes beben alcohol en la fiesta.}
\item \esp{Un alumno insulta a su compañero.}
\item \esp{Los profesores hablan del problema.}
\item \esp{La víctima duerme mal.}
\end{enumerate}
\tcblower
\textbf{Raisonnement} : pour chaque phrase, on conjugue \esp{estar} selon le sujet, puis on forme le gérondif du verbe principal. Attention aux irréguliers.

\textbf{Solution} :
\begin{enumerate}
\item \esp{Los jóvenes están bebiendo alcohol en la fiesta.} « Les jeunes sont en train de boire de l'alcool à la fête. » Sujet pluriel $\rightarrow$ \esp{están} ; \esp{beber $\rightarrow$ bebiendo} (régulier).
\item \esp{Un alumno está insultando a su compañero.} « Un élève est en train d'insulter son camarade. » Sujet singulier $\rightarrow$ \esp{está} ; \esp{insultar $\rightarrow$ insultando}.
\item \esp{Los profesores están hablando del problema.} « Les professeurs sont en train de parler du problème. » \esp{hablar $\rightarrow$ hablando}.
\item \esp{La víctima está durmiendo mal.} « La victime est en train de mal dormir. » Attention : \esp{dormir} a un gérondif irrégulier : \esp{durmiendo} (o $\rightarrow$ u), et non « dormiendo ».
\end{enumerate}
\textbf{Conclusion} : la structure est toujours \esp{sujeto + estar conjugado + gerundio}. Le gérondif est invariable : il ne s'accorde jamais avec le sujet.
\end{exoresolu}

\courssub{Thème : traduire du français vers l'espagnol}

\begin{methode}[Réussir le thème sur ce thème de société]
\begin{enumerate}
\item \textbf{Analyser la phrase française} : sujet ? verbe ? temps ? Une habitude = présent simple ; une action en cours = \esp{estar + gerundio}.
\item \textbf{Traduire le sujet avec son article} : en espagnol, l'article est obligatoire devant un nom général : « la drogue » = \esp{la droga}, « les jeunes » = \esp{los jóvenes}.
\item \textbf{Choisir le bon verbe} : « fumer » = \esp{fumar} ; « consommer de la drogue » = \esp{consumir drogas} ; « harceler » = \esp{acosar}.
\item \textbf{Placer les accents} : \esp{adicción, violencia, atención, solución}. (Attention : \esp{violencia} ne prend PAS d'accent, c'est un mot grave terminé par une voyelle.)
\item \textbf{Relire} : accord sujet-verbe, accord article-nom (genre et nombre), ponctuation espagnole (¿ ¡) si phrase interrogative ou exclamative.
\end{enumerate}
\textbf{Piège majeur} : « drogue » se traduit \esp{la droga} (féminin, sans accent). « Une addiction » = \esp{una adicción} (avec accent sur la ó).
\end{methode}

\begin{exoresolu}[Thème : phrases à traduire]
\textbf{Énoncé} : Traduis en espagnol :
\begin{enumerate}
\item La violence est un problème grave dans certains collèges.
\item Les jeunes de ce quartier consomment de la drogue.
\item En ce moment, le directeur parle avec les parents de l'élève.
\item Le harcèlement scolaire provoque la peur et la tristesse.
\end{enumerate}
\tcblower
\textbf{Raisonnement et solution} :
\begin{enumerate}
\item Phrase générale $\rightarrow$ présent simple. « violence » = \esp{la violencia} (féminin, sans accent). \esp{La violencia es un problema grave en algunos colegios.} Justification : « certains » devant un nom masculin pluriel = \esp{algunos}.
\item Habitude $\rightarrow$ présent. « ce quartier » = \esp{este barrio}. \esp{Los jóvenes de este barrio consumen drogas.} Justification : \esp{consumir} est régulier ; « de la drogue » en général se traduit sans article : \esp{drogas}.
\item « En ce moment » $\rightarrow$ \esp{estar + gerundio}. \esp{En este momento, el director está hablando con los padres del alumno.} Justification : \esp{de + el = del} (contraction obligatoire) ; \esp{hablar $\rightarrow$ hablando}.
\item Fait général $\rightarrow$ présent. \esp{El acoso escolar provoca el miedo y la tristeza.} Justification : « la peur » = \esp{el miedo} (masculin !) ; « la tristesse » = \esp{la tristeza}.
\end{enumerate}
\textbf{Conclusion} : chaque phrase est vérifiée : article, genre, temps, contraction \esp{del}, accents.
\end{exoresolu}

\courssub{Rédiger un court paragraphe descriptif}

\begin{methode}[Rédiger un paragraphe de 5 à 8 lignes sur une conduite à risque]
\begin{enumerate}
\item \textbf{Phrase d'introduction} : présente le phénomène. \esp{En mi colegio existe un problema : ...}
\item \textbf{Description} : qui ? où ? quand ? Utilise le présent et une scène avec \esp{estar + gerundio}. \esp{Algunos alumnos están fumando detrás de las aulas.}
\item \textbf{Conséquences} : \esp{Por eso... / Como resultado...}
\item \textbf{Solution ou opinion personnelle} : \esp{En mi opinión, debemos... / Es necesario + infinitivo.}
\item \textbf{Conclusion courte} : \esp{Juntos podemos luchar contra este problema.}
\end{enumerate}
\textbf{Lexique utile} : \esp{proteger, prevenir, luchar contra, informar, ayudar a las víctimas, es peligroso, es necesario}.
\end{methode}

\begin{exoresolu}[Rédaction guidée : un paragraphe sur l'absentéisme]
\textbf{Énoncé} : Rédige un paragraphe de 5 à 7 phrases sur l'absentéisme dans ton école : présente le problème, décris une scène, donne une conséquence et propose une solution.
\tcblower
\textbf{Plan suivi} : introduction (phrase 1), description avec \esp{estar + gerundio} (phrases 2-3), conséquence (phrase 4), solution (phrases 5-6), conclusion (phrase 7).

\textbf{Production modèle} :

\esp{En mi colegio de Douala existe un problema grave : el absentismo. Todos los días, algunos alumnos faltan a clase sin permiso. Ahora mismo, dos estudiantes están caminando por el mercado en lugar de estar en el aula. Como resultado, estos alumnos no aprenden bien y repiten el año. En mi opinión, los profesores deben hablar con los padres de estos jóvenes. También es necesario informar a los alumnos sobre la importancia de la educación. Juntos podemos luchar contra el absentismo.}

« Dans mon collège de Douala, il existe un problème grave : l'absentéisme. Tous les jours, certains élèves manquent les cours sans autorisation. En ce moment même, deux élèves se promènent au marché au lieu d'être en classe. En conséquence, ces élèves n'apprennent pas bien et redoublent. À mon avis, les professeurs doivent parler avec les parents de ces jeunes. Il est aussi nécessaire d'informer les élèves sur l'importance de l'éducation. Ensemble, nous pouvons lutter contre l'absentéisme. »

\textbf{Justifications grammaticales} : \esp{están caminando} = action en cours ; \esp{en lugar de + infinitivo} = « au lieu de » ; \esp{es necesario + infinitivo} = structure impersonnelle ; \esp{educación} porte un accent (mot aigu terminé par -n).
\end{exoresolu}

\begin{experience}[Débat de classe : ¿Cómo proteger a los jóvenes?]
Par groupes de quatre, préparez un mini-débat en espagnol :
\begin{enumerate}
\item Chaque groupe choisit une conduite à risque (\esp{el alcohol, el sexting, el acoso escolar, el absentismo}).
\item Chaque élève prépare deux phrases : une description avec \esp{estar + gerundio} (\esp{En este momento, unos jóvenes están...}) et une proposition de solution avec \esp{es necesario + infinitivo} ou \esp{debemos + infinitivo}.
\item Un porte-parole présente les idées du groupe à la classe ; les autres groupes posent une question avec \esp{¿Qué podemos hacer para...?}
\end{enumerate}
Critères de réussite : au moins un gérondif correct, un connecteur logique, une prononciation claire.
\end{experience}

\begin{attention}[Les 5 erreurs qui coûtent des points à l'examen]
\begin{enumerate}
\item \textbf{Oublier l'accent} sur les mots aigus en -ción : on écrit \esp{adicción, atención, solución, educación}, jamais « adiccion ». Sans accent, le mot est fautif. Mais \esp{violencia} ne prend pas d'accent.
\item \textbf{Le faux ami « droga »} : c'est un nom FÉMININ : \esp{la droga, las drogas}. Et « drogué » se dit \esp{drogadicto}, pas « drogado » dans ce sens.
\item \textbf{Le gérondif mal formé} : le gérondif de \esp{fumar} est \esp{fumando} ; mais attention à \esp{dormir $\rightarrow$ durmiendo} et \esp{decir $\rightarrow$ diciendo}. Ne jamais écrire « dormiendo ».
\item \textbf{L'accord sujet-verbe} : \esp{los jóvenes consumen} (et non « consume »). Le verbe s'accorde avec son sujet, même si un complément pluriel suit : \esp{el acoso provoca problemas}.
\item \textbf{Le mauvais choix de temps} : une habitude ou un fait général = présent simple (\esp{los jóvenes beben}) ; une scène en cours = \esp{estar + gerundio} (\esp{están bebiendo}). Mélanger les deux dans la même description est sanctionné.
\end{enumerate}
\end{attention}

\begin{aretenir}[L'essentiel de la leçon]
\begin{itemize}
\item Le lexique des conduites à risque : \esp{la droga, el alcohol, el tabaco, la violencia, el acoso escolar, el sexting, el absentismo, la delincuencia, la adicción}.
\item Milieu scolaire : \esp{en el colegio, en el patio, en el aula} ; milieu social : \esp{en el barrio, en la calle, en el mercado} ; milieu numérique : \esp{en las redes sociales, por el móvil}.
\item Habitude ou fait général $\rightarrow$ présent de l'indicatif : \esp{Muchos jóvenes fuman.}
\item Action en cours $\rightarrow$ \esp{estar + gerundio} : \esp{Ahora están fumando detrás del colegio.}
\item Gérondifs irréguliers à connaître par cœur : \esp{leyendo, durmiendo, diciendo, pidiendo}.
\item Pour décrire une conduite : nommer $\rightarrow$ définir $\rightarrow$ situer $\rightarrow$ montrer la conséquence.
\end{itemize}
\end{aretenir}

\begin{exercice}[Pour t'entraîner]
\begin{enumerate}
\item Conjugue au présent : \esp{fumar (yo)}, \esp{beber (ellos)}, \esp{consumir (nosotros)}, \esp{tener (tú)}, \esp{empezar (ella)}.
\item Mets à la forme \esp{estar + gerundio} : a) \esp{El alumno lee un folleto sobre la droga.} b) \esp{Los jóvenes duermen en clase.} c) \esp{La profesora dice la verdad.}
\item Traduis en espagnol : a) « Le tabac est dangereux pour la santé. » b) « En ce moment, les élèves écoutent une causerie sur la violence. » c) « Certains jeunes manquent les cours tous les jours. »
\item Rédige un paragraphe de 5 à 7 phrases sur \esp{el acoso escolar} dans ton école (problème, scène, conséquence, solution).
\end{enumerate}
\end{exercice}

\begin{corrige}[Exercice « Pour t'entraîner »]
\begin{enumerate}
\item \esp{fumo ; beben ; consumimos ; tienes ; empieza}. Rappel : \esp{tener} et \esp{empezar} sont irréguliers (\esp{tengo/tienes}, \esp{empiezo/empieza}).
\item a) \esp{El alumno está leyendo un folleto sobre la droga.} (gérondif irrégulier \esp{leyendo}) ; b) \esp{Los jóvenes están durmiendo en clase.} (\esp{durmiendo}, o $\rightarrow$ u) ; c) \esp{La profesora está diciendo la verdad.} (\esp{diciendo}, e $\rightarrow$ i).
\item a) \esp{El tabaco es peligroso para la salud.} (« pour » de destination/but = \esp{para}) ; b) \esp{En este momento, los alumnos están escuchando una charla sobre la violencia.} ; c) \esp{Algunos jóvenes faltan a clase todos los días.} (habitude $\rightarrow$ présent simple).
\item Modèle attendu : introduction (\esp{En mi colegio existe...}), scène avec \esp{estar + gerundio} (\esp{Unos alumnos están insultando a un compañero en el patio.}), conséquence (\esp{Como resultado, la víctima tiene miedo.}), solution (\esp{Es necesario informar y proteger a las víctimas.}), conclusion (\esp{Juntos podemos luchar contra el acoso escolar.}).
\end{enumerate}
\end{corrige}

\newpage

\coursec{Exercices}

\courssub{Je vérifie mes connaissances}

\begin{exercice}[QCM : le bon mot]
Pour chaque question, entoure la bonne réponse.
\begin{enumerate}
\item \esp{El acoso escolar} désigne :
\begin{enumerate}
\item l'absentéisme scolaire \item le harcèlement entre élèves \item la consommation de drogues
\end{enumerate}
\item \esp{El sexting} consiste à :
\begin{enumerate}
\item envoyer des images intimes par téléphone \item sécher les cours \item fumer en cachette
\end{enumerate}
\item \esp{El absentismo} signifie :
\begin{enumerate}
\item la violence physique \item l'addiction au tabac \item le fait de manquer souvent les cours
\end{enumerate}
\item \esp{La adicción} est :
\begin{enumerate}
\item une dépendance à une substance ou à un comportement \item une punition scolaire \item un jeu vidéo
\end{enumerate}
\end{enumerate}
\end{exercice}

\begin{exercice}[Vrai ou faux ? Justifie]
Réponds par \esp{verdadero} ou \esp{falso} et justifie ta réponse par une phrase complète en espagnol.
\begin{enumerate}
\item \esp{El tabaco no es una conducta de riesgo porque es legal.}
\item \esp{El acoso escolar puede ocurrir también en las redes sociales.}
\item \esp{Los jóvenes que consumen drogas nunca tienen problemas en la escuela.}
\item \esp{Hablar con los padres y los profesores puede ayudar a prevenir las conductas de riesgo.}
\end{enumerate}
\end{exercice}

\begin{exercice}[Vocabulaire : le mot intrus]
Dans chaque série, barre le mot qui n'appartient pas au champ lexical des conduites à risque et explique ton choix en français.
\begin{enumerate}
\item \esp{la droga / el alcohol / el tabaco / la lectura}
\item \esp{la violencia / el acoso / la amistad / la agresión}
\item \esp{el absentismo / la delincuencia / el deporte / el sexting}
\item \esp{la adicción / la dependencia / la salud / el consumo}
\end{enumerate}
\end{exercice}

\begin{exercice}[Conjugaison : présent ou \esp{estar} + gérondif]
Complète les phrases avec le verbe entre parenthèses au présent de l'indicatif ou à la périphrase \esp{estar} + gérondif, selon le sens (habitude ou action en cours).
\begin{enumerate}
\item Ahora mismo, los alumnos \esp{\_\_\_\_\_\_\_\_} (fumar) detrás del gimnasio.
\item Muchos jóvenes \esp{\_\_\_\_\_\_\_\_} (beber) alcohol los fines de semana.
\item ¡Mira! El director \esp{\_\_\_\_\_\_\_\_} (hablar) con los padres del alumno agresor.
\item Mi hermano nunca \esp{\_\_\_\_\_\_\_\_} (llegar) tarde a clase.
\item En este momento, la policía \esp{\_\_\_\_\_\_\_\_} (investigar) una red de droga cerca del liceo.
\item El acoso escolar \esp{\_\_\_\_\_\_\_\_} (preocupar) mucho a los profesores.
\end{enumerate}
\end{exercice}

\courssub{Je m'entraîne}

\begin{exercice}[Appariements : conduites et définitions]
Associe chaque conduite (1 à 6) à sa définition (a à f).
\begin{enumerate}
\item \esp{la droga}
\item \esp{el sexting}
\item \esp{el absentismo}
\item \esp{el acoso escolar}
\item \esp{la delincuencia}
\item \esp{la adicción}
\end{enumerate}
\begin{enumerate}[label=\alph*.]
\item \esp{Faltar a clase con frecuencia sin justificación.}
\item \esp{Enviar fotos o mensajes íntimos por el móvil o Internet.}
\item \esp{Insultar, golpear o humillar repetidamente a un compañero.}
\item \esp{Cometer actos contra la ley, como robar o agredir.}
\item \esp{Sustancia ilegal que altera el cuerpo y la mente.}
\item \esp{Necesidad irresistible de consumir algo o de repetir un comportamiento.}
\end{enumerate}
\end{exercice}

\begin{exercice}[Texte à trous]
Complète le texte avec les mots suivants : \esp{alcohol -- acoso -- riesgo -- absentismo -- adicción -- violencia}.

\begin{textebox}[En el patio del liceo]
En muchos liceos, los jóvenes adoptan conductas de \_\_\_\_\_\_\_\_. Algunos consumen \_\_\_\_\_\_\_\_ durante las fiestas y desarrollan después una \_\_\_\_\_\_\_\_ difícil de curar. Otros sufren \_\_\_\_\_\_\_\_ escolar : sus compañeros se burlan de ellos todos los días. El \_\_\_\_\_\_\_\_ también es un problema grave : muchos alumnos faltan a clase y se quedan en la calle. Finalmente, la \_\_\_\_\_\_\_\_ entre bandas de jóvenes preocupa a los padres y a los profesores.
\end{textebox}
\end{exercice}

\begin{exercice}[Transformations : de l'habitude à l'action en cours]
Transforme chaque phrase : remplace le présent d'habitude par la périphrase \esp{estar} + gérondif en ajoutant \esp{ahora} ou \esp{en este momento}, comme dans le modèle.

Modèle : \esp{Los jóvenes fuman en el patio.} $\rightarrow$ \esp{Ahora los jóvenes están fumando en el patio.}
\begin{enumerate}
\item \esp{Los alumnos beben alcohol en la fiesta.}
\item \esp{El agresor insulta a su compañero.}
\item \esp{La policía vigila el barrio.}
\item \esp{Los profesores hablan del problema.}
\item \esp{Mi amigo envía mensajes peligrosos.}
\end{enumerate}
\end{exercice}

\begin{exercice}[Traduction : thème]
Traduis en espagnol les phrases suivantes.
\begin{enumerate}
\item Les jeunes sont en train de fumer dans la cour.
\item Le harcèlement inquiète les parents.
\item Beaucoup d'élèves manquent les cours à cause de la drogue.
\item En ce moment, le directeur parle avec les victimes.
\item L'alcool détruit la santé des adolescents.
\end{enumerate}
\end{exercice}

\courssub{Je me prépare au Bac}

\begin{exercice}[Compréhension de texte type Bac]
Lis le texte suivant, puis réponds aux questions en espagnol, avec des phrases complètes.

\begin{textebox}[El acoso escolar : un problema silencioso]
En muchos colegios de las grandes ciudades, algunos alumnos viven una pesadilla diaria : el acoso escolar. Todo empieza con burlas en el patio o insultos en los grupos de WhatsApp. La víctima, a menudo tímida, no se atreve a hablar. Poco a poco, deja de participar en clase, llega tarde o falta a los cursos : es el absentismo, una señal de alarma que los profesores deben detectar rápidamente.

Hoy, el acoso ya no se queda en la puerta del colegio : continúa en las redes sociales, día y noche. Algunos jóvenes reciben mensajes de amenaza o ven sus fotos privadas circulando sin su permiso. Esta violencia psicológica puede empujar a los adolescentes hacia otras conductas de riesgo, como el consumo de alcohol o de drogas, para olvidar su sufrimiento.

Sin embargo, existen soluciones. Cuando un alumno habla con un profesor, con sus padres o con un consejero, la situación puede cambiar. En varios establecimientos, los propios estudiantes crean clubes contra el acoso y organizan campañas de sensibilización. La escuela, la familia y la sociedad deben trabajar juntas para proteger a los jóvenes.
\end{textebox}

\textbf{Questions de comprensión}
\begin{enumerate}
\item \esp{¿Cómo empieza normalmente el acoso escolar, según el texto?}
\item \esp{¿Qué señal de alarma deben detectar los profesores?}
\item \esp{¿Dónde continúa el acoso hoy en día?}
\item \esp{¿Qué otras conductas de riesgo puede provocar el acoso?}
\item \esp{¿Qué soluciones propone el texto?}
\end{enumerate}

\textbf{Vrai ou faux} — Justifie chaque réponse par une citation du texte.
\begin{enumerate}
\item \esp{Las víctimas del acoso hablan fácilmente de su problema.}
\item \esp{El acoso termina cuando el alumno sale del colegio.}
\item \esp{Los estudiantes pueden participar en la lucha contra el acoso.}
\end{enumerate}

\textbf{Questions de langue}
\begin{enumerate}
\item Trouve dans le texte deux mots du champ lexical des conduites à risque et traduis-les en français.
\item Réécris la phrase \esp{Algunos jóvenes reciben mensajes de amenaza} en utilisant \esp{estar} + gérondif avec \esp{ahora}.
\end{enumerate}
\end{exercice}

\begin{exercice}[Version et thème]
\textbf{A. Version.} Traduis en français le paragraphe suivant.

\begin{textebox}[La delincuencia juvenil]
En algunos barrios pobres, la delincuencia juvenil aumenta poco a poco. Los jóvenes sin trabajo ni esperanza roban teléfonos o venden droga para sobrevivir. Las autoridades intentan reaccionar : crean centros deportivos y programas de formación profesional. Pero los especialistas repiten que la verdadera solución es la educación y el diálogo con las familias.
\end{textebox}

\textbf{B. Thème.} Traduis en espagnol.
\begin{enumerate}
\item Dans mon quartier, les jeunes sont en train de boire de l'alcool près du marché.
\item Le harcèlement scolaire est un problème grave dans les lycées.
\item Les parents et les professeurs doivent parler avec les adolescents.
\item En ce moment, la victime pleure dans la cour de récréation.
\item La drogue et le tabac détruisent la vie des jeunes.
\end{enumerate}
\end{exercice}

\begin{exercice}[Production écrite type Bac]
Rédige en espagnol un paragraphe de \textbf{8 à 10 lignes} (environ 80 mots) intitulé :

\begin{center}
\esp{\textbf{Las conductas de riesgo en mi barrio}}
\end{center}

Consignes obligatoires :
\begin{itemize}
\item utilise au moins \textbf{5 mots} du lexique de la leçon (\esp{droga, alcohol, tabaco, violencia, acoso escolar, sexting, absentismo, delincuencia, adicción}) et souligne-les ;
\item emploie au moins \textbf{2 périphrases} \esp{estar} + gérondif pour décrire des actions en cours ;
\item termine par une phrase qui propose une solution.
\end{itemize}
\end{exercice}

\newpage

\coursec{Corrigés des exercices}

\begin{corrige}[Exercice 1 — QCM : le bon mot]
\textbf{Réponses :}
\begin{enumerate}
\item \textbf{b} : \esp{el acoso escolar} désigne le harcèlement entre élèves (insultes, coups, humiliations répétées).
\item \textbf{a} : \esp{el sexting} consiste à envoyer des images ou des messages intimes par téléphone ou par Internet.
\item \textbf{c} : \esp{el absentismo} signifie le fait de manquer souvent les cours sans justification.
\item \textbf{a} : \esp{la adicción} est une dépendance à une substance (drogue, alcool, tabac) ou à un comportement (jeux, réseaux sociaux).
\end{enumerate}
\textbf{Astuce :} attention au faux ami : \esp{la droga} = la drogue, mais \esp{la droguería} = la pharmacie/droguerie ; et \esp{asistir} signifie « assister à » (être présent), d'où \esp{el absentismo} = l'absentéisme.
\end{corrige}

\begin{corrige}[Exercice 2 — Vrai ou faux ? Justifie]
\begin{enumerate}
\item \textbf{Falso.} \esp{El tabaco es una conducta de riesgo aunque es legal, porque daña gravemente la salud y crea adicción.} (Le tabac est une conduite à risque bien qu'il soit légal, car il nuit gravement à la santé et crée une dépendance.)
\item \textbf{Verdadero.} \esp{El acoso escolar puede ocurrir también en las redes sociales, por ejemplo con insultos o amenazas en WhatsApp.} (Le harcèlement scolaire peut aussi se produire sur les réseaux sociaux.)
\item \textbf{Falso.} \esp{Los jóvenes que consumen drogas tienen a menudo problemas en la escuela, como el absentismo o las malas notas.} (Les jeunes qui consomment de la drogue ont souvent des problèmes à l'école.)
\item \textbf{Verdadero.} \esp{Hablar con los padres y los profesores puede ayudar a prevenir las conductas de riesgo porque el diálogo protege a los jóvenes.} (Parler avec les parents et les professeurs peut aider à prévenir les conduites à risque.)
\end{enumerate}
\textbf{Point de langue :} pour justifier, on utilise souvent \esp{porque} + indicatif ; on peut aussi employer \esp{ya que} ou \esp{puesto que} dans un texte écrit.
\end{corrige}

\begin{corrige}[Exercice 3 — Vocabulaire : le mot intrus]
\begin{enumerate}
\item Intrus : \esp{la lectura} (la lecture). C'est une activité saine et éducative, pas une substance à risque comme la drogue, l'alcool ou le tabac.
\item Intrus : \esp{la amistad} (l'amitié). C'est une relation positive ; la violence, le harcèlement et l'agression sont des comportements négatifs.
\item Intrus : \esp{el deporte} (le sport). C'est une activité bénéfique pour la santé, contrairement à l'absentéisme, à la délinquance et au sexting.
\item Intrus : \esp{la salud} (la santé). C'est un bien à protéger ; l'addiction, la dépendance et la consommation renvoient aux conduites à risque.
\end{enumerate}
\textbf{Remarque :} cet exercice montre qu'un champ lexical s'organise par opposition : les mots positifs (\esp{salud, amistad, deporte, lectura}) sont les « antidotes » des conduites à risque.
\end{corrige}

\begin{corrige}[Exercice 4 — Conjugaison : présent ou \esp{estar} + gérondif]
\begin{enumerate}
\item \esp{Ahora mismo, los alumnos \textbf{están fumando} detrás del gimnasio.} — Action en cours (\esp{ahora mismo}) : périphrase \esp{estar} + gérondif.
\item \esp{Muchos jóvenes \textbf{beben} alcohol los fines de semana.} — Habitude (\esp{los fines de semana}) : présent simple.
\item \esp{¡Mira! El director \textbf{está hablando} con los padres del alumno agresor.} — Action en train de se dérouler sous nos yeux.
\item \esp{Mi hermano nunca \textbf{llega} tarde a clase.} — Habitude (\esp{nunca}) : présent simple.
\item \esp{En este momento, la policía \textbf{está investigando} una red de droga cerca del liceo.} — Action en cours.
\item \esp{El acoso escolar \textbf{preocupa} mucho a los profesores.} — Vérité générale : présent simple.
\end{enumerate}
\textbf{Rappel :} les marqueurs de l'action en cours sont \esp{ahora, ahora mismo, en este momento, ¡mira!} ; les marqueurs de l'habitude sont \esp{siempre, nunca, normalmente, los fines de semana}. Gérondifs réguliers : \esp{fumar} $\rightarrow$ \esp{fumando} ; \esp{beber} $\rightarrow$ \esp{bebiendo} ; \esp{investigar} $\rightarrow$ \esp{investigando}.
\end{corrige}

\begin{corrige}[Exercice 5 — Appariements : conduites et définitions]
\begin{center}
\begin{tabular}{|l|l|}
\hline
\rowcolor{popblueL} \textbf{Conduite} & \textbf{Définition} \\
\hline
1. \esp{la droga} & \textbf{e} : sustancia ilegal que altera el cuerpo y la mente \\
\hline
2. \esp{el sexting} & \textbf{b} : enviar fotos o mensajes íntimos por el móvil o Internet \\
\hline
3. \esp{el absentismo} & \textbf{a} : faltar a clase con frecuencia sin justificación \\
\hline
4. \esp{el acoso escolar} & \textbf{c} : insultar, golpear o humillar repetidamente a un compañero \\
\hline
5. \esp{la delincuencia} & \textbf{d} : cometer actos contra la ley, como robar o agredir \\
\hline
6. \esp{la adicción} & \textbf{f} : necesidad irresistible de consumir algo o de repetir un comportamiento \\
\hline
\end{tabular}
\end{center}
\textbf{À retenir :} \esp{faltar a clase} = manquer les cours ; \esp{sin justificación} = sans justification ; \esp{contra la ley} = contre la loi.
\end{corrige}

\begin{corrige}[Exercice 6 — Texte à trous]
Texte complété :
\begin{quote}
\esp{En muchos liceos, los jóvenes adoptan conductas de \textbf{riesgo}. Algunos consumen \textbf{alcohol} durante las fiestas y desarrollan después una \textbf{adicción} difícil de curar. Otros sufren \textbf{acoso} escolar : sus compañeros se burlan de ellos todos los días. El \textbf{absentismo} también es un problema grave : muchos alumnos faltan a clase y se quedan en la calle. Finalmente, la \textbf{violencia} entre bandas de jóvenes preocupa a los padres y a los profesores.}
\end{quote}
\textbf{Traduction :} « Dans de nombreux lycées, les jeunes adoptent des conduites à risque. Certains consomment de l'alcool pendant les fêtes et développent ensuite une addiction difficile à guérir. D'autres subissent le harcèlement scolaire : leurs camarades se moquent d'eux tous les jours. L'absentéisme est aussi un problème grave : beaucoup d'élèves manquent les cours et restent dans la rue. Enfin, la violence entre bandes de jeunes inquiète les parents et les professeurs. »
\textbf{Justification :} \esp{conductas de riesgo} est une expression figée ; \esp{una adicción} exige un nom féminin singulier ; \esp{acoso escolar} est le syntagme appris dans la leçon ; \esp{el absentismo} reprend l'idée de \esp{faltar a clase} ; \esp{la violencia} est justifiée par \esp{bandas}.
\end{corrige}

\begin{corrige}[Exercice 7 — Transformations : de l'habitude à l'action en cours]
\begin{enumerate}
\item \esp{Ahora los alumnos \textbf{están bebiendo} alcohol en la fiesta.}
\item \esp{En este momento el agresor \textbf{está insultando} a su compañero.}
\item \esp{Ahora la policía \textbf{está vigilando} el barrio.}
\item \esp{En este momento los profesores \textbf{están hablando} del problema.}
\item \esp{Ahora mi amigo \textbf{está enviando} mensajes peligrosos.}
\end{enumerate}
\textbf{Points de vigilance :}
\begin{itemize}
\item \esp{beber} $\rightarrow$ \esp{bebiendo} (gérondif régulier en \esp{-iendo}) ;
\item \esp{enviar} $\rightarrow$ \esp{enviando} (pas de changement orthographique au gérondif, contrairement au présent \esp{envío}) ;
\item ne jamais écrire \esp{*están bebendo} : la terminaison des verbes en \esp{-er/-ir} est \esp{-iendo} ;
\item le verbe \esp{estar} s'accorde avec le sujet : \esp{está} (singulier) / \esp{están} (pluriel).
\end{itemize}
\end{corrige}

\begin{corrige}[Exercice 8 — Traduction : thème]
\begin{enumerate}
\item \esp{Los jóvenes están fumando en el patio.} — « sont en train de fumer » = action en cours : \esp{estar} + gérondif.
\item \esp{El acoso escolar preocupa a los padres.} — vérité générale : présent simple ; \esp{preocupar a alguien} = inquiéter quelqu'un.
\item \esp{Muchos alumnos faltan a clase a causa de la droga.} — « à cause de » = \esp{a causa de} ; \esp{faltar a clase} = manquer les cours.
\item \esp{En este momento, el director está hablando con las víctimas.} — « en ce moment » déclenche la périphrase.
\item \esp{El alcohol destruye la salud de los adolescentes.} — \esp{destruir} $\rightarrow$ \esp{destruye} (verbe à changement orthographique : \esp{destruyo, destruyes, destruye}) ; « la santé » = \esp{la salud}.
\end{enumerate}
\textbf{Erreurs fréquentes à éviter :} ne pas traduire « en train de » par un mot séparé ; ne pas écrire \esp{*los jóvenes son fumando} : en espagnol, on n'utilise jamais \esp{ser} + gérondif, mais toujours \esp{estar} + gérondif.
\end{corrige}

\begin{corrige}[Exercice 9 — Compréhension de texte type Bac]
\textbf{Barème indicatif (sur 10) :} comprensión 4 pts (0,8 par question) ; vrai/faux 3 pts (1 par réponse justifiée) ; questions de langue 3 pts (1 + 2).

\textbf{A. Questions de comprensión — réponses modèles :}
\begin{enumerate}
\item \esp{Según el texto, el acoso escolar empieza normalmente con burlas en el patio o insultos en los grupos de WhatsApp.}
\item \esp{Los profesores deben detectar el absentismo : la víctima deja de participar en clase, llega tarde o falta a los cursos.}
\item \esp{Hoy en día, el acoso continúa en las redes sociales, día y noche.}
\item \esp{El acoso puede provocar otras conductas de riesgo, como el consumo de alcohol o de drogas.}
\item \esp{El texto propone hablar con un profesor, con los padres o con un consejero, crear clubes contra el acoso y organizar campañas de sensibilización.}
\end{enumerate}

\textbf{B. Vrai ou faux :}
\begin{enumerate}
\item \textbf{Falso.} \esp{« La víctima, a menudo tímida, no se atreve a hablar. »} (La victime n'ose pas parler.)
\item \textbf{Falso.} \esp{« El acoso ya no se queda en la puerta del colegio : continúa en las redes sociales, día y noche. »}
\item \textbf{Verdadero.} \esp{« Los propios estudiantes crean clubes contra el acoso y organizan campañas de sensibilización. »}
\end{enumerate}

\textbf{C. Questions de langue :}
\begin{enumerate}
\item Exemples : \esp{el absentismo} = l'absentéisme ; \esp{el consumo de drogas} = la consommation de drogues (autres possibilités : \esp{el acoso escolar} = le harcèlement scolaire, \esp{la violencia} = la violence).
\item \esp{Ahora algunos jóvenes están recibiendo mensajes de amenaza.} — \esp{recibir} $\rightarrow$ \esp{recibiendo} (gérondif régulier).
\end{enumerate}

\textbf{Points de vigilance :} répondre par des phrases complètes qui reprennent les mots de la question ; pour le vrai/faux, la citation du texte est obligatoire ; ne pas confondre \esp{atreverse a} (oser) et \esp{atreverse} seul.
\end{corrige}

\begin{corrige}[Exercice 10 — Version et thème]
\textbf{Barème indicatif (sur 10) :} version 5 pts ; thème 5 pts (1 pt par phrase).

\textbf{A. Version — traduction proposée :}

« Dans certains quartiers pauvres, la délinquance juvénile augmente peu à peu. Les jeunes sans travail ni espérance volent des téléphones ou vendent de la drogue pour survivre. Les autorités essaient de réagir : elles créent des centres sportifs et des programmes de formation professionnelle. Mais les spécialistes répètent que la véritable solution est l'éducation et le dialogue avec les familles. »

\textbf{Points de traduction :} \esp{poco a poco} = peu à peu ; \esp{sin trabajo ni esperanza} = sans travail ni espérance ; \esp{sobrevivir} = survivre ; \esp{intentan reaccionar} = essaient de réagir ; \esp{la verdadera solución} = la véritable solution.

\textbf{B. Thème — traductions proposées :}
\begin{enumerate}
\item \esp{En mi barrio, los jóvenes están bebiendo alcohol cerca del mercado.}
\item \esp{El acoso escolar es un problema grave en los liceos.}
\item \esp{Los padres y los profesores deben hablar con los adolescentes.}
\item \esp{En este momento, la víctima está llorando en el patio de recreo.}
\item \esp{La droga y el tabaco destruyen la vida de los jóvenes.}
\end{enumerate}

\textbf{Points de vigilance :} « sont en train de boire » = \esp{están bebiendo} ; « près du marché » = \esp{cerca del mercado} ; « doivent » = \esp{deben} + infinitif ; « la cour de récréation » = \esp{el patio de recreo} ; avec deux sujets (\esp{la droga y el tabaco}), le verbe est au pluriel : \esp{destruyen}.
\end{corrige}

\begin{corrige}[Exercice 11 — Production écrite type Bac]
\textbf{Barème indicatif (sur 10) :} respect des consignes (5 mots du lexique soulignés, 2 périphrases, solution finale) 4 pts ; correction grammaticale 3 pts ; richesse du vocabulaire et cohérence 3 pts.

\textbf{Proposition de rédaction modèle :}

\begin{quote}
\esp{En mi barrio, algunos jóvenes adoptan conductas peligrosas. Por la noche, un grupo de chicos \textbf{está fumando} \underline{tabaco} y \underline{drogas} cerca del mercado, y otros \textbf{están bebiendo} \underline{alcohol} delante de las tiendas. Además, el \underline{acoso escolar} existe también en mi liceo : algunos alumnos insultan a sus compañeros en los grupos de WhatsApp. A causa de estos problemas, muchos estudiantes sufren \underline{adicción} o faltan a clase : es el \underline{absentismo}. La \underline{violencia} entre bandas preocupa mucho a los padres. Para solucionar esta situación, creo que la escuela, la familia y las autoridades deben organizar campañas de sensibilización y actividades deportivas para los jóvenes.}
\end{quote}

(Environ 95 mots ; 6 mots du lexique soulignés ; 2 périphrases \esp{estar} + gérondif ; phrase finale de solution.)

\textbf{Traduction du modèle :} « Dans mon quartier, certains jeunes adoptent des conduites dangereuses. Le soir, un groupe de garçons est en train de fumer du tabac et de la drogue près du marché, et d'autres sont en train de boire de l'alcool devant les magasins. De plus, le harcèlement scolaire existe aussi dans mon lycée : certains élèves insultent leurs camarades dans les groupes WhatsApp. À cause de ces problèmes, beaucoup d'étudiants souffrent d'addiction ou manquent les cours : c'est l'absentéisme. La violence entre bandes inquiète beaucoup les parents. Pour résoudre cette situation, je pense que l'école, la famille et les autorités doivent organiser des campagnes de sensibilisation et des activités sportives pour les jeunes. »

\textbf{Points de vigilance :}
\begin{itemize}
\item vérifier l'accord de \esp{estar} avec le sujet (\esp{está} / \esp{están}) ;
\item gérondifs corrects : \esp{fumando, bebiendo} (jamais \esp{*bebendo}) ;
\item connecteurs utiles pour structurer : \esp{además, a causa de, para solucionar, creo que} ;
\item « je pense que » + indicatif : \esp{creo que la escuela debe...} ; avec \esp{no creo que}, il faudrait le subjonctif ;
\item compter ses mots et relire : accents (\esp{adicción, sensibilización}), \esp{ñ}, ponctuation.
\end{itemize}
\end{corrige}

\newpage

\coursec{Fiche bilan}

\begin{popfigure}
\begin{tikzpicture}[mindmap, grow cyclic, every node/.style={concept, text=white, align=center, font=\small}, concept color=popblue, level 1/.append style={level distance=3.4cm, sibling angle=51.4}, level 2/.append style={level distance=1.8cm, sibling angle=45, font=\scriptsize}]
\node[concept, font=\bfseries\small, minimum size=2.6cm, align=center]{Conductas\\ de riesgo\\ (1\iere{} A)}
  child[concept color=poporange] { node[align=center]{Lexique\\ clé}
    child { node[align=center]{droga,\\ alcohol,\\ tabaco} }
    child { node[align=center]{acoso,\\ sexting,\\ absentismo} }
  }
  child[concept color=popgreen] { node[align=center]{Milieu\\ scolaire}
    child { node[align=center]{aula, patio,\\ ba\~nos} }
  }
  child[concept color=poppurple] { node[align=center]{Milieu\\ social}
    child { node[align=center]{calle, redes\\ sociales, bares} }
  }
  child[concept color=popgold] { node[align=center]{Présent\\ de l'indic.}
    child { node[align=center]{réguliers +\\ irregulares} }
  }
  child[concept color=poppink] { node[align=center]{estar +\\ gérondif}
    child { node[align=center]{ando / iendo\\ = en train de} }
  }
  child[concept color=popgreen!70!black] { node[align=center]{Méthode\\ commentaire}
    child { node[align=center]{lire, relever,\\ analyser, juger} }
  }
  child[concept color=poporange!80!black] { node[align=center]{Production\\ guidée}
    child { node[align=center]{paragraphe\\ descriptif} }
  };
\end{tikzpicture}
\legende{Carte mentale de la leçon EA03 : les conduites à risque des jeunes}
\end{popfigure}

\begin{bilan}

\textbf{1. Lexique clé à connaître par cœur}

\begin{center}
\begin{tabularx}{\linewidth}{|X|X|X|X|}
\hline
\rowcolor{popblueL} Español & Français & Español & Français \\
\hline
la droga / drogarse & la drogue / se droguer & el acoso escolar & le harcèlement scolaire \\
\hline
el alcohol & l'alcool & el sexting & le sexting \\
\hline
el tabaco / fumar & le tabac / fumer & el absentismo & l'absentéisme \\
\hline
la violencia & la violence & la delincuencia & la délinquance \\
\hline
la adicción & l'addiction & el riesgo / arriesgar & le risque / risquer \\
\hline
las redes sociales & les réseaux sociaux & prevenir / la prevención & prévenir / la prévention \\
\hline
\end{tabularx}
\end{center}

\textbf{2. Grammaire à maîtriser}

\begin{center}
\begin{tabularx}{\linewidth}{|X|X|}
\hline
\rowcolor{popgreenL} Règle & Exemple \\
\hline
Présent de l'indicatif : faits habituels, vérités générales, descriptions & \esp{Muchos jóvenes fuman en el patio.} « Beaucoup de jeunes fument dans la cour. » \\
\hline
Radical changeant : e $\rightarrow$ ie, o $\rightarrow$ ue (sauf nous/vous) & \esp{pensar : pienso} ; \esp{dormir : duermo} \\
\hline
Première personne irrégulière : \esp{hago, salgo, pongo, conozco, sé} & \esp{Yo no conozco ese peligro.} \\
\hline
\esp{estar} + gérondif = action en cours (être en train de) & \esp{Los alumnos están hablando del problema.} \\
\hline
Gérondif : -ando (ar) / -iendo (er, ir) ; irréguliers : \esp{leyendo, durmiendo, diciendo} & \esp{Está leyendo un artículo sobre la droga.} \\
\hline
\end{tabularx}
\end{center}

\textbf{3. Méthodes express}

\begin{itemize}
\item \textbf{Commentaire de texte} : 1) lire deux fois et souligner les mots-clés ; 2) identifier thème, source, personnages ; 3) relever le lexique du risque ; 4) analyser la situation (où, qui, quoi) ; 5) donner un jugement personnel argumenté.
\item \textbf{Paragraphe descriptif} : présenter le lieu (\esp{en el colegio, en la calle}), nommer la conduite avec le présent, ajouter une action en cours avec \esp{estar} + gérondif, conclure par une conséquence ou une solution (\esp{la prevención, educar, ayudar}).
\item \textbf{Traduction} : « être en train de » se rend toujours par \esp{estar} + gérondif, jamais par \esp{ser}.
\end{itemize}

\textbf{4. Pièges de francophones}

\begin{itemize}
\item \esp{estar} (état, action en cours) $\neq$ \esp{ser} (caractéristique permanente) : \esp{El joven está enfermo} / \esp{El joven es valiente}.
\item Pas de gérondif après \esp{ser} : jamais \esp{es fumando}.
\item Accent obligatoire : \esp{adicción, prevención, educación}.
\item \esp{las drogas} au pluriel en général ; \esp{drogadicto} = toxicomane.
\end{itemize}

\end{bilan}

\begin{aretenir}[Checklist avant le Bac]
\begin{itemize}
\item[$\square$] Je connais le lexique du risque : \esp{droga, alcohol, tabaco, violencia, acoso escolar, sexting, absentismo, delincuencia, adicción}.
\item[$\square$] Je sais conjuguer les verbes réguliers en \esp{-ar, -er, -ir} au présent sans erreur.
\item[$\square$] Je maîtrise les verbes à radical changeant (\esp{pensar $\rightarrow$ pienso, dormir $\rightarrow$ duermo}).
\item[$\square$] Je connais les premières personnes irrégulières : \esp{hago, salgo, pongo, conozco, sé}.
\item[$\square$] Je forme le gérondif correctement : \esp{hablando, comiendo, escribiendo, leyendo, durmiendo}.
\item[$\square$] J'emploie \esp{estar} + gérondif pour décrire une action en cours.
\item[$\square$] Je fais la différence entre \esp{ser} et \esp{estar}.
\item[$\square$] Je sais repérer dans un texte le milieu scolaire (\esp{aula, patio, profesores}) et le milieu social (\esp{calle, redes sociales, amigos}).
\item[$\square$] Je sais rédiger un paragraphe descriptif de 5 à 8 lignes sur une conduite à risque.
\item[$\square$] Je n'oublie ni les accents (\esp{adicción, prevención}) ni les signes ¿ ¡ dans mes productions.
\item[$\square$] Je peux proposer des solutions : \esp{prevenir, educar, informar, ayudar a los jóvenes}.
\item[$\square$] Je relis ma copie pour vérifier les accords et la ponctuation espagnole.
\end{itemize}
\end{aretenir}

\findecours

\end{document}
